% Les tensions alternatives sinusoïdales — PCT 3ème — Cours signé OnBuch+
% Programme officiel MINESEC. Compiler avec : tectonic N13-tensions-alternatives-sinusoidales.tex
\newcommand{\DOCMATIERE}{PCT}
\newcommand{\DOCNIVEAU}{3ème}
\newcommand{\DOCMODULE}{Module 2 : Électricité}
\newcommand{\DOCLECON}{N13}
\newcommand{\DOCTITRE}{Les tensions alternatives sinusoïdales}
% OnBuch+ — préambule « cours signés » ENRICHI (compilé avec tectonic / XeLaTeX)
% Système visuel « Pop » d'OnBuch+ : fond crème (jamais blanc), bordures encre
% épaisses, ombres « dures » décalées, titres Archivo Black en capitales.
% Version MATHÉMATIQUES : graphes (pgfplots/tikz), boîtes pédagogiques
% (définition, propriété, méthode, attention, le savais-tu, exercice résolu,
% exercice, TP, bilan), page de garde et signature OnBuch+ ancrée en bas.
%
% Macros attendues dans main.tex AVANT \input{preamble} :
%   \DOCMATIERE  \DOCNIVEAU  \DOCMODULE  \DOCLECON (numéro)  \DOCTITRE
\documentclass[11pt]{article}

\usepackage{fontspec}
\usepackage[french]{babel}
\usepackage[a4paper,margin=2cm,top=3.4cm,bottom=2.8cm,headheight=1.4cm,footskip=1.3cm]{geometry}
\usepackage{amsmath,amssymb,amsfonts}
\usepackage{mathtools} % \xmapsto, \coloneqq... souvent utilisés par les modèles
\usepackage{cancel} % simplifications barrées (bilans de forces)
% \abs{z} (module, valeur absolue) et \norm{u} (norme), utilisés par les modèles
\providecommand{\abs}{}\let\abs\relax\DeclarePairedDelimiter{\abs}{\lvert}{\rvert}
\providecommand{\norm}{}\let\norm\relax\DeclarePairedDelimiter{\norm}{\lVert}{\rVert}
\usepackage[table]{xcolor} % [table] : fournit aussi \rowcolors (alternance de lignes)
\usepackage{pagecolor}
\usepackage{tikz}
\usetikzlibrary{arrows.meta,positioning,calc,shapes.geometric,shapes.misc,shapes.arrows,decorations.pathmorphing,decorations.pathreplacing,decorations.markings,patterns,shadows,mindmap,trees,fit,backgrounds,matrix,intersections,angles,quotes}
\usepackage{pgfplots}
\usepackage[version=4]{mhchem} % équations nucléaires \ce{^{235}_{92}U}
\usepackage[european,siunitx]{circuitikz} % schémas électriques
\pgfplotsset{compat=1.18}
\usepgfplotslibrary{fillbetween,groupplots} % aires entre courbes (calcul intégral)
\usepackage{enumitem}
\usepackage{fancyhdr}
% [most] charge toutes les bibliothèques tcolorbox (listings, minted, raster,
% documentation...) et gonfle inutilement la mémoire TeX sur un document long
% (~300 boîtes) : on ne charge que ce qui est réellement utilisé.
\usepackage[skins,breakable]{tcolorbox}
\usepackage{graphicx}
\usepackage{colortbl}
\usepackage{booktabs}
\usepackage{tabularx}
\usepackage{multirow}
\usepackage{diagbox} % case d'en-tête diagonale des tableaux à double entrée
% Colonnes extensibles centrée / gauche / droite (C, L, R), souvent utilisées par les modèles
\newcolumntype{C}{>{\centering\arraybackslash}X}
\newcolumntype{L}{>{\raggedright\arraybackslash}X}
\newcolumntype{R}{>{\raggedleft\arraybackslash}X}
\usepackage{array}
\usepackage{multicol}
\usepackage{siunitx}
% parse-numbers=false : accepte aussi \SI{2,5 \times 10^{-3}}{...}
\sisetup{locale=FR,output-decimal-marker={,},group-minimum-digits=4,parse-numbers=false}
\DeclareSIUnit\ohm{\ensuremath{\symup{\Omega}}} % U+2126 absent de la police
\DeclareSIUnit\division{div} % réglages d'oscilloscope (ms/div, V/div)
\DeclareSIUnit\tr{tr} % tours (vitesse de rotation, ex. \SI{1500}{\tr\per\minute})
\DeclareSIUnit\molar{M} % concentration molaire (ex. \SI{3}{\molar}, \SI{1}{\micro\molar})
\DeclareSIUnit\an{an} % année (le modèle confond parfois avec \year, un compteur TeX) — ex. \SI{5}{\kilo\meter\per\an}
\DeclareSIUnit\FCFA{FCFA} % franc CFA (ex. \SI{500}{\FCFA\per\kilo\gram})
\DeclareSIUnit\degreeBrix{\textdegree Brix} % teneur en sucre d'un sirop/jus (réfractométrie)
\DeclareSIUnit\month{mois} % le modèle utilise parfois \month, un compteur TeX, comme unité de durée
\DeclareSIUnit\microliter{\micro\liter} % le modèle écrit parfois \microliter en un seul token
\DeclareSIUnit\million{millions} % dénombrement (ex. \SI{15}{\million\per\mL} spermatozoïdes)
\usepackage{qrcode}
% \degree : commande fréquemment utilisée par les modèles pour un angle ou une
% température en dehors de \SI{}{} (non fournie par les paquets chargés).
\providecommand{\degree}{^{\circ}}
% \qed (fin de démonstration), utilisé par les modèles sans amsthm
\providecommand{\qed}{\hfill$\square$}
% \barre : double barre des valeurs interdites dans un tableau de variations (tkz-tab)
\providecommand{\barre}{\ensuremath{\|}}
% environnement proof (amsthm non chargé) utilisé par les modèles
\ifdefined\proof\else\newenvironment{proof}[1][Démonstration]{\par\noindent\textit{#1.}\ }{\qed\par}\fi
\usepackage{lastpage}
\usepackage{hyperref}
\hypersetup{hidelinks}

% ============================================================
% Palette « Pop » OnBuch+ — mêmes hex que la classe P (plus_ui.dart)
% ============================================================
\definecolor{poporange}{HTML}{F59321}
\definecolor{popdark}{HTML}{E07A0C}
\definecolor{popcream}{HTML}{FFF6E8}
\definecolor{popink}{HTML}{1C1714}
\definecolor{poppurple}{HTML}{7A5AE0}
\definecolor{popgreen}{HTML}{2E9E6A}
\definecolor{poppink}{HTML}{E0457B}
\definecolor{popblue}{HTML}{2F7DE1}
\definecolor{popgold}{HTML}{C98A12}
\definecolor{popmuted}{HTML}{8A7F76}
\definecolor{popline}{HTML}{EADFD2}
\definecolor{popgray}{HTML}{8A7F76}
\colorlet{popgrey}{popgray}
% Alias tolérants (le modèle nomme parfois les couleurs différemment)
\colorlet{popwhite}{white}
\colorlet{popblack}{popink}
\colorlet{popred}{poppink}
\colorlet{popyellow}{popgold}
% Teintes claires (fonds de tableaux/encadrés) — le modèle nomme parfois
% les couleurs avec un suffixe L (ex. popblueL) sans les avoir définies.
\colorlet{poporangeL}{poporange!20}
\colorlet{popdarkL}{popdark!20}
\colorlet{poppurpleL}{poppurple!20}
\colorlet{popgreenL}{popgreen!20}
\colorlet{poppinkL}{poppink!20}
\colorlet{popblueL}{popblue!20}
\colorlet{popgoldL}{popgold!20}
\colorlet{popredL}{popred!20}
\colorlet{popgrayL}{popgray!20}
% Teintes claires pour fonds de tableaux / zones
\colorlet{popblueL}{popblue!10!white}
\colorlet{poporangeL}{poporange!14!white}
\colorlet{popgreenL}{popgreen!12!white}
\colorlet{poppurpleL}{poppurple!12!white}
\colorlet{poppinkL}{poppink!12!white}
\colorlet{popgoldL}{popgold!14!white}

\pagecolor{popcream}

% ============================================================
% Typographie
% ============================================================
\defaultfontfeatures{Ligatures=TeX}
\setmainfont{PlusJakartaSans-Medium.ttf}[Path=../../../fonts/,BoldFont=PlusJakartaSans-Bold.ttf]
% Maths en sans-serif (Fira Math) avec chiffres et lettres droites de Plus
% Jakarta, pour que les formules se fondent dans le texte.
\usepackage[math-style=ISO,bold-style=ISO]{unicode-math}
\setmathfont{FiraMath-Regular.otf}[Path=../../../fonts/]
\setmathfont{PlusJakartaSans-Medium.ttf}[Path=../../../fonts/,range={up/num,up/{Latin,latin},\mathup}]
\usepackage{icomma} % virgule décimale sans espace en mode math
% Caractères mathématiques Unicode absents des polices de texte (Plus Jakarta,
% Archivo Black) : rendus par leur équivalent mathématique, en texte comme en
% formule — sinon ils s'affichent en carré vide (ex. « Arithmétique dans ℤ »).
\usepackage{newunicodechar}
\newunicodechar{ℝ}{\ensuremath{\mathbb{R}}}
\newunicodechar{ℤ}{\ensuremath{\mathbb{Z}}}
\newunicodechar{ℂ}{\ensuremath{\mathbb{C}}}
\newunicodechar{ℕ}{\ensuremath{\mathbb{N}}}
\newunicodechar{ℚ}{\ensuremath{\mathbb{Q}}}
\newunicodechar{𝔻}{\ensuremath{\mathbb{D}}}
\newunicodechar{α}{\ensuremath{\alpha}}
\newunicodechar{β}{\ensuremath{\beta}}
\newunicodechar{γ}{\ensuremath{\gamma}}
\newunicodechar{δ}{\ensuremath{\delta}}
\newunicodechar{ε}{\ensuremath{\varepsilon}}
\newunicodechar{θ}{\ensuremath{\theta}}
\newunicodechar{λ}{\ensuremath{\lambda}}
\newunicodechar{μ}{\ensuremath{\mu}}
\newunicodechar{σ}{\ensuremath{\sigma}}
\newunicodechar{τ}{\ensuremath{\tau}}
\newunicodechar{φ}{\ensuremath{\varphi}}
\newunicodechar{ψ}{\ensuremath{\psi}}
\newunicodechar{ω}{\ensuremath{\omega}}
\newunicodechar{Σ}{\ensuremath{\Sigma}}
% majuscules produites par \MakeUppercase dans les titres (α-aminés -> Α)
\newunicodechar{Α}{\ensuremath{\alpha}}
\newunicodechar{Β}{\ensuremath{\beta}}
\newunicodechar{Γ}{\ensuremath{\Gamma}}
\newunicodechar{Δ}{\ensuremath{\Delta}}
\newunicodechar{Ε}{\ensuremath{\varepsilon}}
\newunicodechar{Θ}{\ensuremath{\Theta}}
\newunicodechar{Λ}{\ensuremath{\Lambda}}
\newunicodechar{Μ}{\ensuremath{\mu}}
\newunicodechar{Τ}{\ensuremath{\tau}}
\newunicodechar{Φ}{\ensuremath{\Phi}}
\newunicodechar{Ψ}{\ensuremath{\Psi}}
\newunicodechar{Ω}{\ensuremath{\Omega}}
\newunicodechar{↦}{\ensuremath{\mapsto}}
\newunicodechar{∈}{\ensuremath{\in}}
\newunicodechar{∉}{\ensuremath{\notin}}
\newunicodechar{∧}{\ensuremath{\wedge}}
\newunicodechar{⇔}{\ensuremath{\Leftrightarrow}}
\newunicodechar{⟺}{\ensuremath{\Longleftrightarrow}}
\newunicodechar{⇒}{\ensuremath{\Rightarrow}}
\newunicodechar{⟹}{\ensuremath{\Longrightarrow}}
\newunicodechar{∘}{\ensuremath{\circ}}
\newunicodechar{∪}{\ensuremath{\cup}}
\newunicodechar{∩}{\ensuremath{\cap}}
\newunicodechar{≡}{\ensuremath{\equiv}}
\newunicodechar{⊂}{\ensuremath{\subset}}
\newunicodechar{⊥}{\ensuremath{\perp}}
\newunicodechar{∃}{\ensuremath{\exists}}
\newunicodechar{∀}{\ensuremath{\forall}}
\newunicodechar{∛}{\ensuremath{\sqrt[3]{\,}}}
\newunicodechar{∜}{\ensuremath{\sqrt[4]{\,}}}
\newunicodechar{⃗}{\ensuremath{{}^{\to}}}
\newunicodechar{ⁿ}{\ifmmode^{n}\else\textsuperscript{\ensuremath{n}}\fi}
\newunicodechar{ˣ}{\ifmmode^{x}\else\textsuperscript{\ensuremath{x}}\fi}
\newunicodechar{ᵇ}{\ifmmode^{b}\else\textsuperscript{\ensuremath{b}}\fi}
\newunicodechar{ʳ}{\ifmmode^{r}\else\textsuperscript{\ensuremath{r}}\fi}
\newunicodechar{ᵘ}{\ifmmode^{u}\else\textsuperscript{\ensuremath{u}}\fi}
\newunicodechar{ʸ}{\ifmmode^{y}\else\textsuperscript{\ensuremath{y}}\fi}
\newunicodechar{ᵃ}{\ifmmode^{a}\else\textsuperscript{\ensuremath{a}}\fi}
\newunicodechar{ᵉ}{\ifmmode^{e}\else\textsuperscript{\ensuremath{e}}\fi}
\newunicodechar{ᵏ}{\ifmmode^{k}\else\textsuperscript{\ensuremath{k}}\fi}
\newunicodechar{ᵖ}{\ifmmode^{p}\else\textsuperscript{\ensuremath{p}}\fi}
\newunicodechar{ᴺ}{\ifmmode^{N}\else\textsuperscript{\ensuremath{N}}\fi}
\newunicodechar{ᶜ}{\ifmmode^{c}\else\textsuperscript{\ensuremath{c}}\fi}
\newunicodechar{ʰ}{\ifmmode^{h}\else\textsuperscript{\ensuremath{h}}\fi}
\newunicodechar{ᵠ}{\ifmmode^{q}\else\textsuperscript{\ensuremath{q}}\fi}
\newunicodechar{ₙ}{\ifmmode_{n}\else\textsubscript{\ensuremath{n}}\fi}
\newunicodechar{ₐ}{\ifmmode_{a}\else\textsubscript{\ensuremath{a}}\fi}
\newunicodechar{ᵢ}{\ifmmode_{i}\else\textsubscript{\ensuremath{i}}\fi}
\newunicodechar{ᵣ}{\ifmmode_{r}\else\textsubscript{\ensuremath{r}}\fi}
\newunicodechar{ₖ}{\ifmmode_{k}\else\textsubscript{\ensuremath{k}}\fi}
\newunicodechar{ᵦ}{\ifmmode_{\beta}\else\textsubscript{\ensuremath{\beta}}\fi}
\newunicodechar{ₓ}{\ifmmode_{x}\else\textsubscript{\ensuremath{x}}\fi}
\newfontfamily\popdisplay{ArchivoBlack-Regular.ttf}[Path=../../../fonts/]
\newfontfamily\poplogo{SpaceGrotesk-Bold.ttf}[Path=../../../fonts/]
\newfontfamily\popsemi{PlusJakartaSans-SemiBold.ttf}[Path=../../../fonts/]
\newfontfamily\popxbold{PlusJakartaSans-ExtraBold.ttf}[Path=../../../fonts/]
\newfontfamily\popxboldls{PlusJakartaSans-ExtraBold.ttf}[Path=../../../fonts/,LetterSpace=3.0]

\setlist[enumerate]{leftmargin=1.7em,itemsep=5pt,topsep=4pt}
\setlist[itemize]{leftmargin=1.7em,itemsep=4pt,topsep=4pt,label={\color{poporange}\textbullet}}
\setlength{\parindent}{0pt}
\setlength{\parskip}{5pt}
\renewcommand{\arraystretch}{1.25}

% pgfplots : style Pop par défaut
\pgfplotsset{
  every axis/.append style={axis line style={popink,line width=1pt},tick style={popink},
    grid style={popline,line width=0.5pt},label style={font=\small},tick label style={font=\footnotesize}},
  popcurve/.style={poporange,line width=1.8pt,no markers,smooth},
}
% Même style disponible dans un \draw TikZ simple (hors pgfplots)
\tikzset{popcurve/.style={poporange,line width=1.8pt}}
\tikzset{popgrid/.style={popline,very thin}}
\colorlet{popcurve}{poporange} % le nom du style est parfois utilisé comme couleur

% ============================================================
% Pill / squiggle
% ============================================================
\newcommand{\poppill}[3]{%
  \tikz[baseline=-0.55ex]\node[draw=popink,line width=1.2pt,rounded corners=2.6pt,%
    fill=#2,inner xsep=6.5pt,inner ysep=3pt]{%
    \popxboldls\fontsize{7}{7}\selectfont\color{#3}\MakeUppercase{#1}};%
}
\newcommand{\popsquiggle}[1]{%
  \tikz[baseline=-0.4ex]{
    \draw[#1,line width=2.6pt,line cap=round]
      plot [smooth] coordinates {(0,0.09) (0.34,0.01) (0.68,0.21) (1.02,0.01) (1.36,0.10)};
  }%
}
% Mot-clé surligné (marqueur orange sous le texte)
\newcommand{\cle}[1]{{\popxbold\color{popdark}#1}}

% ============================================================
% En-tête / pied
% ============================================================
\pagestyle{fancy}
\fancyhf{}
\renewcommand{\headrulewidth}{0pt}
\renewcommand{\footrulewidth}{0pt}
\lhead{\raisebox{-0.15ex}{\poplogo\fontsize{13}{13}\selectfont\color{popink}OnBuch\color{poporange}+}}
\rhead{\poppill{\DOCMATIERE\ \textbullet\ \DOCNIVEAU\ \textbullet\ Leçon \DOCLECON}{white}{popink}}
\lfoot{\popxbold\fontsize{7.6}{7.6}\selectfont\color{popmuted}\MakeUppercase{Cours signé OnBuch+}\ \textbullet\ {\popsemi\DOCTITRE}}
\rfoot{\tikz[baseline=-0.6ex]\node[rounded corners=8.5pt,fill=popink,minimum height=17pt,minimum width=17pt,inner xsep=5pt,inner sep=0]{\popxbold\fontsize{8}{8}\selectfont\color{popcream}\thepage\,/\,\pageref*{LastPage}};}

% ============================================================
% Titres
% ============================================================
\newcommand{\chaptitre}[2]{%
  \begin{tcolorbox}[enhanced,colback=poporange,colframe=popink,boxrule=2.6pt,arc=11pt,
      drop shadow=popink,left=15pt,right=15pt,top=12pt,bottom=11pt,before skip=3pt,after skip=18pt]
    \poppill{#2}{popink}{popcream}\par\vspace{9pt}
    {\raggedright\hyphenpenalty=10000\exhyphenpenalty=10000\popdisplay\fontsize{22}{24}\selectfont\color{popcream}\MakeUppercase{#1}\par}
  \end{tcolorbox}
}
% Section numérotée : pastille orange avec le numéro + titre Archivo
\newcounter{popsec}
\newcommand{\coursec}[1]{%
  \stepcounter{popsec}\setcounter{popsub}{0}%
  \par\vspace{16pt}\noindent
  \begin{minipage}{\linewidth}
  \tikz[baseline=-0.7ex]\node[draw=popink,line width=1.6pt,fill=poporange,rounded corners=4pt,minimum size=22pt,inner sep=0]{\popdisplay\fontsize{12}{12}\selectfont\color{popcream}\thepopsec};%
  \hspace{8pt}{\popdisplay\fontsize{15}{17}\selectfont\color{popink}\MakeUppercase{#1}}\par
  \vspace{4pt}\noindent{\color{poporange}\rule{36pt}{3.6pt}}
  \end{minipage}\par\nopagebreak\vspace{8pt}
}
\newcounter{popsub}
\newcommand{\courssub}[1]{%
  \stepcounter{popsub}%
  \par\vspace{9pt}\noindent{\popxbold\fontsize{11.5}{13}\selectfont\color{popink}{\color{poporange}\thepopsec.\thepopsub}\hspace{6pt}#1}\par\nopagebreak\vspace{4pt}
}

% ============================================================
% Boîtes pédagogiques — même recette : fond blanc, bordure épaisse colorée,
% ombre dure encre, badge plein. Titre optionnel [#1].
% ============================================================
\tcbset{popbase/.style={breakable,enhanced,colback=white,boxrule=2.3pt,arc=9pt,
  shadow={3pt}{-3pt}{0pt}{popink},left=14pt,right=14pt,top=11pt,bottom=11pt,before skip=11pt,after skip=15pt,
  fontupper=\popsemi\fontsize{10.6}{14.5}\selectfont\color{popink},
  fontlower=\popsemi\fontsize{10.6}{14.5}\selectfont\color{popink}}}
% Coche dessinée (le glyphe ✓ n'existe pas dans les polices du document)
\newcommand{\popcheck}[1]{\tikz[baseline=-0.5ex]\draw[#1,line width=1.6pt,line cap=round,line join=round](-0.13,0.0)--(-0.04,-0.09)--(0.13,0.1);}
\providecommand{\coche}{\popcheck{popgreen}}
\providecommand{\resultat}[1]{\par\smallskip\noindent\fcolorbox{popgreen}{popgreenL}{\parbox{\dimexpr\linewidth-2\fboxsep-2\fboxrule\relax}{\textbf{Résultat.} #1}}\par\smallskip}
\newcommand{\popboxhead}[3]{\poppill{#1}{#2}{white}\ifx\relax#3\relax\else\hspace{6pt}{\popxbold\fontsize{10.6}{12}\selectfont\color{popink}#3}\fi\par\vspace{7pt}}

\newtcolorbox{aretenir}[1][]{popbase,colframe=popgreen,before upper={\popboxhead{À retenir}{popgreen}{#1}}}
\newtcolorbox{exemplebox}[1][]{popbase,colframe=poporange,before upper={\popboxhead{Exemple}{poporange}{#1}}}
\newtcolorbox{definition}[1][]{popbase,colframe=popblue,before upper={\popboxhead{Définition}{popblue}{#1}}}
\newtcolorbox{propriete}[1][]{popbase,colframe=poppurple,before upper={\popboxhead{Propriété}{poppurple}{#1}}}
\newtcolorbox{methode}[1][]{popbase,colframe=popgold,before upper={\popboxhead{Méthode}{popgold}{#1}}}
\newtcolorbox{attention}[1][]{popbase,colframe=poppink,before upper={\popboxhead{Attention}{poppink}{#1}}}
\newtcolorbox{experience}[1][]{popbase,colframe=popblue,colback=popblueL,before upper={\popboxhead{Expérience}{popblue}{#1}}}
% « Le savais-tu ? » : carte violette pleine (contexte, culture, Cameroun)
\newtcolorbox{savaistu}[1][]{popbase,colback=poppurple,colframe=popink,
  fontupper=\popsemi\fontsize{10.4}{14.2}\selectfont\color{white},
  before upper={\poppill{Le savais-tu\,?}{white}{poppurple}\ifx\relax#1\relax\else\hspace{6pt}{\popxbold\color{white}#1}\fi\par\vspace{7pt}}}
% Exercice résolu : énoncé en haut, \tcblower puis la solution sur fond crème
\newcounter{popexo}\newcounter{popexr}
\newtcolorbox{exoresolu}[1][]{popbase,colframe=poporange,colbacklower=popcream,
  segmentation style={popink,line width=1.2pt,dashed},
  before upper={\stepcounter{popexr}\popboxhead{Exercice résolu \thepopexr}{poporange}{#1}},
  before lower={\poppill{Solution}{popink}{popcream}\par\vspace{6pt}}}
% Exercice (non résolu en place) — la correction est dans la partie Corrigés
\newtcolorbox{exercice}[1][]{popbase,colframe=popink,
  before upper={\stepcounter{popexo}\popboxhead{Exercice \thepopexo}{popink}{#1}}}
% Corrigé
\newtcolorbox{corrige}[1][]{popbase,colframe=popgreen,colback=popgreenL,
  before upper={\popboxhead{Corrigé}{popgreen}{#1}}}
% Objectifs / prérequis / bilan
\newtcolorbox{objectifs}{popbase,colframe=popink,colback=poporangeL,
  before upper={\popboxhead{Objectifs de la leçon}{popink}{}}}
\newtcolorbox{prerequis}{popbase,colframe=popink,colback=popblueL,
  before upper={\popboxhead{Prérequis}{popblue}{}}}
\newtcolorbox{bilan}{popbase,colframe=popink,colback=white,boxrule=2.8pt,
  before upper={\popboxhead{Fiche bilan}{poporange}{L'essentiel à connaître pour l'examen}}}
% Figure encadrée avec légende
\newtcolorbox{popfigure}[1][]{enhanced,breakable=false,colback=white,colframe=popline,boxrule=1.4pt,arc=8pt,
  left=10pt,right=10pt,top=10pt,bottom=8pt,before skip=10pt,after skip=12pt,halign=center,
  lower separated=false,halign lower=center,
  fontlower=\popsemi\fontsize{9}{11.5}\selectfont\color{popmuted},
  #1}
\newcommand{\legende}[1]{\par\vspace{6pt}{\popsemi\fontsize{9}{11.5}\selectfont\color{popmuted}#1}}

% ============================================================
% Signature OnBuch+
% ============================================================
\newcommand{\poplogoinline}[1]{{\poplogo\fontsize{#1}{#1}\selectfont\color{popink}OnBuch\color{poporange}+}}
\newcommand{\signatureonbuch}{%
  \begin{tcolorbox}[enhanced,colback=popink,colframe=popink,boxrule=2.2pt,arc=10pt,
      drop shadow=poporange,left=14pt,right=14pt,top=10pt,bottom=10pt,before skip=0pt,after skip=0pt]
    \tikz[baseline=-0.7ex]\node[circle,fill=popgreen,draw=popcream,line width=1.3pt,minimum size=22pt,inner sep=0]{\popcheck{white}};%
    \hspace{9pt}\begin{minipage}[c]{0.62\linewidth}
      {\popxbold\fontsize{12}{13}\selectfont\color{popcream}Signé\ \textbullet\ {\poplogo OnBuch\color{poporange}+}}\par\vspace{2pt}
      {\raggedright\popsemi\fontsize{8}{10.5}\selectfont\color{popline}\MakeUppercase{Équipe pédagogique OnBuch+}\\\MakeUppercase{Conforme au programme officiel MINESEC}\par}
    \end{minipage}\hfill
    {\poplogo\fontsize{22}{22}\selectfont\color{popcream}OnBuch\color{poporange}+}
  \end{tcolorbox}%
}

% ============================================================
% Page de garde
% #1 titre · #2 matière · #3 niveau · #4 module · #5 n° leçon · #6 durée ·
% #7 description / objectifs (texte ou itemize)
% ============================================================
\newcommand{\pagegarde}[7]{%
  \thispagestyle{empty}
  \newgeometry{a4paper,margin=1.7cm,top=1.6cm,bottom=1.4cm}%
  \begin{tikzpicture}[remember picture,overlay]
    \fill[poporange!18] ([xshift=-3.2cm,yshift=-3.6cm]current page.north east) circle (4.6cm);
    \fill[poppurple!14] ([xshift=2.4cm,yshift=7.6cm]current page.south west) circle (3.8cm);
  \end{tikzpicture}%
  {\poplogo\fontsize{30}{30}\selectfont\color{popink}OnBuch\color{poporange}+}\hfill
  \raisebox{0.6ex}{\poppill{Cours signé \textbullet\ Premium}{popink}{popcream}}\par
  \vspace{1pt}\popsquiggle{poppurple}\par
  \vspace{8pt}
  \begin{tcolorbox}[enhanced,colback=poporange,colframe=popink,boxrule=2.8pt,arc=13pt,
      drop shadow=popink,left=18pt,right=18pt,top=12pt,bottom=13pt,after skip=10pt]
    \poppill{#2 \textbullet\ #3}{popink}{popcream}\hspace{5pt}\poppill{#4}{white}{popink}\par\vspace{8pt}
    {\popdisplay\fontsize{14}{14}\selectfont\color{popink}LEÇON #5}\par\vspace{4pt}
    {\raggedright\hyphenpenalty=10000\exhyphenpenalty=10000\popdisplay\fontsize{25}{27}\selectfont\color{popcream}\MakeUppercase{#1}\par}
    \vspace{8pt}
    {\popxbold\fontsize{9.5}{9.5}\selectfont\color{popink}\MakeUppercase{Durée conseillée : #6}}
  \end{tcolorbox}
  \begin{tcolorbox}[enhanced,colback=white,colframe=popink,boxrule=2.2pt,arc=10pt,
      drop shadow=popink,left=16pt,right=16pt,top=10pt,bottom=10pt,after skip=8pt]
    \poppill{Dans cette leçon}{poppurple}{white}\par\vspace{5pt}
    \popsemi\fontsize{9.3}{12.6}\selectfont\color{popink}#7
  \end{tcolorbox}
  \noindent\begin{minipage}[t]{0.487\linewidth}
    \begin{tcolorbox}[enhanced,colback=popgreen,colframe=popink,boxrule=2.2pt,arc=10pt,
        drop shadow=popink,left=11pt,right=11pt,top=10pt,bottom=10pt,height=3.1cm,valign=center]
      \tcbox[colback=white,colframe=popink,boxrule=1.3pt,arc=5pt,boxsep=0pt,nobeforeafter,
        left=3pt,right=3pt,top=3pt,bottom=3pt,valign=center]{\qrcode[height=1.9cm]{https://play.google.com/store/apps/details?id=com.luvvix.onbuch&hl=fr}}%
      \hspace{8pt}\begin{minipage}[c]{0.5\linewidth}
        {\popdisplay\fontsize{11}{12.5}\selectfont\color{white}\MakeUppercase{Télécharge OnBuch}}\par\vspace{4pt}
        {\raggedright\popsemi\fontsize{8.4}{11}\selectfont\color{white}Gratuit \textbullet\ Google Play\\Tuteur IA, annales, concours, résultats\par}
      \end{minipage}
    \end{tcolorbox}
  \end{minipage}\hfill
  \begin{minipage}[t]{0.487\linewidth}
    \begin{tcolorbox}[enhanced,colback=poppurple,colframe=popink,boxrule=2.2pt,arc=10pt,
        drop shadow=popink,left=11pt,right=11pt,top=10pt,bottom=10pt,height=3.1cm,valign=center]
      \tikz[baseline=-0.7ex]\node[circle,fill=white,draw=popink,line width=1.3pt,minimum size=1.6cm,inner sep=0]{\popdisplay\fontsize{24}{24}\selectfont\color{poppurple}+};%
      \hspace{8pt}\begin{minipage}[c]{0.6\linewidth}
        {\popdisplay\fontsize{11}{12.5}\selectfont\color{white}\MakeUppercase{Passe à OnBuch+}}\par\vspace{4pt}
        {\raggedright\popsemi\fontsize{8.4}{11}\selectfont\color{white}Tous les cours signés, Tuteur IA illimité, fiches PDF — onglet OnBuch+ de l'app\par}
      \end{minipage}
    \end{tcolorbox}
  \end{minipage}
  \vspace{8pt}
  \popcontenu
  \vfill
  \signatureonbuch
  \restoregeometry
  \newpage
}

% Rangée « Ce cours contient » (page de garde)
\newcommand{\poptile}[3]{% #1 couleur #2 gros texte #3 légende
  \begin{tcolorbox}[enhanced,colback=#1,colframe=popink,boxrule=2pt,arc=9pt,shadow={2.5pt}{-2.5pt}{0pt}{popink},
      left=6pt,right=6pt,top=9pt,bottom=9pt,halign=center,valign=center,height=1.75cm,width=\linewidth,nobeforeafter]
    {\popdisplay\fontsize{12.5}{13}\selectfont\color{white}\MakeUppercase{#2}}\par\vspace{3pt}
    {\centering\popsemi\fontsize{7.8}{9.5}\selectfont\color{white}#3\par}
  \end{tcolorbox}}
\newcommand{\popcontenu}{%
  \par\noindent{\popxboldls\fontsize{7.5}{7.5}\selectfont\color{popmuted}CE COURS SIGNÉ CONTIENT}\par\vspace{7pt}
  \noindent\begin{minipage}[t]{0.235\linewidth}\poptile{popblue}{Cours}{complet, illustré, pas à pas}\end{minipage}\hfill
  \begin{minipage}[t]{0.235\linewidth}\poptile{poporange}{Méthodes}{et exercices résolus}\end{minipage}\hfill
  \begin{minipage}[t]{0.235\linewidth}\poptile{poppink}{Exercices}{type examen + corrigés}\end{minipage}\hfill
  \begin{minipage}[t]{0.235\linewidth}\poptile{popgreen}{Fiche bilan}{l'essentiel à retenir}\end{minipage}\par}

% Sommaire
\newtcolorbox{sommaire}{enhanced,colback=white,colframe=popink,boxrule=2.2pt,arc=10pt,
  drop shadow=popink,left=18pt,right=18pt,top=13pt,bottom=13pt,before skip=6pt,after skip=20pt,
  before upper={\poppill{Sommaire}{poppurple}{white}\par\vspace{9pt}\popsemi\fontsize{10.6}{14.5}\selectfont\color{popink}}}

% Signature de fin de cours — poussée tout en bas de la dernière page
\newcommand{\findecours}{%
  \par\vfill\nopagebreak
  \begin{center}\popsquiggle{poporange}\end{center}\vspace{4pt}
  \signatureonbuch
}
\AtBeginDocument{\def\llbracket{\mathopen{[\mkern-3.5mu[}}\def\rrbracket{\mathclose{]\mkern-3.5mu]}}} % intervalles d'entiers (glyphe ⟦ absent de la police)
% Glyphes absents de Fira Math : redéfinis à partir de glyphes disponibles
\AtBeginDocument{%
  \def\setminus{\mathbin{\backslash}}\let\smallsetminus\setminus
  \def\vdots{\mathord{\vbox{\baselineskip3.2pt\lineskiplimit0pt\kern4pt\hbox{.}\hbox{.}\hbox{.}}}}%
  \def\ddots{\mathinner{\mkern1mu\raise6.5pt\hbox{.}\mkern2mu\raise3.6pt\hbox{.}\mkern2mu\raise0.7pt\hbox{.}\mkern1mu}}%
  \def\triangle{\mathord{\scalebox{0.9}{$\symup{\Delta}$}}}%
  \def\nabla{\mathord{\raisebox{\depth}{\scalebox{1}[-1]{$\symup{\Delta}$}}}}%
  \def\checkmark{\popcheck{popdark}}%
  \def\star{\mathbin{\ast}}\def\bigstar{\mathord{\ast}}%
  \def\diamond{\mathbin{\scalebox{0.6}{$\Diamond$}}}%
  \def\bigcup{\mathop{\mathchoice{\raisebox{-0.3em}{\scalebox{1.7}{$\cup$}}}{\scalebox{1.25}{$\cup$}}{\cup}{\cup}}\displaylimits}%
  \def\bigcap{\mathop{\mathchoice{\raisebox{-0.3em}{\scalebox{1.7}{$\cap$}}}{\scalebox{1.25}{$\cap$}}{\cap}{\cap}}\displaylimits}%
  \def\lesseqqgtr{\mathrel{\lessgtr}}\def\bigoplus{\mathop{\oplus}\displaylimits}%
}
\providecommand{\ssi}{si et seulement si} % abréviation fréquente
\AtBeginDocument{\providecommand{\wideparen}{\overparen}} % arc de cercle

\begin{document}

\pagegarde{Les tensions alternatives sinusoïdales}{PCT}{3ème}{Module 2 : Électricité}{N13}{1 séance (1 h chacune)}{Cette leçon introduit les tensions alternatives sinusoïdales, leur visualisation à l'oscilloscope et leurs grandeurs caractéristiques (période, fréquence, tension maximale et efficace). Elle permet à l'élève de comprendre la nature de la tension du secteur et d'effectuer des mesures et des calculs simples en toute sécurité.
\vspace{4pt}
\begin{multicols}{2}\small
\begin{itemize}
\item À la fin de cette leçon, je sais distinguer une tension continue d'une tension alternative et reconnaître une tension sinusoïdale
\item régler et utiliser un oscilloscope pour visualiser une tension alternative
\item définir la période et la fréquence d'une tension périodique et citer leurs unités
\item exploiter l'oscillogramme d'une tension sinusoïdale pour déterminer sa période et sa fréquence
\item définir la tension maximale et la tension efficace et citer la relation qui les lie
\item calculer une tension efficace à partir de la tension maximale (et inversement) dans des situations concrètes
\end{itemize}\end{multicols}}

\begin{sommaire}
\begin{multicols}{2}
\begin{enumerate}
\item Tensions continues et tensions alternatives
\item Visualisation d'une tension : l'oscilloscope
\item Période et fréquence d'une tension sinusoïdale
\item Tension maximale et tension efficace
\item Exploitation complète d'un oscillogramme et sécurité
\item Activité d'intégration
\item Méthodes et exercices résolus
\item Exercices
\item Corrigés des exercices
\item Fiche bilan
\end{enumerate}
\end{multicols}
\end{sommaire}

\begin{objectifs}
\begin{itemize}
\item À la fin de cette leçon, je sais distinguer une tension continue d'une tension alternative et reconnaître une tension sinusoïdale
\item régler et utiliser un oscilloscope pour visualiser une tension alternative
\item définir la période et la fréquence d'une tension périodique et citer leurs unités
\item exploiter l'oscillogramme d'une tension sinusoïdale pour déterminer sa période et sa fréquence
\item définir la tension maximale et la tension efficace et citer la relation qui les lie
\item calculer une tension efficace à partir de la tension maximale (et inversement) dans des situations concrètes
\item interpréter les indications « 220 V – 50 Hz » portées sur les installations et appareils électriques
\item appliquer les consignes de sécurité lors de l'utilisation du secteur et de l'oscilloscope
\end{itemize}
\end{objectifs}

\begin{prerequis}
\begin{itemize}
\item Notion de tension électrique et son unité, le volt (leçon sur les tensions continues)
\item Utilisation du voltmètre et du multimètre (positions AC et DC)
\item Notion de courant alternatif et de générateur basse fréquence (GBF)
\item Lecture et exploitation d'un graphique en fonction du temps
\item Notion de grandeur périodique (exemples de la vie courante : pendule, rotation)
\end{itemize}
\end{prerequis}

\newpage

\coursec{Tensions continues et tensions alternatives}

\courssub{Situation d'accroche : le mystère des \SI{220}{V}, \SI{50}{Hz}}

Le soir tombe sur Yaoundé. Dans le quartier Mvog-Ada, maman branche le chargeur du téléphone portable sur la prise murale. Sur l'étiquette du chargeur, on lit : « Entrée : \SI{220}{V} — \SI{50}{Hz} ». À Douala, à Bafoussam, à Maroua, toutes les prises du réseau ENEO affichent les mêmes deux nombres. Mais que signifient-ils exactement ?

Voici une question plus troublante encore : la tension de la prise \emph{varie en permanence} — elle monte, descend, s'inverse \SI{50}{fois par seconde} — et pourtant le téléphone ne grille pas. Pourquoi ? Et lorsqu'un technicien de la SONREL branche un appareil appelé \cle{oscilloscope} sur le réseau pour « voir » la tension, qu'observe-t-il sur l'écran ?

Dans les leçons précédentes, nous avons travaillé avec des piles et des accumulateurs : la tension était \emph{continue}, stable, fidèle. Avec le réseau électrique, tout change. Cette leçon nous apprend à \textbf{décrire, mesurer et interpréter les tensions alternatives sinusoïdales} qui alimentent nos foyers. Commençons par distinguer les deux grandes familles de tensions : les tensions continues et les tensions alternatives.

\begin{attention}[Danger du secteur \SI{220}{V} — Consigne de sécurité]
La tension des prises murales (\SI{220}{V} alternatif) est \textbf{dangereuse} : elle peut provoquer des brûlures graves, un arrêt cardiaque ou la mort. \textbf{Ne jamais} brancher un oscilloscope scolaire directement sur une prise secteur. Au laboratoire, on utilise exclusivement le \textbf{GBF (Générateur Basse Fréquence)} qui délivre une tension de sécurité (généralement $< \SI{20}{V}$). L'oscilloscope doit rester sur sa prise de terre (prise à 3 broches).
\end{attention}

\courssub{Rappel : la tension continue, une tension fidèle}

Tu connais déjà bien les piles et les accumulateurs : ce sont eux qui alimentent la lampe torche, la radio, la télécommande ou la moto-taxi (pour son démarreur). Ces générateurs possèdent deux bornes bien distinctes : une borne \cle{positive} ($+$) et une borne \cle{négative} ($-$).

\begin{definition}[Tension continue]
Une tension est dite \cle{continue} lorsqu'elle garde une \textbf{valeur constante} au cours du temps et conserve \textbf{toujours le même signe}. Elle est délivrée par les piles, les accumulateurs (batteries) et les panneaux solaires.
\end{definition}

Si l'on mesure la tension d'une pile neuve de \SI{4,5}{V} à différents instants — maintenant, dans une minute, dans une heure — on trouve toujours environ \SI{4,5}{V}. Sur un graphique représentant la tension $u$ en fonction du temps $t$, on obtient une \textbf{droite horizontale} : rien ne bouge.

\begin{attention}[« Continue » ne veut pas dire « éternelle »]
Une pile qui s'use voit sa tension diminuer lentement (\SI{4,5}{V}, puis \SI{4,2}{V}, puis \SI{3,8}{V}...). On parle malgré tout de tension continue, car elle \textbf{ne change jamais de signe} : la borne $+$ reste la borne $+$. Le mot « continu » signifie « qui ne s'inverse pas », pas « qui dure toujours ».
\end{attention}

\courssub{Expérience d'introduction : comparer une pile et un GBF à l'oscilloscope}

Pour découvrir expérimentalement la différence entre les deux familles de tensions, réalisons au laboratoire la manipulation suivante, en démarche d'investigation.

\begin{experience}[Pile ou GBF : que voit l'oscilloscope ?]
\textbf{Matériel} : un oscilloscope, une pile plate de \SI{4,5}{V}, un générateur basse fréquence (GBF), des fils de connexion.

\textbf{Protocole} :
\begin{enumerate}
\item Brancher la pile sur la voie d'entrée de l'oscilloscope (borne $+$ sur l'entrée, borne $-$ sur la masse). Observer la trace.
\item Débrancher la pile, puis brancher le GBF réglé sur le signal « sinusoïdal ». Observer la nouvelle trace.
\item Régler successivement le GBF sur « triangulaire » puis sur « créneaux » (signal rectangulaire). Observer à chaque fois.
\end{enumerate}

\textbf{Observations} :
\begin{itemize}
\item Avec la pile : la trace est une \textbf{droite horizontale}, décalée vers le haut par rapport au zéro. Elle ne bouge pas au cours du temps.
\item Avec le GBF : la trace \textbf{ondulée} monte et descend régulièrement, en passant tantôt au-dessus, tantôt en dessous du zéro. Le motif se répète sans fin.
\end{itemize}

\textbf{Interprétation} : la pile délivre une tension qui garde la même valeur et le même signe : c'est une tension \emph{continue}. Le GBF délivre une tension qui varie et \emph{change de signe} au cours du temps : c'est une tension \emph{alternative}.

\textbf{Conclusion} : il existe deux grandes familles de tensions : les tensions continues (piles, batteries) et les tensions alternatives (GBF, prises du réseau).
\end{experience}

\courssub{Tension alternative, tension périodique : les définitions}

L'expérience précédente nous a montré une trace qui monte, descend, passe au-dessus puis en dessous du zéro. Que signifie concrètement ce « changement de signe » ? Pendant une partie du temps, c'est la borne A du générateur qui joue le rôle de borne $+$ ; un instant plus tard, c'est la borne B ! Le courant dans le circuit change alors de sens, comme une rivière qui refluerait régulièrement. C'est exactement ce qui se passe dans les fils de ta maison.

\begin{definition}[Tension alternative]
Une tension est dite \cle{alternative} lorsqu'elle \textbf{change de signe} au cours du temps : elle est tantôt positive, tantôt négative.
\end{definition}

\begin{definition}[Tension périodique et période]
Une tension est dite \cle{périodique} lorsqu'elle \textbf{se reproduit identique à elle-même} à intervalles de temps égaux. La plus petite durée au bout de laquelle le motif se répète s'appelle la \cle{période} (notée $T$, en secondes \SI{}{s}).
\end{definition}

La tension du réseau ENEO est à la fois alternative \emph{et} périodique : son motif se répète \SI{50}{fois par seconde}, sans jamais se fatiguer. On dit que sa \cle{fréquence} $f$ vaut \SI{50}{Hz} (Hertz).

\begin{attention}[Le piège du mot « alternative »]
Erreur très fréquente au BEPC : croire que « alternative » signifie simplement « variable ». C'est faux !
\begin{itemize}
\item Une tension qui varie entre \SI{3}{V} et \SI{5}{V} en restant toujours positive est \textbf{variable}, mais elle n'est \textbf{pas alternative}.
\item Pour être alternative, une tension doit \textbf{changer de signe} : passer par exemple de $+\SI{5}{V}$ à $-\SI{5}{V}$.
\end{itemize}
Retiens : \emph{alternative} = \emph{qui change de signe}. Toute tension alternative est variable, mais toute tension variable n'est pas alternative.
\end{attention}

\courssub{Les différentes formes de tensions alternatives}

L'oscilloscope révèle que toutes les tensions alternatives ne se ressemblent pas. Le GBF du laboratoire peut en produire trois formes classiques, que tout élève de 3ème doit savoir reconnaître :

\begin{itemize}
\item la tension \cle{sinusoïdale} : une ondulation douce et régulière, comme une vague ; c'est la forme de la tension du réseau ENEO ;
\item la tension \cle{triangulaire} : elle monte et descend en lignes droites, comme un toit de case répété ;
\item la tension \cle{en créneaux} (ou rectangulaire) : elle saute brutalement d'une valeur positive à une valeur négative, comme des marches d'escalier.
\end{itemize}

\begin{popfigure}
\begin{center}
\begin{tikzpicture}[scale=0.8]
% ---- Oscillogramme 1 : tension continue ----
\begin{scope}
\draw[popline,step=0.5,very thin] (0,-1.2) grid (4.5,1.2);
\draw[->,popdark] (0,0) -- (4.7,0) node[right,font=\small]{$t$};
\draw[->,popdark] (0,-1.3) -- (0,1.4) node[above,font=\small]{$u$};
\draw[popblue,very thick] (0,0.8) -- (4.5,0.8);
\node[font=\small\bfseries,text=popdark] at (2.25,-1.7) {a) Tension continue};
\end{scope}
% ---- Oscillogramme 2 : sinusoidale ----
\begin{scope}[xshift=5.5cm]
\draw[popline,step=0.5,very thin] (0,-1.2) grid (4.5,1.2);
\draw[->,popdark] (0,0) -- (4.7,0) node[right,font=\small]{$t$};
\draw[->,popdark] (0,-1.3) -- (0,1.4) node[above,font=\small]{$u$};
\draw[poporange,very thick,domain=0:4.5,samples=120] plot (\x,{0.9*sin(360*\x/2.25)});
\node[font=\small\bfseries,text=popdark] at (2.25,-1.7) {b) Tension sinusoïdale};
\end{scope}
% ---- Oscillogramme 3 : creneaux ----
\begin{scope}[xshift=11cm]
\draw[popline,step=0.5,very thin] (0,-1.2) grid (4.5,1.2);
\draw[->,popdark] (0,0) -- (4.7,0) node[right,font=\small]{$t$};
\draw[->,popdark] (0,-1.3) -- (0,1.4) node[above,font=\small]{$u$};
\draw[popgreen,very thick] (0,0.8) -- (1.125,0.8) -- (1.125,-0.8) -- (2.25,-0.8) -- (2.25,0.8) -- (3.375,0.8) -- (3.375,-0.8) -- (4.5,-0.8);
\node[font=\small\bfseries,text=popdark] at (2.25,-1.7) {c) Tension en créneaux};
\end{scope}
\end{tikzpicture}
\end{center}
\legende{Trois oscillogrammes comparés. a) La tension continue (pile) est une droite horizontale : valeur et signe constants. b) La tension sinusoïdale ondule en douceur de part et d'autre du zéro. c) La tension en créneaux saute brutalement d'une valeur positive à une valeur négative. Les tensions b) et c) sont alternatives et périodiques.}
\end{popfigure}

\begin{exemplebox}[Pile ou prise murale ?]
\begin{itemize}
\item La tension d'une pile plate de \SI{4,5}{V} est une tension \textbf{continue} : elle garde une valeur constante et ne change jamais de signe. Sa trace à l'oscilloscope est une droite horizontale.
\item La tension d'une prise murale ENEO est une tension \textbf{alternative sinusoïdale} : elle varie en douceur et change de signe \SI{50}{fois par seconde}. Sa trace est une belle vague régulière.
\end{itemize}
\end{exemplebox}

\courssub{D'où viennent les tensions alternatives ?}

D'où vient cette tension qui ondule sans cesse dans nos prises ? Elle est produite par des machines appelées \cle{alternateurs}. Le principe est simple et génial : lorsqu'on fait \textbf{tourner un aimant devant une bobine de fil de cuivre} (ou une bobine devant un aimant), il naît aux bornes de la bobine une tension alternative. Plus l'aimant tourne régulièrement, plus la tension produite est une belle sinusoïde.

\begin{popfigure}
\begin{center}
\begin{tikzpicture}[scale=1.0]
% Aimant tournant
\draw[popdark,thick,fill=poporangeL] (-1.6,-0.5) rectangle (0,0.5);
\draw[popdark,thick,fill=popblueL] (0,-0.5) rectangle (1.6,0.5);
\node[font=\bfseries,text=popdark] at (-0.8,0) {N};
\node[font=\bfseries,text=popdark] at (0.8,0) {S};
\draw[->,popgold,very thick] (2.2,1.1) arc (60:120:2.2);
\node[font=\small,text=popdark] at (0,1.9) {rotation};
\fill[popdark] (0,0) circle (0.06);
% Axe
\draw[popdark,dashed] (0,-1.6) -- (0,1.6);
% Bobine
\begin{scope}[xshift=4.2cm]
\draw[popdark,thick] (-0.9,-1.1) rectangle (0.9,1.1);
\foreach \y in {-0.8,-0.4,0,0.4,0.8}{
  \draw[poppurple,thick] (-0.9,\y) arc (180:360:0.9 and 0.25);
  \draw[poppurple,thick] (-0.9,\y) arc (180:0:0.9 and 0.25);
}
\node[font=\small,text=popdark] at (0,-1.6) {bobine de cuivre};
% Fils de sortie
\draw[poppurple,thick] (0.9,-0.4) -- (2.1,-0.4) node[right,font=\small,text=popdark]{$u$ alternative};
\draw[poppurple,thick] (0.9,0.4) -- (2.1,0.4);
\end{scope}
% Fleche champ
\draw[->,popmuted,thick] (1.9,0) -- (3.0,0);
\node[font=\small,text=popmuted] at (2.45,0.35) {champ};
\node[font=\small,text=popmuted] at (2.45,-0.4) {magnétique};
\end{tikzpicture}
\end{center}
\legende{Schéma simplifié d'un alternateur. Un aimant (pôles N et S) mis en rotation devant une bobine de fil de cuivre y fait naître une tension alternative : c'est le phénomène d'induction électromagnétique, étudié plus tard. C'est le principe de toutes nos centrales.}
\end{popfigure}

Au Cameroun, les sources de tensions alternatives sont partout autour de nous :

\begin{itemize}
\item les \textbf{alternateurs des centrales hydroélectriques} : à Song Loulou et à Édéa, sur le fleuve Sanaga, la chute d'eau fait tourner d'énormes turbines, qui entraînent des alternateurs géants ;
\item les \textbf{groupes électrogènes} : quand ENEO déleste, le groupe électrogène du quartier prend le relais ; son moteur (essence ou gasoil) fait tourner un alternateur qui délivre une tension alternative de \SI{220}{V} ;
\item le \textbf{GBF du laboratoire} : un petit générateur électronique qui produit des tensions alternatives de forme et de fréquence réglables, parfait pour nos expériences.
\end{itemize}

\begin{savaistu}[Le barrage de Song Loulou, géant de la Sanaga]
Le barrage hydroélectrique de \textbf{Song Loulou}, sur le fleuve Sanaga, est l'une des plus grandes centrales du Cameroun. Son principe est une magnifique chaîne de transformations d'énergie : l'eau retenue en hauteur possède de l'\emph{énergie potentielle} ; en tombant, elle la transforme en \emph{énergie de mouvement} qui fait tourner les turbines ; celles-ci entraînent les alternateurs, qui la transforment en \emph{énergie électrique} sous forme d'une tension alternative sinusoïdale de \SI{50}{Hz}. Avec la centrale d'Édéa, Song Loulou alimente une grande partie du sud du pays, jusqu'à ta prise murale !
\end{savaistu}

\begin{exoresolu}[Continue ou alternative ? À toi de classer]
Énoncé. On observe à l'oscilloscope les tensions délivrées par quatre appareils : (1) une batterie de moto ; (2) une prise murale ENEO ; (3) un panneau solaire ; (4) un groupe électrogène. Pour chacun, indiquer si la tension est continue ou alternative, en justifiant.
\tcblower
Solution détaillée.
\begin{enumerate}
\item \textbf{Batterie de moto} : tension \emph{continue}. C'est un accumulateur : la tension garde une valeur constante (environ \SI{12}{V}) et ne change jamais de signe. Trace : droite horizontale.
\item \textbf{Prise murale ENEO} : tension \emph{alternative sinusoïdale}. Elle provient des alternateurs des centrales (Song Loulou, Édéa) : elle change de signe \SI{50}{fois par seconde}. Trace : sinusoïde.
\item \textbf{Panneau solaire} : tension \emph{continue}. Comme une pile, il possède une borne $+$ et une borne $-$ fixes ; sa tension reste positive tant qu'il est éclairé.
\item \textbf{Groupe électrogène} : tension \emph{alternative}. Son moteur entraîne un alternateur (aimant tournant devant une bobine), qui produit une tension sinusoïdale de \SI{220}{V}, comme le réseau.
\end{enumerate}
\textbf{Astuce} : pile, batterie, panneau solaire $\rightarrow$ continu ; alternateur, prise du réseau, groupe électrogène, GBF $\rightarrow$ alternatif.
\end{exoresolu}

\begin{aretenir}[L'essentiel de cette partie]
\begin{itemize}
\item Une tension \cle{continue} (pile, accumulateur, panneau solaire) garde une valeur constante et le même signe ; sa trace à l'oscilloscope est une droite horizontale.
\item Une tension \cle{alternative} \textbf{change de signe} au cours du temps ; elle est \cle{périodique} si elle se reproduit identique à elle-même à intervalles égaux. La plus petite durée de répétition est la \cle{période} $T$.
\item Les tensions alternatives peuvent être \textbf{sinusoïdales}, \textbf{triangulaires} ou \textbf{en créneaux}.
\item La tension des prises du réseau camerounais (ENEO) est \textbf{alternative sinusoïdale} de fréquence \SI{50}{Hz} ; elle est produite par les \textbf{alternateurs} des centrales hydroélectriques (Song Loulou, Édéa) ou des groupes électrogènes.
\item Piège : « alternative » ne veut pas dire « variable », mais « qui change de signe ».
\item Sécurité : \textbf{ne jamais} connecter un oscilloscope scolaire directement sur le secteur \SI{220}{V} (danger de mort). On utilise le GBF (Basse Tension de Sécurité).
\end{itemize}
\end{aretenir}

Maintenant que nous savons distinguer les deux familles de tensions, il nous faut un instrument pour les « voir » et les mesurer avec précision : c'est l'oscilloscope, que nous découvrirons dans la partie suivante.

\coursec{Visualisation d'une tension : l'oscilloscope}

Nous venons de distinguer les tensions continues (comme celle d'une pile) des tensions alternatives (comme celle du réseau ENEO), qui changent de valeur et de sens au cours du temps. Mais une tension est invisible : on ne peut pas la voir à l'œil nu, contrairement à la flamme d'un feu ou au déplacement d'une moto-taxi. Comment alors \emph{observer} la tension du secteur qui monte, descend, puis s'inverse 50 fois par seconde ? Comment vérifier qu'une tension est bien continue ou alternative ? Il nous faut un instrument capable de \cle{dessiner} la tension : cet instrument s'appelle l'\cle{oscilloscope}.

\begin{savaistu}[Un appareil partout autour de toi]
L'oscilloscope est l'« œil » de l'électricien et de l'électronicien. Le technicien qui répare ton téléviseur à Yaoundé, l'ingénieur d'ENEO qui contrôle le réseau, le mécanicien qui diagnostique l'allumage d'une moto : tous utilisent un oscilloscope pour \emph{voir} les tensions. À l'hôpital, l'écran qui affiche les battements du cœur (électrocardiogramme) fonctionne sur le même principe.
\end{savaistu}

\courssub{Rôle de l'oscilloscope}

\begin{definition}[L'oscilloscope]
L'\cle{oscilloscope} est un appareil de mesure qui permet de \cle{visualiser} sur un écran l'\cle{évolution d'une tension électrique au cours du temps}. La courbe obtenue sur l'écran s'appelle un \cle{oscillogramme}.
\end{definition}

L'idée est simple : l'oscilloscope transforme une tension (invisible) en un dessin (visible). Sur l'écran :
\begin{itemize}
\item l'\cle{axe vertical} représente la valeur de la tension $u$ (en volts) ;
\item l'\cle{axe horizontal} représente le temps $t$ (en secondes ou millisecondes).
\end{itemize}
Ainsi, l'oscillogramme est la représentation graphique de la fonction $u(t)$ : c'est exactement comme la courbe de température qu'un médecin trace pour suivre un malade, sauf qu'ici c'est la tension qui est « suivie » instant après instant.

\begin{exemplebox}[Intuition : le stylo et le ruban]
Imagine un stylo qui se déplace verticalement au rythme de la tension : il monte quand la tension augmente, il descend quand elle diminue. Si pendant ce temps on fait défiler un ruban de papier horizontalement à vitesse constante sous le stylo, celui-ci trace une courbe : c'est l'oscillogramme. L'oscilloscope fait exactement cela, mais avec un faisceau d'électrons (ou un écran numérique) à la place du stylo.
\end{exemplebox}

\courssub{Description simplifiée de l'oscilloscope}

Observons la face avant d'un oscilloscope scolaire. On y distingue trois parties essentielles.

\begin{aretenir}[Les trois parties de la face avant]
\begin{enumerate}
\item \textbf{L'écran gradué} : il est divisé en carreaux appelés \cle{divisions} (div). Un axe horizontal et un axe vertical, gradués en divisions, servent de repères pour les mesures.
\item \textbf{La voie d'entrée (voie Y)} : ce sont deux bornes de branchement : la borne d'entrée \cle{Y} (souvent rouge) et la borne \cle{masse} (souvent noire, repérée par le symbole de terre). La tension à étudier se branche entre ces deux bornes.
\item \textbf{Les boutons de réglage} : parmi eux, deux sont indispensables au BEPC :
\begin{itemize}
\item le bouton de \cle{sensibilité verticale} (en V/div) ;
\item le bouton de \cle{durée de balayage} (en ms/div).
\end{itemize}
\end{enumerate}
\end{aretenir}

\begin{popfigure}
\begin{center}
\begin{tikzpicture}[scale=1]
% Boîtier
\draw[fill=popcream,draw=popdark,line width=1pt,rounded corners=4pt] (0,0) rectangle (10.5,6.5);
% Écran
\draw[fill=popblue!8,draw=popdark,line width=1pt] (0.5,1.2) rectangle (6.3,6);
% Grille de l'écran
\foreach \x in {1.08,1.66,2.24,2.82,3.4,3.98,4.56,5.14,5.72}
  \draw[popline,thin] (\x,1.2) -- (\x,6);
\foreach \y in {1.8,2.4,3.0,3.6,4.2,4.8,5.4}
  \draw[popline,thin] (0.5,\y) -- (6.3,\y);
% Axes centraux plus marqués
\draw[popmuted,line width=0.8pt] (3.4,1.2) -- (3.4,6);
\draw[popmuted,line width=0.8pt] (0.5,3.6) -- (6.3,3.6);
% Courbe sinusoïdale sur l'écran
\draw[poporange,line width=1.4pt] plot[domain=0.5:6.3,samples=120] (\x,{3.6+1.2*sin(deg(2*pi*(\x-0.5)/2.9))});
% Étiquette écran
\node[popdark,font=\small] at (3.4,0.85) {Écran gradué en divisions};
% Bouton sensibilité verticale
\draw[fill=popblueL,draw=popdark] (7.6,4.6) circle (0.5);
\draw[popdark,line width=1pt] (7.6,4.6) -- (7.25,5.0);
\node[popdark,font=\small,align=center] at (8.9,4.6) {Sensibilité verticale\\ (V/div)};
% Bouton balayage
\draw[fill=popgreenL,draw=popdark] (7.6,3.2) circle (0.5);
\draw[popdark,line width=1pt] (7.6,3.2) -- (8.0,3.55);
\node[popdark,font=\small,align=center] at (8.9,3.2) {Durée de balayage\\ (ms/div)};
% Bornes d'entrée
\draw[fill=red!70,draw=popdark] (7.6,1.7) circle (0.18);
\node[popdark,font=\small] at (8.15,1.7) {Y};
\draw[fill=black!70,draw=popdark] (7.6,0.9) circle (0.18);
\node[popdark,font=\small] at (8.6,0.9) {Masse};
\end{tikzpicture}
\end{center}
\legende{Face avant simplifiée d'un oscilloscope scolaire : écran gradué en divisions, bouton de sensibilité verticale (V/div), bouton de durée de balayage (ms/div), bornes d'entrée Y (rouge) et masse (noire).}
\end{popfigure}

\courssub{Les deux réglages fondamentaux}

\textbf{La sensibilité verticale.}\par
Le bouton de sensibilité verticale indique combien de volts représente \emph{une division verticale} de l'écran. C'est l'« échelle » de l'axe des tensions.

\begin{definition}[Sensibilité verticale]
La \cle{sensibilité verticale}, notée $S_v$, est le nombre de volts représentés par une division verticale de l'écran. Elle s'exprime en volts par division (V/div). La tension correspondant à une déviation de $y$ divisions vaut :
\[ u = S_v \times y \]
\end{definition}

\begin{exemplebox}[Exemple chiffré]
La sensibilité verticale est réglée sur $S_v = \SI{2}{V/div}$. Le spot se déplace de $y = 3$ divisions vers le haut. La tension mesurée est :
\[ u = S_v \times y = 2 \times 3 = \SI{6}{V} \]
Si le spot était descendu de 3 divisions sous l'axe, on lirait $u = \SI{-6}{V}$ : le signe indique le sens de la tension.
\end{exemplebox}

\textbf{La durée de balayage.}\par
Le bouton de durée de balayage indique combien de temps représente \emph{une division horizontale} de l'écran. C'est l'« échelle » de l'axe des temps.

\begin{definition}[Durée de balayage]
La \cle{durée de balayage}, notée $B$, est la durée représentée par une division horizontale de l'écran. Elle s'exprime en millisecondes par division (ms/div). Une étendue horizontale de $x$ divisions correspond à une durée :
\[ t = B \times x \]
\end{definition}

\begin{exemplebox}[Exemple chiffré]
La durée de balayage est réglée sur $B = \SI{5}{ms/div}$. Un motif de la courbe s'étend sur $x = 4$ divisions horizontales. La durée correspondante est :
\[ t = B \times x = 5 \times 4 = \SI{20}{ms} \]
\end{exemplebox}

\begin{attention}[Ne confonds pas les deux réglages !]
Erreur très fréquente au BEPC : confondre la \textbf{sensibilité verticale} (en V/div) avec la \textbf{durée de balayage} (en ms/div). Retiens bien :
\begin{itemize}
\item \textbf{vertical} $\rightarrow$ \textbf{tension} $\rightarrow$ volts par division (V/div) ;
\item \textbf{horizontal} $\rightarrow$ \textbf{temps} $\rightarrow$ millisecondes par division (ms/div).
\end{itemize}
Une hauteur verticale ne donne JAMAIS une durée, et une étendue horizontale ne donne JAMAIS une tension. Vérifie toujours l'unité du bouton avant de calculer.
\end{attention}

\courssub{Branchement de l'oscilloscope}

\begin{propriete}[Règle de branchement]
Pour visualiser une tension, on branche l'oscilloscope \cle{en dérivation} (comme un voltmètre) aux bornes du dipôle ou du générateur à étudier : la tension se branche entre la borne d'entrée \cle{Y} et la borne \cle{masse}. L'oscilloscope ne se branche jamais en série dans le circuit.
\end{propriete}

\begin{experience}[Observation : que voit-on sans tension ?]
\textbf{Matériel} : un oscilloscope, un générateur basse fréquence (GBF), des fils de connexion.

\textbf{Protocole} :
\begin{enumerate}
\item Allumer l'oscilloscope sans rien brancher sur la voie Y.
\item Observer l'écran.
\item Brancher ensuite un GBF entre la borne Y et la masse, puis observer à nouveau.
\end{enumerate}

\textbf{Observations} :
\begin{itemize}
\item Sans tension appliquée (ou voie réglée sur la position « 0 »), on observe une \cle{ligne horizontale} au milieu de l'écran (ou un simple \cle{spot} lumineux si le balayage est coupé).
\item Avec le GBF branché, la ligne se transforme en une courbe qui monte et descend : c'est l'oscillogramme de la tension du GBF.
\end{itemize}

\textbf{Interprétation} : la ligne horizontale correspond à une tension nulle et constante : $u = \SI{0}{V}$ à chaque instant. Elle sert de \cle{référence} (le « zéro ») pour les mesures. Dès qu'une tension variable est appliquée, le spot se déplace verticalement au rythme de cette tension, et le balayage horizontal étale ce mouvement dans le temps : la courbe apparaît.
\end{experience}

\begin{popfigure}
\begin{center}
\begin{tikzpicture}[scale=1]
% GBF
\draw[fill=popgreenL,draw=popdark,line width=1pt,rounded corners=3pt] (0,0) rectangle (3.4,2.6);
\node[popdark,font=\small\bfseries] at (1.7,2.25) {GBF};
\draw[fill=white,draw=popdark] (1.7,1.5) circle (0.55);
\draw[poporange,line width=1pt] plot[domain=1.2:2.2,samples=40] (\x,{1.5+0.18*sin(deg(2*pi*(\x-1.2)/0.5))});
\draw[fill=red!70,draw=popdark] (1.0,0.45) circle (0.14);
\draw[fill=black!70,draw=popdark] (2.4,0.45) circle (0.14);
\node[popdark,font=\footnotesize] at (1.0,0.15) {$+$};
\node[popdark,font=\footnotesize] at (2.4,0.15) {$-$};
% Oscilloscope
\draw[fill=popcream,draw=popdark,line width=1pt,rounded corners=3pt] (6.2,0) rectangle (11,3.4);
\draw[fill=popblue!8,draw=popdark] (6.5,1.1) rectangle (9.3,3.1);
\foreach \x in {6.9,7.3,7.7,8.1,8.5,8.9}
  \draw[popline,thin] (\x,1.1) -- (\x,3.1);
\foreach \y in {1.5,1.9,2.3,2.7}
  \draw[popline,thin] (6.5,\y) -- (9.3,\y);
\draw[poporange,line width=1.2pt] plot[domain=6.5:9.3,samples=100] (\x,{2.1+0.6*sin(deg(2*pi*(\x-6.5)/1.9))});
\draw[fill=red!70,draw=popdark] (9.9,0.75) circle (0.14);
\node[popdark,font=\footnotesize] at (10.35,0.75) {Y};
\draw[fill=black!70,draw=popdark] (9.9,0.3) circle (0.14);
\node[popdark,font=\footnotesize] at (10.5,0.3) {M};
\node[popdark,font=\small\bfseries] at (8.6,3.6) {Oscilloscope};
% Fils de connexion
\draw[red!70,line width=1.4pt] (1.0,0.45) -- (1.0,-0.7) -- (9.9,-0.7) -- (9.9,0.61);
\draw[black!70,line width=1.4pt] (2.4,0.45) -- (2.4,-1.1) -- (10.6,-1.1) -- (10.6,0.3) -- (10.04,0.3);
\end{tikzpicture}
\end{center}
\legende{Branchement d'un GBF sur la voie Y d'un oscilloscope : la borne $+$ du GBF est reliée à l'entrée Y (fil rouge), la borne $-$ à la masse (fil noir). L'oscilloscope est branché en dérivation, comme un voltmètre.}
\end{popfigure}

\begin{methode}[Visualiser une tension à l'oscilloscope : les 5 étapes]
\begin{enumerate}
\item \textbf{Brancher hors tension} : relier la tension à étudier entre la borne Y et la borne masse de l'oscilloscope, circuit éteint.
\item \textbf{Régler le zéro} : sans tension appliquée, amener la ligne horizontale (ou le spot) sur l'axe horizontal central de l'écran à l'aide du bouton de position verticale.
\item \textbf{Choisir la sensibilité verticale} (V/div) : régler le bouton pour que la courbe occupe une grande partie de l'écran \emph{sans en sortir}. Si la courbe dépasse, augmenter les V/div.
\item \textbf{Choisir la durée de balayage} (ms/div) : régler le bouton pour observer un ou deux motifs complets de la courbe, bien étalés horizontalement.
\item \textbf{Lire les divisions} : compter le nombre de divisions verticales (pour la tension) et horizontales (pour la durée), puis multiplier par le calibre correspondant.
\end{enumerate}
\end{methode}

\courssub{Lecture d'un oscillogramme}

Lire un oscillogramme, c'est « traduire » le dessin de l'écran en valeurs de tension et de durée. La règle est toujours la même : \emph{compter des divisions, puis multiplier par le calibre}.

\begin{aretenir}[Les deux lectures fondamentales]
\begin{itemize}
\item \textbf{Hauteur verticale} $\rightarrow$ \textbf{valeur de la tension} :
\[ u = \text{sensibilité verticale (V/div)} \times \text{nombre de divisions verticales} \]
\item \textbf{Étendue horizontale} $\rightarrow$ \textbf{durée} :
\[ t = \text{durée de balayage (ms/div)} \times \text{nombre de divisions horizontales} \]
\end{itemize}
\end{aretenir}

\begin{exoresolu}[Lecture complète d'un oscillogramme]
Sur l'écran d'un oscilloscope, on visualise la tension délivrée par un GBF. Les réglages sont : sensibilité verticale $S_v = \SI{2}{V/div}$ ; durée de balayage $B = \SI{5}{ms/div}$. Le sommet de la courbe se situe à 3 divisions au-dessus de l'axe horizontal, et un motif complet de la courbe s'étend sur 4 divisions horizontales.

\begin{enumerate}
\item Quelle est la valeur de la tension au sommet de la courbe ?
\item Quelle est la durée d'un motif complet ?
\end{enumerate}
\tcblower
\textbf{Solution détaillée.}

\textbf{1. Tension au sommet.} La hauteur se lit verticalement : on utilise la sensibilité verticale.
\begin{align*}
u &= S_v \times y \\
u &= 2 \times 3 \\
u &= \SI{6}{V}
\end{align*}
La tension au sommet de la courbe vaut \SI{6}{V}.

\textbf{2. Durée d'un motif.} L'étendue se lit horizontalement : on utilise la durée de balayage.
\begin{align*}
t &= B \times x \\
t &= 5 \times 4 \\
t &= \SI{20}{ms}
\end{align*}
Un motif complet dure \SI{20}{ms}, soit \SI{0,02}{s}.

\textbf{Vérification des unités} : V/div $\times$ div donne bien des volts ; ms/div $\times$ div donne bien des millisecondes. Les divisions s'annulent : c'est un bon réflexe pour contrôler son calcul.
\end{exoresolu}

\courssub{Sécurité : un appareil précieux et des règles strictes}

L'oscilloscope est un appareil coûteux et délicat. Sa manipulation obéit à des consignes de sécurité qu'il faut connaître et respecter, pour ta protection et celle du matériel.

\begin{attention}[Consignes de sécurité avec l'oscilloscope]
\begin{itemize}
\item \textbf{Ne jamais brancher directement le secteur} (la prise de courant ENEO, \SI{220}{V}) sur un oscilloscope scolaire sans précaution particulière et sans l'autorisation du professeur : risque d'électrocution et de destruction de l'appareil.
\item \textbf{La borne « masse » de l'oscilloscope est reliée à la terre} par la prise de l'appareil. Un mauvais branchement peut donc créer un court-circuit. Respecte scrupuleusement les consignes du professeur et ne modifies jamais un montage sous tension.
\item \textbf{Vérifier les calibres} avant la mise sous tension : choisir d'abord une sensibilité verticale élevée (par exemple \SI{5}{V/div}), puis l'ajuster pour que la courbe reste sur l'écran.
\item \textbf{Manipuler hors tension, puis mettre sous tension} : on réalise toujours les branchements circuit éteint ; on ne met sous tension qu'après vérification du montage par le professeur.
\end{itemize}
\end{attention}

\begin{savaistu}[Pourquoi la masse est-elle reliée à la terre ?]
Dans les installations électriques, la \cle{terre} est un fil de protection relié au sol. Elle protège les utilisateurs en évacuant les courants de défaut. Comme la masse de l'oscilloscope est reliée à la terre, brancher sa borne de masse sur un point dangereux d'un circuit revient à relier ce point directement au sol : d'où le risque de court-circuit si l'on ne suit pas les consignes. C'est pour cela qu'au laboratoire, on utilise des générateurs basse fréquence (GBF) délivrant des tensions faibles et sans danger.
\end{savaistu}

\begin{exercice}[À toi de jouer]
Un élève visualise une tension avec les réglages suivants : sensibilité verticale $S_v = \SI{0,5}{V/div}$, durée de balayage $B = \SI{2}{ms/div}$.
\begin{enumerate}
\item Le spot se déplace de 4 divisions vers le haut. Quelle est la valeur de la tension ?
\item Deux sommets consécutifs de la courbe sont séparés de 5 divisions horizontales. Quelle durée les sépare ?
\item Citer deux consignes de sécurité à respecter lors de cette manipulation.
\end{enumerate}
\end{exercice}

\begin{corrige}[Exercice 1]
\textbf{1.} $u = S_v \times y = 0{,}5 \times 4 = \SI{2}{V}$.

\textbf{2.} $t = B \times x = 2 \times 5 = \SI{10}{ms}$.

\textbf{3.} Par exemple : réaliser les branchements hors tension et faire vérifier le montage par le professeur avant la mise sous tension ; ne jamais brancher le secteur directement sur l'oscilloscope ; vérifier les calibres avant de commencer.
\end{corrige}

En résumé, l'oscilloscope est l'instrument qui rend la tension visible : branché en dérivation entre la voie Y et la masse, il dessine la tension en fonction du temps grâce à deux réglages — la sensibilité verticale (V/div) pour les tensions et la durée de balayage (ms/div) pour les durées. Maintenant que nous savons lire un oscillogramme, nous pouvons l'utiliser pour étudier la tension alternative sinusoïdale : c'est ce que nous ferons en mesurant sa \emph{période} et sa \emph{fréquence} dans la prochaine section.

\coursec{Période et fréquence d'une tension sinusoïdale}

Dans la section précédente, tu as appris à brancher un oscilloscope et à observer sur son écran la courbe d'une tension. Lorsque nous avons branché le générateur basse fréquence (GBF), la courbe obtenue n'était pas n'importe quelle courbe : c'était une vague régulière, qui monte, descend, remonte, redescend\ldots{} toujours de la même façon. Regarde bien cette courbe : elle \emph{se répète}. C'est précisément cette répétition que nous allons maintenant exploiter pour définir deux grandeurs fondamentales de l'électricité : la \cle{période} et la \cle{fréquence}. Ce sont elles qui permettent de répondre à une question très concrète : pourquoi le réseau ENEO est-il annoncé comme un réseau « \SI{220}{V} -- \SI{50}{Hz} » ?

\courssub{Observation expérimentale : une courbe qui se répète}

\begin{experience}[Observer la répétition du motif à l'oscilloscope]
\textbf{Matériel :} un GBF réglé sur « signal sinusoïdal », un oscilloscope, deux fils de connexion.

\textbf{Protocole :}
\begin{enumerate}
\item Régler le GBF sur une fréquence de l'ordre de \SI{100}{Hz} et une amplitude moyenne.
\item Brancher la sortie du GBF sur la voie Y de l'oscilloscope.
\item Régler la durée de balayage (par exemple \SI{5}{ms/div}) et la sensibilité verticale jusqu'à obtenir une courbe stable et bien étalée.
\item Observer attentivement la courbe : repérer la plus petite partie de la courbe qui, recopiée bout à bout, permettrait de reconstruire tout le tracé.
\end{enumerate}

\textbf{Observations :} la courbe a la forme d'une vague régulière appelée \cle{sinusoïde}. On y distingue une portion qui se reproduit \emph{identiquement à elle-même} : par exemple, la portion qui part d'un passage par zéro en montant, atteint un sommet, redescend en passant par zéro, descend jusqu'à un creux, puis revient à zéro en montant. Cette portion se répète deux, trois, quatre fois sur l'écran.

\textbf{Interprétation :} la tension délivrée par le GBF est \cle{périodique} : elle reprend exactement les mêmes valeurs à intervalles de temps égaux. La plus petite portion de courbe qui se répète est appelée le \cle{motif élémentaire} (ou simplement \emph{motif}).
\end{experience}

\begin{aretenir}[Motif élémentaire]
Une tension périodique est une tension dont la courbe est formée d'un \cle{motif} qui se répète identique à lui-même à intervalles de temps égaux. Le motif élémentaire est la \textbf{plus petite} partie de la courbe qui se répète.
\end{aretenir}

\begin{popfigure}
\begin{center}
\begin{tikzpicture}
\begin{axis}[
    width=0.95\linewidth, height=5.2cm,
    axis lines=middle,
    xlabel={$t$}, ylabel={$u$},
    xlabel style={at={(ticklabel* cs:1)}, anchor=west},
    ylabel style={at={(ticklabel* cs:1)}, anchor=south},
    xtick=\empty, ytick=\empty,
    xmin=-0.3, xmax=4.6, ymin=-1.7, ymax=1.9,
    clip=false
]
\addplot[popcurve,domain=0:4.2,samples=300]{sin(deg(3.14159*x))};
% double flèche pour T
\draw[<->, thick, poporange] (axis cs:0,1.45) -- (axis cs:2,1.45);
\node[poporange, font=\bfseries] at (axis cs:1,1.65) {$T$};
\draw[<->, thick, poporange] (axis cs:2,1.45) -- (axis cs:4,1.45);
\node[poporange, font=\bfseries] at (axis cs:3,1.65) {$T$};
% lignes pointillées
\draw[dashed, popmuted] (axis cs:0,0) -- (axis cs:0,1.45);
\draw[dashed, popmuted] (axis cs:2,0) -- (axis cs:2,1.45);
\draw[dashed, popmuted] (axis cs:4,0) -- (axis cs:4,1.45);
\node[popmuted, font=\small, below] at (axis cs:0,0) {$t_0$};
\node[popmuted, font=\small, below] at (axis cs:2,0) {$t_0+T$};
\node[popmuted, font=\small, below] at (axis cs:4,0) {$t_0+2T$};
\end{axis}
\end{tikzpicture}
\end{center}
\legende{Tension sinusoïdale $u(t)$ tracée sur deux périodes. Le motif élémentaire (en bleu) se répète : la durée $T$ d'un motif est la même partout. La double flèche orange matérialise une période.}
\end{popfigure}

\courssub{La période $T$ : durée du motif}

L'intuition est simple : si le motif se répète, il met un certain \emph{temps} à se dérouler. Ce temps, c'est la période. Sur l'oscillogramme, plus le balayage est lent, plus le motif paraît étalé ; mais la durée réelle du motif, elle, ne change pas : elle ne dépend que de la tension étudiée.

\begin{definition}[Période d'une tension périodique]
La \cle{période} d'une tension périodique, notée $T$, est la \textbf{durée du plus petit motif} qui se répète.

Son unité dans le système international est la \textbf{seconde} (symbole : s). On utilise très souvent le sous-multiple : la \textbf{milliseconde} (ms), avec
\[ \SI{1}{ms} = \SI{0,001}{s} = 10^{-3}\ \text{s} \qquad \text{et} \qquad \SI{1}{s} = \SI{1000}{ms}. \]
\end{definition}

\begin{exemplebox}[La période, un temps très court]
Sur le réseau ENEO, la tension du secteur effectue un motif complet en \SI{20}{ms}, c'est-à-dire \SI{0,02}{s}. C'est environ 25 fois plus rapide qu'un battement de paupière ! C'est pourquoi une lampe alimentée par le secteur semble briller en continu, alors que sa tension s'annule 100 fois par seconde.
\end{exemplebox}

\courssub{Mesurer la période sur un oscillogramme}

Comment connaître $T$ concrètement ? L'oscilloscope est aussi un \emph{chronomètre} : l'axe horizontal de l'écran est gradué en temps grâce au réglage de la \cle{durée de balayage} (notée souvent $b$), affichée en ms/div ou s/div. Chaque division horizontale (chaque grand carreau de l'écran) correspond à cette durée.

\begin{propriete}[Mesure de la période sur l'écran]
Si un motif complet occupe $x$ divisions horizontales sur l'écran et si la durée de balayage vaut $b$, alors la période est :
\[ T = x \times b \]
avec $x$ en divisions (sans unité) et $b$ la durée de balayage (par exemple en ms/div), ce qui donne $T$ dans l'unité de $b$ (par exemple en ms).
\end{propriete}

\begin{popfigure}
\begin{center}
\begin{tikzpicture}[scale=0.72]
% écran
\draw[fill=popcream, rounded corners=4pt] (-0.6,-3.6) rectangle (10.6,3.6);
% grille
\draw[step=1, popline!60, thin] (0,-3) grid (10,3);
\draw[popline] (0,0) -- (10,0);
\draw[popline] (5,-3) -- (5,3);
% sinusoïde : période 4 divisions. sin(90*x) en degrés.
% Zéros montants à x=0, 4, 8. Sommet x=1. Creux x=3.
\draw[very thick, popblue, domain=0:10, samples=200] plot (\x, {2.2*sin(90*\x)});
% mesure d'une période : de x=0 à x=4 (passages par zéro montants)
\draw[dashed, popmuted] (0,0) -- (0,-3.2);
\draw[dashed, popmuted] (4,0) -- (4,-3.2);
\draw[<->, very thick, poporange] (0,-3.0) -- (4,-3.0);
\node[poporange, font=\small\bfseries, below] at (2,-3.05) {1 motif = 4 divisions};
% sommet et creux
\fill[popblue] (1,2.2) circle (2.5pt);
\fill[popblue] (3,-2.2) circle (2.5pt);
\node[popblue, font=\small, above] at (1,2.25) {sommet};
\node[popblue, font=\small, below] at (3,-2.25) {creux};
% réglage balayage
\node[draw=popdark, rounded corners=2pt, fill=white, font=\small, anchor=north west] at (-0.5,-3.8) {Balayage : $b = \SI{5}{ms/div}$};
\node[draw=poporange, rounded corners=2pt, fill=white, font=\small, anchor=north east] at (10.5,-3.8) {$T = 4 \times 5 = \SI{20}{ms}$};
\end{tikzpicture}
\end{center}
\legende{Mesure de la période sur l'écran de l'oscilloscope. Un motif complet (entre deux passages par zéro en montant, traits pointillés) occupe 4 divisions horizontales. Avec un balayage de \SI{5}{ms/div}, on obtient $T = 4 \times 5 = \SI{20}{ms}$.}
\end{popfigure}

\courssub{La fréquence $f$ : nombre de motifs par seconde}

Imaginons maintenant deux GBF : le premier produit une tension dont le motif dure \SI{20}{ms}, le second un motif de \SI{2}{ms}. Le second « vibre » dix fois plus vite : ses motifs se succèdent plus rapidement. Pour décrire cette rapidité de répétition, on introduit une nouvelle grandeur : la fréquence.

\begin{definition}[Fréquence d'une tension périodique]
La \cle{fréquence} d'une tension périodique, notée $f$, est le \textbf{nombre de périodes (de motifs) par seconde}.

Son unité est le \textbf{hertz} (symbole : Hz), en hommage au physicien allemand Heinrich Hertz. Une tension de fréquence $f = \SI{50}{Hz}$ effectue donc 50 motifs complets chaque seconde.

Multiples utiles : le kilohertz ($\SI{1}{kHz} = \SI{1000}{Hz}$) et le mégahertz ($\SI{1}{MHz} = 10^{6}\ \text{Hz}$).
\end{definition}

Raisonnons pour établir la relation entre $f$ et $T$. Si chaque motif dure $T$ secondes, alors en une seconde il se déroule exactement $\dfrac{1}{T}$ motifs. Par exemple, si $T = \SI{0,02}{s}$, il se déroule $\dfrac{1}{0,02} = 50$ motifs par seconde. D'où la propriété fondamentale :

\begin{propriete}[Relation entre fréquence et période]
La fréquence $f$ et la période $T$ d'une tension périodique sont liées par :
\[ f = \frac{1}{T} \qquad \text{et de manière équivalente} \qquad T = \frac{1}{f} \]
avec $T$ en \textbf{secondes} (s) et $f$ en \textbf{hertz} (Hz). Cette relation, admise comme définition de la fréquence, se vérifie numériquement sur n'importe quel oscillogramme. Elle est exigible au BEPC.
\end{propriete}

\begin{exemplebox}[Le réseau électrique camerounais]
La tension délivrée par ENEO dans les prises de courant a une fréquence $f = \SI{50}{Hz}$. Sa période vaut donc :
\[ T = \frac{1}{f} = \frac{1}{50} = \SI{0,02}{s} = \SI{20}{ms}. \]
C'est exactement la valeur mesurée sur l'oscillogramme de la figure précédente : la théorie et la mesure sont en accord.
\end{exemplebox}

\begin{savaistu}[Pourquoi 50 Hz ?]
Le choix de \SI{50}{Hz} pour les réseaux électriques d'Afrique, d'Europe et d'Asie est historique : c'est un compromis trouvé au début du XX\textsuperscript{e} siècle entre les pertes dans les lignes, le fonctionnement des moteurs et le confort d'éclairage. Les Amériques ont choisi \SI{60}{Hz}. Au laboratoire du collège, le GBF permet de produire des tensions de fréquences très variées, réglables de quelques hertz à quelques kilohertz (jusqu'à \SI{20}{kHz} environ), ce qui couvre la gamme des sons audibles par l'oreille humaine : c'est d'ailleurs ainsi qu'on alimente un haut-parleur en TP de physique !
\end{savaistu}

\courssub{Méthode : déterminer $T$ puis $f$ à partir d'un oscillogramme}

\begin{methode}[Exploiter un oscillogramme pour trouver la période et la fréquence]
\begin{enumerate}
\item \textbf{Repérer le motif} : identifier sur la courbe la plus petite portion qui se répète (par exemple entre deux passages par zéro dans le même sens, ou entre deux sommets consécutifs).
\item \textbf{Compter les divisions} : compter le nombre $x$ de divisions horizontales occupées par un motif complet.
\item \textbf{Calculer la période} : multiplier par la durée de balayage : $T = x \times b$.
\item \textbf{Convertir en secondes} si nécessaire : $T$ (en s) $= T$ (en ms) $\times 10^{-3}$.
\item \textbf{Calculer la fréquence} : appliquer $f = \dfrac{1}{T}$ avec $T$ en secondes ; le résultat est en hertz.
\item \textbf{Vérifier la cohérence} : une période plus petite doit donner une fréquence plus grande.
\end{enumerate}
\end{methode}

\begin{exoresolu}[Oscillogramme d'un GBF]
\textbf{Énoncé.} Sur l'écran d'un oscilloscope branché sur un GBF, on observe une tension sinusoïdale. Un motif complet occupe $x = 4$ divisions horizontales. La durée de balayage est réglée sur $b = \SI{5}{ms/div}$.
\begin{enumerate}
\item Déterminer la période $T$ de cette tension, en ms puis en s.
\item En déduire sa fréquence $f$.
\item Que devient la fréquence si, le GBF étant réglé différemment, un motif n'occupe plus que 2 divisions (avec le même balayage) ?
\end{enumerate}
\tcblower
\textbf{Solution détaillée.}
\begin{enumerate}
\item La période est la durée d'un motif :
\[ T = x \times b = 4 \times 5 = \SI{20}{ms}. \]
Conversion en secondes : $T = 20 \times 10^{-3} = \SI{0,02}{s}$.
\item On applique la relation fondamentale avec $T$ en secondes :
\[ f = \frac{1}{T} = \frac{1}{0,02} = \SI{50}{Hz}. \]
La tension du GBF a la même fréquence que celle du réseau ENEO.
\item Nouvelle période : $T' = 2 \times 5 = \SI{10}{ms} = \SI{0,01}{s}$. Nouvelle fréquence :
\[ f' = \frac{1}{T'} = \frac{1}{0,01} = \SI{100}{Hz}. \]
Le motif est deux fois plus court, la fréquence est deux fois plus grande : c'est cohérent, car $f$ et $T$ varient en sens inverse.
\end{enumerate}
\end{exoresolu}

\begin{attention}[L'erreur classique : les millisecondes oubliées]
L'erreur la plus fréquente au BEPC est d'appliquer $f = \dfrac{1}{T}$ avec $T$ \textbf{en millisecondes}. C'est faux ! La formule exige $T$ en \textbf{secondes}.

Exemple : si $T = \SI{20}{ms}$, écrire $f = \dfrac{1}{20} = \SI{0,05}{Hz}$ est une double erreur (valeur et physique : une tension de \SI{0,05}{Hz} mettrait 20 s à faire un motif !). Il faut d'abord convertir : $T = \SI{0,02}{s}$, puis $f = \dfrac{1}{0,02} = \SI{50}{Hz}$.

Autres pièges à éviter :
\begin{itemize}
\item Compter les divisions d'un \emph{demi-motif} (par exemple entre un sommet et un creux) au lieu d'un motif complet : cela donne une période deux fois trop petite.
\item Confondre la durée de balayage (réglage horizontal, en ms/div) avec la sensibilité verticale (en V/div), qui sert à mesurer les tensions, pas les durées.
\item Oublier l'unité dans la réponse finale : une fréquence s'exprime en hertz, une période en secondes.
\end{itemize}
\end{attention}

\begin{exercice}[À toi de jouer]
Sur l'écran d'un oscilloscope, un motif de tension sinusoïdale occupe 5 divisions horizontales avec un balayage de $\SI{2}{ms/div}$.
\begin{enumerate}
\item Calculer la période $T$ de cette tension en ms, puis en s.
\item Calculer sa fréquence $f$.
\item Cette tension est-elle celle du réseau ENEO ? Justifier.
\end{enumerate}
\end{exercice}

\begin{corrige}[Exercice : À toi de jouer]
\begin{enumerate}
\item $T = x \times b = 5 \times 2 = \SI{10}{ms}$, soit $T = \SI{0,01}{s}$.
\item $f = \dfrac{1}{T} = \dfrac{1}{0,01} = \SI{100}{Hz}$.
\item Non : le réseau ENEO a une fréquence de \SI{50}{Hz}. Ici $f = \SI{100}{Hz}$, soit le double : cette tension vient d'un autre générateur (par exemple un GBF).
\end{enumerate}
\end{corrige}

\begin{aretenir}[L'essentiel de la section]
\begin{itemize}
\item Une tension périodique est formée d'un \cle{motif} qui se répète identique à lui-même.
\item La \cle{période} $T$ est la durée d'un motif ; elle s'exprime en secondes (s) ou en millisecondes (ms).
\item Sur l'oscillogramme : $T = x \times b$ (nombre de divisions $\times$ durée de balayage).
\item La \cle{fréquence} $f$ est le nombre de périodes par seconde ; elle s'exprime en hertz (Hz).
\item Relation fondamentale : $f = \dfrac{1}{T}$, avec $T$ en secondes et $f$ en hertz.
\item Réseau camerounais : $f = \SI{50}{Hz}$, donc $T = \SI{20}{ms}$.
\end{itemize}
\end{aretenir}

\coursec{Tension maximale et tension efficace}

Dans la section précédente, nous avons appris à mesurer le \emph{rythme} d'une tension alternative sinusoïdale : sa période $T$ et sa fréquence $f$. Nous savons donc désormais que la tension du secteur ENEO « bat » 50 fois par seconde. Mais une question reste en suspens : quelle est la \emph{valeur} de cette tension ? Quand on dit « le secteur délivre \SI{220}{V} », de quelle valeur parle-t-on exactement, puisque la tension change à chaque instant, passant par zéro, montant, redescendant, devenant négative ? C'est toute l'ambition de cette section : définir deux grandeurs qui caractérisent la \emph{hauteur} de la sinusoïde, la \cle{tension maximale} et la \cle{tension efficace}, et découvrir la relation simple qui les relie.

\courssub{Observation : une tension qui varie entre deux bornes}

\begin{experience}[Observer les valeurs extrêmes d'une tension sinusoïdale]
\textbf{Matériel :} un générateur basse fréquence (GBF), un oscilloscope, des fils de connexion.

\textbf{Protocole :}
\begin{enumerate}
\item Régler le GBF sur « signal sinusoïdal » et le brancher aux bornes de la voie Y de l'oscilloscope.
\item Régler la sensibilité verticale sur \SI{2}{V/div} et la durée de balayage de façon à observer deux ou trois périodes.
\item Observer attentivement la courbe obtenue.
\end{enumerate}

\textbf{Observations :} la courbe monte jusqu'à un sommet, redescend en passant par zéro, descend jusqu'à un creux, puis remonte. Les sommets sont tous à la même hauteur au-dessus de l'axe horizontal ; les creux sont tous à la même profondeur en dessous. La hauteur d'un sommet est exactement égale à la profondeur d'un creux.

\textbf{Interprétation :} la tension $u(t)$ varie en permanence, mais elle reste toujours comprise entre deux valeurs extrêmes opposées : une valeur maximale positive, notée $U_{max}$, et une valeur minimale négative, égale à $-U_{max}$. La sinusoïde est donc « enfermée » dans une bande horizontale dont les bords sont $+U_{max}$ et $-U_{max}$.
\end{experience}

\begin{popfigure}
\begin{center}
\begin{tikzpicture}
\begin{axis}[
  width=0.95\linewidth, height=6.2cm,
  axis lines=middle,
  xlabel={$t$ (ms)}, ylabel={$u$ (V)},
  xlabel style={at={(ticklabel* cs:1)}, anchor=west},
  ylabel style={at={(ticklabel* cs:1)}, anchor=south},
  xmin=-0.5, xmax=42, ymin=-9.5, ymax=9.5,
  xtick={0,10,20,30,40},
  ytick={-8,-5.7,0,5.7,8},
  yticklabels={$-U_{max}$, $-U$, $0$, $U$, $U_{max}$},
  tick label style={font=\small},
  grid=none,
  clip=false
]
\addplot[popcurve,domain=0:40,samples=300]{8*sin(360*x/20)};
\addplot[dashed,poporange,thick,domain=0:40]{8};
\addplot[dashed,poporange,thick,domain=0:40]{-8};
\addplot[dashed,popgreen,thick,domain=0:40]{5.7};
\addplot[dashed,popgreen,thick,domain=0:40]{-5.7};
\node[poporange,anchor=west,font=\small] at (axis cs:40.5,8) {$U_{max}=8$ V};
\node[poporange,anchor=west,font=\small] at (axis cs:40.5,-8) {$-U_{max}$};
\node[popgreen,anchor=west,font=\small] at (axis cs:40.5,5.7) {$U \approx 5{,}7$ V};
\end{axis}
\end{tikzpicture}
\end{center}
\legende{Tension sinusoïdale de tension maximale $U_{max} = \SI{8}{V}$. La courbe reste comprise entre $+U_{max}$ et $-U_{max}$ (pointillés orange). Le niveau de la tension efficace $U = U_{max}/\sqrt{2} \approx \SI{5,7}{V}$ est repéré en pointillés verts.}
\end{popfigure}

\courssub{La tension maximale $U_{max}$}

\begin{definition}[Tension maximale]
La \cle{tension maximale} (ou \emph{amplitude}) d'une tension alternative sinusoïdale, notée $U_{max}$, est la \textbf{plus grande valeur} atteinte par la tension au cours du temps. Elle s'exprime en \textbf{volts} (V). La tension minimale vaut $-U_{max}$.
\end{definition}

La tension maximale se mesure à l'\textbf{oscilloscope} : on compte le nombre de divisions verticales entre l'axe horizontal (le zéro) et un sommet de la courbe, puis on multiplie par la sensibilité verticale choisie.

\begin{methode}[Mesurer $U_{max}$ à l'oscilloscope]
\begin{enumerate}
\item Repérer la position du zéro (trace horizontale quand aucune tension n'est appliquée).
\item Compter le nombre $n$ de divisions entre le zéro et un sommet de la sinusoïde.
\item Lire la sensibilité verticale $S_V$ affichée sur le bouton de réglage (en V/div).
\item Calculer : $U_{max} = n \times S_V$.
\end{enumerate}
\end{methode}

\begin{popfigure}
\begin{center}
\begin{tikzpicture}[scale=0.72]
% Cadre de l'écran
\draw[fill=popdark,rounded corners=4pt] (-0.6,-4.6) rectangle (9.6,4.6);
\draw[fill=black!85,rounded corners=2pt] (0,-4) rectangle (9,4);
% Grille
\foreach \x in {0,1,...,9} \draw[popline!40,thin] (\x,-4) -- (\x,4);
\foreach \y in {-4,-3,...,4} \draw[popline!40,thin] (0,\y) -- (9,\y);
\draw[popline!70] (0,0) -- (9,0);
\draw[popline!70] (4.5,-4) -- (4.5,4);
% Sinusoïde : amplitude 3 divisions
\draw[popgreen,very thick] plot[domain=0:9,samples=200] (\x,{3*sin(360*\x/4.5)});
% Mesure de Umax
\draw[<->,poporange,very thick] (9.35,0) -- (9.35,3) node[midway,right,font=\small\bfseries]{$n = 3$ div};
\fill[poporange] (1.125,3) circle (2.2pt);
\node[poporange,font=\small,anchor=south] at (1.125,3.15) {sommet};
% Annotations
\node[white,font=\small,anchor=west] at (-0.5,4.9) {Écran de l'oscilloscope};
\node[poporange,font=\small,anchor=west] at (0.1,-4.9) {Sensibilité verticale : $S_V = \SI{2}{V/div}$};
\node[popgreen,font=\small\bfseries,anchor=west] at (4.9,-4.9) {$U_{max} = 3 \times 2 = \SI{6}{V}$};
\end{tikzpicture}
\end{center}
\legende{Mesure de la tension maximale sur l'écran gradué : le sommet se situe à $n = 3$ divisions au-dessus du zéro. Avec une sensibilité $S_V = \SI{2}{V/div}$, on obtient $U_{max} = 3 \times 2 = \SI{6}{V}$.}
\end{popfigure}

\courssub{La tension efficace $U$ : le constat expérimental}

Branchons maintenant, sur le \emph{même} générateur, un voltmètre réglé en position \textbf{AC} (courant alternatif, symbole $\sim$).

\begin{experience}[Comparer l'oscilloscope et le voltmètre]
\textbf{Protocole :}
\begin{enumerate}
\item Régler le GBF pour obtenir à l'oscilloscope une sinusoïde de tension maximale $U_{max} = \SI{8}{V}$.
\item Sans toucher au réglage du GBF, brancher un voltmètre en position AC aux bornes du générateur.
\item Lire la valeur affichée par le voltmètre.
\end{enumerate}

\textbf{Observations :} le voltmètre indique une valeur \emph{constante}, environ \SI{5,7}{V}, alors que la tension varie en permanence entre $-8$ V et $+8$ V ! Cette valeur est \emph{plus petite} que $U_{max}$.

\textbf{Interprétation :} le voltmètre en position AC n'indique ni la tension maximale, ni la tension à un instant donné (qui change tout le temps), mais une valeur particulière appelée \cle{tension efficace}. C'est une sorte de « valeur moyenne utile » de la tension.
\end{experience}

\begin{definition}[Tension efficace]
La \cle{tension efficace} d'une tension alternative sinusoïdale, notée $U$, est la valeur indiquée par un \textbf{voltmètre branché en position AC} ($\sim$). Elle s'exprime en \textbf{volts} (V).

C'est la tension efficace qui caractérise l'\textbf{effet énergétique} de la tension : une tension alternative de tension efficace $U$ produit, dans une résistance, le \emph{même échauffement} qu'une tension continue de même valeur $U$.
\end{definition}

Pour bien comprendre cette idée d'« équivalence énergétique », imagine deux lampes identiques. La première est alimentée par une tension continue de \SI{12}{V} (une batterie) ; la deuxième est alimentée par une tension alternative sinusoïdale de tension efficace \SI{12}{V}. Les deux lampes brillent \emph{exactement de la même façon}, car elles reçoivent la même énergie. Pourtant, la tension alternative atteint en réalité des crêtes de près de \SI{17}{V} ! C'est pourquoi la tension efficace est la grandeur vraiment utile pour l'électrotechnique : c'est elle qui dit « combien la tension vaut » du point de vue de l'énergie.

\courssub{La relation entre $U_{max}$ et $U$}

Reprendre l'expérience précédente avec plusieurs réglages du GBF permet de dresser le tableau suivant :

\begin{center}
\begin{tabularx}{\linewidth}{|c|X|X|X|X|}
\hline
\rowcolor{popblueL}
$U_{max}$ mesurée à l'oscilloscope (V) & 2,0 & 4,0 & 6,0 & 8,0 \\
\hline
$U$ mesurée au voltmètre (V) & 1,4 & 2,8 & 4,2 & 5,7 \\
\hline
Rapport $U_{max}/U$ & \num{1,41} & \num{1,41} & \num{1,41} & \num{1,41} \\
\hline
\end{tabularx}
\end{center}

Le rapport $U_{max}/U$ est \emph{toujours le même} : environ 1,41, c'est-à-dire $\sqrt{2}$. Ce résultat expérimental se généralise.

\begin{propriete}[Relation entre tension maximale et tension efficace]
Pour une tension alternative \textbf{sinusoïdale}, la tension maximale et la tension efficace sont liées par la relation (admise en classe de 3ème, vérifiée expérimentalement) :
\[ U_{max} = U \times \sqrt{2} \qquad \text{soit} \qquad U = \frac{U_{max}}{\sqrt{2}} \]
avec $\sqrt{2} \approx 1,414$.
\end{propriete}

\begin{attention}[Cette relation n'est valable que pour une tension sinusoïdale]
La formule $U_{max} = U\sqrt{2}$ ne s'applique \textbf{que si la tension est sinusoïdale}. Pour un signal triangulaire ou rectangulaire, le rapport entre $U_{max}$ et $U$ est différent. Au BEPC, on te précisera toujours que la tension est sinusoïdale avant d'utiliser cette formule.
\end{attention}

\begin{exemplebox}[De l'oscilloscope au voltmètre]
Sur l'écran d'un oscilloscope, on mesure une tension maximale $U_{max} = \SI{8}{V}$. La tension efficace vaut alors :
\[ U = \frac{U_{max}}{\sqrt{2}} = \frac{8}{1,414} \approx \SI{5,7}{V} \]
Un voltmètre branché en position AC sur ce même générateur afficherait donc environ \SI{5,7}{V}.
\end{exemplebox}

\courssub{Application au réseau électrique : que signifie « 220 V -- 50 Hz » ?}

Regarde la plaque signalétique d'une bouilloire, d'un fer à repasser ou d'un chargeur de téléphone : tu y lis presque toujours « \SI{220}{V} -- \SI{50}{Hz} ». Nous pouvons maintenant décoder entièrement cette inscription :

\begin{itemize}
\item \textbf{\SI{220}{V}} est la \cle{tension efficace} $U$ du réseau pour lequel l'appareil est prévu ;
\item \textbf{\SI{50}{Hz}} est la \cle{fréquence} $f$ de cette tension (période $T = 1/f = \SI{0,02}{s} = \SI{20}{ms}$).
\end{itemize}

Mais alors, quelle est la tension \emph{maximale} du secteur ? Appliquons la relation :
\[ U_{max} = U \times \sqrt{2} = 220 \times 1,414 \approx \SI{311}{V} \]

La tension du secteur atteint donc en réalité $+\SI{311}{V}$ et $-\SI{311}{V}$ à chaque période, soit 100 fois par seconde (50 crêtes positives et 50 crêtes négatives) ! La valeur « \SI{220}{V} » n'est qu'une valeur efficace, celle qui rend compte de l'échauffement produit dans les appareils.

\begin{savaistu}[Pourquoi 220 V et pas 311 V sur les appareils ?]
C'est bien la \textbf{tension efficace} \SI{220}{V} qui figure sur les appareils électriques, car c'est elle qui détermine leur fonctionnement normal : puissance consommée, échauffement, luminosité. Pourtant, la tension du secteur ENEO culmine réellement à $\pm \SI{311}{V}$ à chaque instant de crête ! Cette distinction est capitale pour les fabricants de composants électroniques : un condensateur ou un isolant branché sur le secteur doit supporter \SI{311}{V}, et non \SI{220}{V}. C'est aussi pour cela que le secteur est si dangereux : même si l'on parle de « 220 volts », les pointes réelles dépassent 300 volts, largement de quoi provoquer une électrocution mortelle. Ne touche jamais une prise, un fil dénudé ou l'intérieur d'un appareil branché !
\end{savaistu}

\courssub{Exercice résolu et erreurs à éviter}

\begin{exoresolu}[Exploiter les deux appareils de mesure]
Un élève branche un GBF sur un oscilloscope réglé à la sensibilité $S_V = \SI{5}{V/div}$. Le sommet de la sinusoïde se situe à 2 divisions au-dessus du zéro.
\begin{enumerate}
\item Calculer la tension maximale $U_{max}$.
\item En déduire la valeur qu'indiquerait un voltmètre branché en position AC aux bornes du GBF.
\item Le même élève branche ensuite un appareil portant la mention « \SI{220}{V} -- \SI{50}{Hz} » sur le secteur. Calculer la tension maximale du secteur.
\end{enumerate}
\tcblower
\textbf{Solution détaillée.}
\begin{enumerate}
\item La tension maximale se calcule à partir de l'oscillogramme :
\[ U_{max} = n \times S_V = 2 \times 5 = \SI{10}{V} \]
\item Le voltmètre en position AC indique la tension efficace :
\[ U = \frac{U_{max}}{\sqrt{2}} = \frac{10}{1,414} \approx \SI{7,1}{V} \]
\item Pour le secteur, on connaît la tension efficace $U = \SI{220}{V}$, donc :
\[ U_{max} = U \times \sqrt{2} = 220 \times 1,414 \approx \SI{311}{V} \]
\end{enumerate}
\end{exoresolu}

\begin{attention}[Ne confonds jamais $U_{max}$ et $U$ !]
C'est l'erreur la plus fréquente au BEPC. Retiens bien :
\begin{itemize}
\item L'\textbf{oscilloscope} permet de mesurer la tension \textbf{maximale} : $U_{max} = n \times S_V$ (nombre de divisions $\times$ sensibilité).
\item Le \textbf{voltmètre en position AC} mesure la tension \textbf{efficace} $U$.
\item Les deux sont liées par $U_{max} = U\sqrt{2}$, donc $U_{max}$ est \emph{toujours plus grande} que $U$ (environ 1,4 fois plus grande).
\end{itemize}
Si un énoncé te donne « \SI{220}{V} » pour le secteur, il s'agit de la tension \emph{efficace} : il faut multiplier par $\sqrt{2}$ pour obtenir $U_{max}$. À l'inverse, si l'on te donne une mesure faite à l'oscilloscope, il s'agit de $U_{max}$ : il faut diviser par $\sqrt{2}$ pour obtenir $U$. Vérifie toujours la cohérence de ta réponse : $U_{max} > U$.
\end{attention}

\begin{aretenir}[L'essentiel de la section]
\begin{itemize}
\item Une tension sinusoïdale varie entre $+U_{max}$ et $-U_{max}$.
\item La \cle{tension maximale} $U_{max}$ est la plus grande valeur atteinte ; elle se mesure à l'\textbf{oscilloscope} : $U_{max} = n \times S_V$.
\item La \cle{tension efficace} $U$ est la valeur indiquée par le \textbf{voltmètre en position AC} ; elle caractérise l'effet énergétique (échauffement) de la tension.
\item Pour une tension sinusoïdale : $U_{max} = U\sqrt{2}$, soit $U = U_{max}/\sqrt{2}$, avec $\sqrt{2} \approx 1,414$.
\item « \SI{220}{V} -- \SI{50}{Hz} » signifie : tension efficace \SI{220}{V} (soit $U_{max} \approx \SI{311}{V}$) et fréquence \SI{50}{Hz}.
\end{itemize}
\end{aretenir}

\begin{exercice}[À toi de jouer]
\begin{enumerate}
\item Sur un oscillogramme obtenu avec la sensibilité $S_V = \SI{2}{V/div}$, le sommet de la sinusoïde est à 3,5 divisions au-dessus du zéro. Calculer $U_{max}$ puis la tension efficace $U$.
\item Un voltmètre en position AC branché sur une prise indique \SI{220}{V}. Quelle est la valeur minimale atteinte par la tension du secteur ?
\item Un élève affirme : « Le voltmètre indique \SI{6}{V}, donc la tension maximale est \SI{6}{V}. » Corriger son erreur et donner la bonne valeur de $U_{max}$.
\end{enumerate}
\end{exercice}

\begin{corrige}[Exercice : À toi de jouer]
\begin{enumerate}
\item $U_{max} = n \times S_V = 3,5 \times 2 = \SI{7}{V}$. Puis $U = U_{max}/\sqrt{2} = 7/1,414 \approx \SI{5,0}{V}$.
\item La valeur minimale est $-U_{max} = -U\sqrt{2} = -220 \times 1,414 \approx -\SI{311}{V}$.
\item Le voltmètre en position AC mesure la tension \emph{efficace}, pas la tension maximale. La tension maximale vaut $U_{max} = U\sqrt{2} = 6 \times 1,414 \approx \SI{8,5}{V}$.
\end{enumerate}
\end{corrige}

\coursec{Exploitation complète d'un oscillogramme et sécurité}

Dans la section précédente, tu as découvert deux façons de caractériser une tension alternative sinusoïdale : la \cle{tension maximale} $U_{\text{max}}$, qui se lit directement sur l'écran de l'oscilloscope, et la \cle{tension efficace} $U$, celle qu'indique le voltmètre et qui produit le même échauffement qu'une tension continue de même valeur. Tu as aussi retenu la relation fondamentale :

\[ U = \dfrac{U_{\text{max}}}{\sqrt{2}} \]

Mais dans la pratique, au BEPC comme au laboratoire, on ne te donne pas les valeurs : on te montre un \cle{oscillogramme} et c'est à toi d'en extraire toutes les informations. Cette dernière section te donne la méthode complète, pas à pas, puis elle se termine par un sujet capital : la sécurité électrique, car la tension du secteur n'est pas un jouet.

\courssub{Méthode : exploiter un oscillogramme en quatre étapes}

Un oscillogramme, c'est une photographie de la tension au cours du temps. Pour l'exploiter, il faut d'abord connaître les \cle{réglages} de l'appareil : sans eux, les divisions de l'écran ne veulent rien dire. Deux réglages sont indispensables :

\begin{itemize}
\item la \cle{sensibilité verticale} (ou calibre vertical), notée $S_v$, exprimée en volts par division (V/div) : elle indique combien de volts représente une division verticale ;
\item le \cle{balayage} (ou base de temps), noté $B$, exprimé en millisecondes par division (ms/div) : il indique quelle durée représente une division horizontale.
\end{itemize}

\begin{methode}[Exploiter un oscillogramme : les quatre étapes]
Pour déterminer toutes les caractéristiques d'une tension sinusoïdale à partir de son oscillogramme, procède toujours dans le même ordre :
\begin{enumerate}
\item \textbf{Lire les réglages} : note la sensibilité verticale $S_v$ (en V/div) et le balayage $B$ (en ms/div). Ces deux informations sont toujours données dans l'énoncé ou affichées sur l'écran.
\item \textbf{Mesurer la tension maximale} : compte le nombre de divisions $y$ entre l'axe horizontal (position de repos) et le sommet de la courbe, puis calcule $U_{\text{max}} = y \times S_v$.
\item \textbf{Mesurer la période} : compte le nombre de divisions $x$ correspondant à un motif complet (par exemple entre deux sommets consécutifs), puis calcule $T = x \times B$. Convertis le résultat en secondes.
\item \textbf{Calculer la fréquence et la tension efficace} : $f = \dfrac{1}{T}$ (avec $T$ en secondes, $f$ en hertz) et $U = \dfrac{U_{\text{max}}}{\sqrt{2}}$.
\end{enumerate}
\end{methode}

\begin{attention}[Les pièges classiques de l'exploitation d'un oscillogramme]
\begin{itemize}
\item Ne confonds jamais l'\cle{amplitude} (distance entre l'axe et un sommet) avec la hauteur totale crête-à-crête (distance entre le sommet haut et le sommet bas). La tension maximale se mesure à partir de l'axe, pas entre les deux crêtes !
\item Pour la période, mesure un \textbf{motif complet} : d'un sommet au sommet suivant, ou d'un passage par zéro dans un sens au passage par zéro suivant dans le même sens.
\item Convertis toujours la période en \textbf{secondes} avant de calculer la fréquence : $T = \SI{20}{ms} = \SI{0,020}{s}$. Si tu oublies la conversion, ta fréquence sera mille fois trop grande !
\item Arrondis avec bon sens : $\sqrt{2} \approx 1,414$, donc $\dfrac{1}{\sqrt{2}} \approx 0,707$.
\end{itemize}
\end{attention}

\courssub{Étude de cas guidée : l'oscillogramme d'une tension de GBF}

Au laboratoire du collège, le professeur branche un \cle{GBF} (générateur basse fréquence) sur la voie A de l'oscilloscope. L'écran affiche une belle sinusoïde. Voici les réglages et les mesures :

\begin{itemize}
\item sensibilité verticale : $S_v = \SI{5}{V/div}$ ;
\item balayage : $B = \SI{2}{ms/div}$ ;
\item amplitude observée : $y = 2$ divisions ;
\item un motif complet s'étend sur $x = 10$ divisions.
\end{itemize}

Appliquons la méthode des quatre étapes.

\textbf{Étape 1 : les réglages.} Ils sont lus : $S_v = \SI{5}{V/div}$ et $B = \SI{2}{ms/div}$.

\textbf{Étape 2 : la tension maximale.} La courbe monte de 2 divisions au-dessus de l'axe :
\[ U_{\text{max}} = y \times S_v = 2 \times 5 = \SI{10}{V} \]

\textbf{Étape 3 : la période.} Un motif complet occupe 10 divisions horizontales :
\[ T = x \times B = 10 \times 2 = \SI{20}{ms} = \SI{0,020}{s} \]

\textbf{Étape 4 : la fréquence et la tension efficace.}
\begin{align*}
f &= \dfrac{1}{T} = \dfrac{1}{0,020} = \SI{50}{Hz} \\
U &= \dfrac{U_{\text{max}}}{\sqrt{2}} = \dfrac{10}{1,414} \approx \SI{7,1}{V}
\end{align*}

\begin{popfigure}
\begin{center}
\begin{tikzpicture}
\begin{axis}[
  width=0.92\linewidth, height=6.2cm,
  axis lines=middle,
  xlabel={$t$ (ms)}, ylabel={$u$ (V)},
  xlabel style={at={(ticklabel* cs:1)}, anchor=west},
  ylabel style={at={(ticklabel* cs:1)}, anchor=south},
  xmin=-2, xmax=26, ymin=-13, ymax=15,
  xtick={0,5,10,15,20,25},
  ytick={-10,-7.1,0,7.1,10},
  yticklabels={$-U_{\text{max}}$,$-U$,$0$,$U$,$U_{\text{max}}$},
  grid=both, grid style={popline!40},
  clip=false
]
\addplot[popcurve,domain=0:25,samples=200]{10*sin(360*x/20)};
\addplot[dashed,poporange,thick] coordinates {(-2,10) (25,10)};
\addplot[dashed,popgreen,thick] coordinates {(-2,7.1) (25,7.1)};
\draw[<->,popblue,thick] (axis cs:22,0) -- (axis cs:22,10) node[midway,right,popblue]{\small $U_{\text{max}}=10$ V};
\draw[<->,popdark,thick] (axis cs:5,12.5) -- (axis cs:25,12.5) node[midway,above,popdark]{\small $T = 20$ ms};
\addplot[mark=*,popink] coordinates {(5,10)};
\addplot[mark=*,popink] coordinates {(25,10)};
\node[popgreen,anchor=west] at (axis cs:25.5,7.1) {\small $U \approx 7,1$ V};
\node[poporange,anchor=west] at (axis cs:25.5,10) {\small $U_{\text{max}}$};
\end{axis}
\end{tikzpicture}
\end{center}
\legende{Oscillogramme récapitulatif de la tension du GBF : la sinusoïde (bleu) atteint la tension maximale $U_{\text{max}} = \SI{10}{V}$, le motif se répète toutes les $T = \SI{20}{ms}$, et la tension efficace $U \approx \SI{7,1}{V}$ (pointillés verts) se situe en dessous de $U_{\text{max}}$.}
\end{popfigure}

\courssub{Vérification expérimentale : l'oscilloscope face au voltmètre}

Comment être sûr que la relation $U = \dfrac{U_{\text{max}}}{\sqrt{2}}$ est correcte ? Par l'expérience, bien sûr ! C'est la démarche d'investigation : on confronte la théorie à la mesure.

\begin{experience}[Comparaison oscilloscope -- voltmètre]
\textbf{Matériel} : un GBF, un oscilloscope, un voltmètre, des fils de connexion.

\textbf{Protocole} :
\begin{enumerate}
\item Règle le GBF pour obtenir l'oscillogramme de l'étude de cas précédente ($U_{\text{max}} = \SI{10}{V}$, $f = \SI{50}{Hz}$).
\item Sans modifier le GBF, branche le voltmètre \textbf{en parallèle} en \textbf{position AC} (courant alternatif, symbole $\sim$) aux bornes du GBF.
\item Lis la valeur affichée par le voltmètre et compare-la au calcul $U = \dfrac{U_{\text{max}}}{\sqrt{2}}$.
\end{enumerate}

\textbf{Observations} : le voltmètre indique environ $\SI{7,1}{V}$, alors que l'oscilloscope montre une tension maximale de $\SI{10}{V}$.

\textbf{Interprétation} : les deux appareils ne mesurent pas la même chose. L'oscilloscope visualise la tension instantanée et permet de lire $U_{\text{max}}$ ; le voltmètre en position AC indique directement la \cle{tension efficace} $U$. La valeur mesurée ($\approx \SI{7,1}{V}$) correspond au calcul $\dfrac{10}{\sqrt{2}} \approx \SI{7,1}{V}$ : la relation $U = \dfrac{U_{\text{max}}}{\sqrt{2}}$ est vérifiée expérimentalement.

\textbf{Conclusion} : pour une tension sinusoïdale, la tension efficace indiquée par le voltmètre est égale à la tension maximale divisée par $\sqrt{2}$.
\end{experience}

\begin{exoresolu}[Exploitation complète d'un oscillogramme]
Sur l'écran d'un oscilloscope, on observe la tension délivrée par un GBF. Réglages : sensibilité verticale $S_v = \SI{5}{V/div}$, balayage $B = \SI{2}{ms/div}$. La courbe culmine à 2 divisions au-dessus de l'axe et un motif complet occupe 10 divisions.
\begin{enumerate}
\item Détermine la tension maximale $U_{\text{max}}$.
\item Détermine la période $T$ puis la fréquence $f$.
\item Calcule la tension efficace $U$ et dis ce qu'indiquerait un voltmètre branché aux bornes du GBF.
\end{enumerate}
\tcblower
\textbf{Solution détaillée.}
\begin{enumerate}
\item Tension maximale : $U_{\text{max}} = y \times S_v = 2 \times 5 = \SI{10}{V}$.
\item Période : $T = x \times B = 10 \times 2 = \SI{20}{ms} = \SI{0,020}{s}$.

Fréquence : $f = \dfrac{1}{T} = \dfrac{1}{0,020} = \SI{50}{Hz}$.
\item Tension efficace : $U = \dfrac{U_{\text{max}}}{\sqrt{2}} = \dfrac{10}{1,414} \approx \SI{7,1}{V}$.

Le voltmètre en position AC indiquerait environ $\SI{7,1}{V}$, car il mesure la tension efficace.
\end{enumerate}
\end{exoresolu}

\begin{exercice}[À toi de jouer]
Un oscillogramme est obtenu avec les réglages suivants : $S_v = \SI{2}{V/div}$ et $B = \SI{5}{ms/div}$. L'amplitude est de 3 divisions et un motif occupe 4 divisions. Détermine $U_{\text{max}}$, $T$, $f$ et $U$.
\end{exercice}

\begin{corrige}[Exercice : À toi de jouer]
$U_{\text{max}} = 3 \times 2 = \SI{6}{V}$ ; $T = 4 \times 5 = \SI{20}{ms} = \SI{0,020}{s}$ ; $f = \dfrac{1}{0,020} = \SI{50}{Hz}$ ; $U = \dfrac{6}{\sqrt{2}} \approx \SI{4,2}{V}$.
\end{corrige}

\courssub{Les dangers de la tension du secteur}

Jusqu'ici, nous avons travaillé avec le GBF du laboratoire, qui délivre des tensions de quelques volts : sans danger. Mais la tension du \cle{secteur}, celle des prises murales alimentées par ENEO, est toute autre : sa valeur efficace est d'environ $\SI{220}{V}$ et sa fréquence de $\SI{50}{Hz}$. Sa tension maximale atteint même $U_{\text{max}} = 220 \times \sqrt{2} \approx \SI{311}{V}$ !

\begin{attention}[Le secteur 220 V peut tuer]
Une tension efficace de $\SI{220}{V}$ est \textbf{mortelle} pour l'être humain. Le passage d'un courant électrique à travers le corps provoque une \cle{électrisation} : brûlures, tétanisation des muscles (impossibilité de lâcher le conducteur), arrêt de la respiration, fibrillation du cœur. Le danger dépend de l'intensité qui traverse le corps, mais dès quelques dizaines de milliampères, le cœur est menacé. C'est pourquoi \textbf{toutes les expériences de la classe utilisent uniquement le GBF} (quelques volts) et jamais la tension du secteur. Ne reproduis jamais chez toi un montage branché sur une prise murale.
\end{attention}

Voici les consignes de sécurité à respecter absolument, à la maison comme au laboratoire :

\begin{aretenir}[Consignes de sécurité électrique]
\begin{itemize}
\item Ne jamais toucher un fil \cle{dénudé} (dont l'isolant est arraché) ni un appareil électrique ouvert.
\item Ne jamais utiliser un appareil électrique les \textbf{mains mouillées} ou les pieds dans l'eau : l'eau rend le corps beaucoup plus conducteur.
\item Ne jamais insérer d'objet (tournevis, fourchette, trombone) dans les trous d'une prise de courant.
\item Ne pas surcharger une prise ou une rallonge avec trop d'appareils branchés en même temps.
\item En cas de problème (odeur de brûlé, étincelles, fil abîmé), ne pas intervenir soi-même : \textbf{faire appel à un adulte ou à un technicien qualifié}.
\end{itemize}
\end{aretenir}

\begin{popfigure}
\begin{center}
\begin{tikzpicture}[scale=1]
% Mur
\fill[popcream] (-5.2,-3.4) rectangle (5.2,3.4);
% Plaque de la prise
\fill[white,draw=popline,rounded corners=4pt,line width=1pt] (-1.6,-2.2) rectangle (1.6,2.2);
% Corps de la prise
\fill[popblueL,draw=popblue,line width=1.2pt] (0,0) circle (1.25);
\fill[white,draw=popblue,line width=1pt] (0,0) circle (0.95);
% Trous
\fill[popdark] (-0.4,0.15) circle (0.13);
\fill[popdark] (0.4,0.15) circle (0.13);
% Terre
\draw[popdark,line width=2pt] (0,0.95) -- (0,1.15);
\node[popmuted] at (0,-1.75) {\small Prise 220 V -- 50 Hz};
% Pictogramme 1 : danger électrique
\begin{scope}[shift={(-3.6,1.6)}]
\draw[popgold,fill=popgoldL,line width=1.2pt] (0,1) -- (-1.05,-0.8) -- (1.05,-0.8) -- cycle;
\draw[popdark,fill=popdark] (0.08,0.55) -- (-0.28,-0.05) -- (0.02,-0.05) -- (-0.08,-0.6) -- (0.28,0.05) -- (-0.02,0.05) -- cycle;
\node[popdark] at (0,-1.25) {\small Danger};
\end{scope}
% Pictogramme 2 : interdiction mains mouillées
\begin{scope}[shift={(-3.6,-1.3)}]
\draw[popink,fill=white,line width=1.6pt] (0,0) circle (0.85);
\draw[popink,line width=1.6pt] (-0.6,0.6) -- (0.6,-0.6);
\draw[popblue,line width=1.4pt] (0,0.45) .. controls (-0.3,0.05) and (-0.25,-0.25) .. (0,-0.45) .. controls (0.25,-0.25) and (0.3,0.05) .. (0,0.45);
\node[popdark] at (0,-1.25) {\small Mains mouillées};
\end{scope}
% Pictogramme 3 : fusible/disjoncteur
\begin{scope}[shift={(3.6,1.6)}]
\draw[popgreen,fill=popgreenL,line width=1.2pt,rounded corners=3pt] (-0.9,-0.7) rectangle (0.9,0.7);
\draw[popgreen,line width=1.4pt] (-1.3,0) -- (-0.9,0) (0.9,0) -- (1.3,0);
\draw[popgreen,line width=1.4pt] (-0.5,0) -- (-0.25,0.25) -- (0,0) -- (0.25,-0.25) -- (0.5,0);
\node[popdark] at (0,-1.25) {\small Fusible};
\end{scope}
% Pictogramme 4 : appeler un adulte
\begin{scope}[shift={(3.6,-1.3)}]
\draw[poppurple,fill=poppurpleL,line width=1.2pt] (0,0) circle (0.85);
\node[poppurple] at (0,0) {\Large\textbf{!}};
\node[popdark] at (0,-1.25) {\small Alerter un adulte};
\end{scope}
\end{tikzpicture}
\end{center}
\legende{Une prise murale du secteur (220 V -- 50 Hz) entourée des quatre réflexes de sécurité : reconnaître le danger électrique, ne jamais manipuler les mains mouillées, compter sur les fusibles et disjoncteurs, et alerter un adulte en cas de problème.}
\end{popfigure}

\courssub{Les dispositifs de protection : fusibles et disjoncteurs}

Heureusement, l'installation électrique d'une maison n'est pas laissée sans défense. Tu l'as vu dans la leçon sur les installations domestiques : deux dispositifs veillent en permanence.

\begin{definition}[Fusible et disjoncteur]
Un \cle{fusible} est un petit fil conducteur calibré, placé en série dans le circuit. Si l'intensité dépasse une valeur limite (court-circuit ou surcharge), il \textbf{fond} et coupe le circuit. Un \cle{disjoncteur} joue le même rôle de protection, mais il est \textbf{réarmable} : après avoir « sauté », il suffit de le réenclencher une fois la panne réparée. Le \cle{disjoncteur différentiel}, installé en tête de l'installation, protège en plus les personnes contre les fuites de courant vers la terre.
\end{definition}

Ces appareils protègent les installations et les personnes, mais ils ne dispensent pas de prudence : un fusible ne réagit qu'à une intensité trop grande, alors qu'une intensité déjà dangereuse pour le corps humain peut passer sans le faire fondre. La première protection, c'est ton comportement.

\begin{savaistu}[220 V ou 110 V ? Cela dépend des pays]
Tous les pays n'ont pas choisi le même secteur ! Au Cameroun et dans la plupart des pays d'Europe et d'Afrique, le secteur délivre $\SI{220}{V}$ à $\SI{50}{Hz}$. Mais aux États-Unis, au Canada ou au Japon, on trouve du $\SI{110}{V}$ à $\SI{60}{Hz}$. Un appareil conçu pour le 220 V branché sur du 110 V fonctionnera mal (une lampe éclairera faiblement), et surtout un appareil prévu pour 110 V branché sur du 220 V sera \textbf{endommagé}, voire détruit, car il recevra une tension deux fois trop grande. C'est pourquoi les voyageurs vérifient toujours la tension indiquée sur la plaque signalétique de leurs appareils, et utilisent parfois un adaptateur-transformateur.
\end{savaistu}

\courssub{Pour aller plus loin : la valeur moyenne d'une tension sinusoïdale (hors programme)}

Voici un complément qui t'aidera à comprendre un comportement curieux du voltmètre. \emph{Cette notion dépasse le programme strict de la classe de 3ème, mais elle éclaire la suite.}

Observe la sinusoïde de la figure récapitulative : pendant une demi-période, la tension est positive ; pendant l'autre demi-période, elle est négative, et \textbf{exactement symétrique}. Les deux parties se compensent parfaitement. On dit que la \cle{valeur moyenne} d'une tension sinusoïdale sur une période complète est \textbf{nulle} :

\[ U_{\text{moyenne}} = \SI{0}{V} \]

Conséquence pratique : si tu branches un voltmètre en position \textbf{DC} (courant continu, symbole $=$) aux bornes d'un GBF ou du secteur, il indiquera $\SI{0}{V}$, alors que la tension existe bel et bien ! Ce n'est pas une panne : le voltmètre en position DC mesure la valeur moyenne, qui est nulle pour une tension alternative sinusoïdale. Pour mesurer une tension alternative, il faut donc toujours choisir la position \textbf{AC} ($\sim$), qui donne la tension efficace.

\begin{attention}[Bien choisir la position du voltmètre]
Position \textbf{DC} ($=$) pour les tensions continues (pile, batterie) ; position \textbf{AC} ($\sim$) pour les tensions alternatives (GBF, secteur). Un voltmètre en position DC branché sur une tension alternative indique $0$ V : cela ne signifie pas qu'il n'y a pas de tension, mais que sa valeur moyenne est nulle. Inversement, en position AC sur une pile, l'indication est fausse. Réfléchis toujours à la nature de la tension avant de brancher.
\end{attention}

\begin{aretenir}[L'essentiel de la section]
\begin{itemize}
\item Pour exploiter un oscillogramme : 1) lire les réglages ($S_v$ et $B$), 2) mesurer $U_{\text{max}} = y \times S_v$, 3) mesurer $T = x \times B$, 4) calculer $f = \dfrac{1}{T}$ et $U = \dfrac{U_{\text{max}}}{\sqrt{2}}$.
\item Le voltmètre en position AC indique la tension efficace $U$ ; l'oscilloscope permet de lire $U_{\text{max}}$. Les deux sont liés par $U = \dfrac{U_{\text{max}}}{\sqrt{2}}$, vérifié par l'expérience.
\item La tension du secteur ($\SI{220}{V}$ efficaces, $\SI{50}{Hz}$) est \textbf{dangereuse, voire mortelle}. En classe, on utilise uniquement le GBF.
\item Sécurité : jamais de fils dénudés, jamais les mains mouillées, jamais d'objets dans les prises ; en cas de problème, alerter un adulte ou un technicien.
\item Les fusibles et les disjoncteurs protègent l'installation en coupant le circuit en cas de surcharge ou de court-circuit.
\item (Hors programme) La valeur moyenne d'une tension sinusoïdale sur une période est nulle : un voltmètre en position DC indique alors $\SI{0}{V}$.
\end{itemize}
\end{aretenir}

Tu sais désormais lire un oscillogramme comme un technicien et, surtout, tu connais les gestes qui protègent. Retiens cette idée simple : l'électricité est une formidable servante à la maison, dans les ateliers et les hôpitaux, à condition de la respecter. Cette leçon sur les tensions alternatives sinusoïdales t'a donné les clés pour comprendre le courant qui alimente tout le Cameroun, de la SONARA aux prises de ta salle de classe.

\newpage

\coursec{Activité d'intégration}

\coursec{Les tensions alternatives sinusoïdales}

\courssub{Objectifs de l'activité}

À la fin de cette activité, tu seras capable de :
\begin{itemize}
\item exploiter un \cle{oscillogramme} pour déterminer la \cle{tension maximale} $U_{max}$, la \cle{période} $T$ et la \cle{fréquence} $f$ d'une tension alternative sinusoïdale ;
\item calculer la \cle{tension efficace} $U$ et la comparer à la valeur lue sur un voltmètre en position AC ;
\item vérifier expérimentalement la relation $U_{max} = U\sqrt{2}$ ;
\item confronter des mesures aux normes du réseau électrique camerounais (\SI{220}{V} -- \SI{50}{Hz}) et conclure sur la conformité d'une installation ;
\item citer et justifier les consignes de sécurité liées aux mesures électriques.
\end{itemize}

\courssub{Situation}

\textbf{La tension du quartier est-elle normale ?}

Dans le quartier Nkolndongo à Yaoundé, les élèves du club scientifique du collège remarquent que certaines ampoules grillent souvent et que la bouilloire de la cantine met parfois très longtemps à chauffer. Leur professeur de PCT leur propose de vérifier si la tension électrique délivrée dans le quartier est conforme aux normes du réseau ENEO : \SI{220}{V} (tension efficace) et \SI{50}{Hz} (fréquence).

Pour travailler en toute sécurité, le professeur ne branche pas directement l'oscilloscope sur la prise murale : il utilise un \cle{GBF} (générateur basses fréquences) réglé pour \emph{simuler} le secteur, c'est-à-dire pour délivrer une tension alternative sinusoïdale de mêmes caractéristiques que celle du réseau. Il remet à chaque groupe trois documents.

\begin{popfigure}
\begin{tikzpicture}[scale=0.85, >=Stealth]
% Grille (1 cm = 1 division)
\draw[popline, step=1cm] (0,0) grid (10,6);
% Axes
\draw[->, thick, popdark] (0,3) -- (10.3,3) node[right] {$t$ (div)};
\draw[->, thick, popdark] (0,0) -- (0,6.3) node[above] {$U$ (div)};
% Ligne zéro (axe horizontal central)
\draw[dashed, popmuted, thick] (0,3) -- (10,3);
% Courbe sinusoïdale : période 4 div, amplitude 3.1 div, centrée sur y=3
% y = 3 + 3.1 * sin( (2*pi/4) * x ) = 3 + 3.1 * sin(pi/2 * x)
\draw[popcurve, very thick, domain=0:10, samples=300] plot (\x, {3 + 3.1*sin(90*\x)});
% Annotation Amplitude (verticale à x=0.5)
\draw[<->, poporange, thick] (0.5,3) -- (0.5,6.1) node[midway, right=2pt, poporange, font=\bfseries] {$3,1$ div};
\draw[poporange, dashed, thin] (0.5,6.1) -- (0,6.1);
% Annotation Période (horizontale en bas)
\draw[<->, popblue, thick] (0,2.3) -- (4,2.3) node[midway, below=2pt, popblue, font=\bfseries] {$4$ div};
\draw[popblue, dashed, thin] (0,2.3) -- (0,2);
\draw[popblue, dashed, thin] (4,2.3) -- (4,2);
% Repère du sommet
\fill[popink] (1,6.1) circle (3pt) node[above right, popink, font=\small] {Sommet ($y=3,1$)};
% Cadre des réglages
\node[align=left, anchor=north east, fill=popcream, draw=poporange, thick, rounded corners, inner sep=6pt] at (9.8, 5.9) {
    \textbf{Réglages de l'oscilloscope} \\
    Sensibilité verticale : $S_v = \SI{100}{V/div}$ \\
    Balayage (base de temps) : $B = \SI{5}{ms/div}$
};
\end{tikzpicture}
\legende{Oscillogramme visualisé sur l'écran : une sinusoïde d'amplitude $3,1$ divisions et de période $4$ divisions.}
\end{popfigure}

\textbf{Document 1 : l'oscillogramme.} L'oscilloscope est branché aux bornes du GBF. Les réglages sont ceux indiqués sur la figure ci-dessus. On observe une sinusoïde dont le sommet se situe à $3,1$ divisions au-dessus de l'axe horizontal, et dont un motif complet (une période) s'étend sur $4$ divisions horizontales.

\textbf{Document 2 : le multimètre.} Un multimètre, branché aux bornes du GBF et réglé en position voltmètre AC (alternatif), affiche : $U_{lu} = \SI{218}{V}$.

\textbf{Document 3 : l'étiquette de la bouilloire.} Sur la plaque signalétique de la bouilloire de la cantine, on lit : « \SI{220}{V} -- \SI{50}{Hz} -- \SI{2000}{W} ».

\begin{popfigure}
\begin{tikzpicture}[>=Stealth, node distance=1.2cm and 1.5cm, thick]
% Styles
\tikzset{
    appareil/.style={draw=popdark, thick, rounded corners, minimum height=1.8cm, minimum width=2.8cm, fill=popcream, align=center},
    jonction/.style={circle, fill=popdark, inner sep=1.5pt},
    fil/.style={draw=popdark, very thick},
    etiquette/.style={font=\small, popdark, align=left}
}
% Nœuds
\node[appareil, align=center] (gbf){GBF\\Générateur BF\\$\sim$ \SI{50}{Hz}};
\node[jonction, right=of gbf] (j1) {};
\node[appareil, above=of j1, align=center] (osc){Oscilloscope\\Entrée CH1\\$S_v=\SI{100}{V/div}$\\$B=\SI{5}{ms/div}$};
\node[appareil, below=of j1, align=center] (mul){Multimètre\\Mode V$\sim$ (AC)\\Calibre \SI{600}{V}};
% Fils
\draw[fil] (gbf.east) -- (j1);
\draw[fil] (j1) -- (osc);
\draw[fil] (j1) -- (mul);
% Mise à la terre (symboles simples)
\draw[fil] (gbf.south) -- ++(0,-0.6) -- ++(-0.3,0) -- ++(0.6,0) -- ++(-0.3,0) -- ++(0,-0.15) -- ++(-0.2,0) -- ++(0.4,0);
\draw[fil] (osc.south) -- ++(0,-0.6) -- ++(-0.3,0) -- ++(0.6,0) -- ++(-0.3,0) -- ++(0,-0.15) -- ++(-0.2,0) -- ++(0.4,0);
\draw[fil] (mul.south) -- ++(0,-0.6) -- ++(-0.3,0) -- ++(0.6,0) -- ++(-0.3,0) -- ++(0,-0.15) -- ++(-0.2,0) -- ++(0.4,0);
% Légendes
\node[etiquette, above=0.1cm of gbf.north] {\textbf{Montage de simulation du secteur}};
\node[etiquette, right=0.2cm of osc.east, anchor=west, align=center]{Visualise la forme d'onde\\Mesure $U_{max}$, $T$};
\node[etiquette, right=0.2cm of mul.east, anchor=west, align=center]{Mesure $U$ (efficace)\\Affiche \SI{218}{V}};
\end{tikzpicture}
\legende{Montage expérimental : le GBF simule le réseau ENEO ; l'oscilloscope et le multimètre sont branchés en dérivation.}
\end{popfigure}

\textbf{Problème à résoudre :} la tension simulée (qui représente celle du quartier) est-elle conforme aux normes du réseau camerounais et aux indications de la bouilloire ?

\courssub{Consignes}

\begin{enumerate}
\item À partir du document 1, calculer la tension maximale $U_{max}$ de la tension observée.
\item Toujours à partir du document 1, déterminer la période $T$ de cette tension, puis calculer sa fréquence $f$ en hertz.
\item Calculer la tension efficace $U$ correspondant à cette tension maximale.
\item Comparer la valeur de $U$ obtenue à la question 3 avec l'indication $U_{lu}$ du multimètre (document 2). Ce résultat était-il prévisible ? Justifier.
\item Vérifier, avec les valeurs numériques obtenues, que la relation $U_{max} = U\sqrt{2}$ est bien satisfaite.
\item Comparer $U$ et $f$ aux normes du réseau camerounais et aux indications portées sur l'étiquette de la bouilloire (document 3). La tension du quartier est-elle conforme ?
\item Rédiger un court rapport (5 à 8 lignes) présentant les mesures, les calculs et la conclusion, en citant \textbf{deux consignes de sécurité} à respecter lors de mesures sur une installation électrique réelle.
\end{enumerate}

\courssub{Ressources à mobiliser}

\begin{aretenir}[Formules utiles]
\begin{itemize}
\item Tension maximale lue sur l'oscillogramme : $U_{max} = S_v \times y$, où $S_v$ est la sensibilité verticale (en V/div) et $y$ le nombre de divisions entre l'axe et un sommet.
\item Période : $T = B \times x$, où $B$ est le balayage (en ms/div) et $x$ le nombre de divisions d'un motif complet.
\item Fréquence : $f = \dfrac{1}{T}$, avec $T$ en secondes (s) et $f$ en hertz (Hz).
\item Tension efficace d'une tension sinusoïdale : $U = \dfrac{U_{max}}{\sqrt{2}}$, c'est-à-dire $U_{max} = U\sqrt{2}$.
\item Le voltmètre en position AC mesure directement la tension \cle{efficace}.
\end{itemize}
\end{aretenir}

\begin{attention}[Conversions indispensables]
Avant de calculer la fréquence, la période doit être convertie en \textbf{secondes} : $\SI{1}{ms} = \SI{0,001}{s} = 10^{-3}\ \text{s}$. Oublier cette conversion est l'erreur la plus fréquente au BEPC !
\end{attention}

\courssub{Démarche et corrigé commenté}

\begin{exoresolu}[Diagnostic de la tension du quartier]
Énoncé : un oscillogramme (sensibilité $\SI{100}{V/div}$, balayage $\SI{5}{ms/div}$) montre une sinusoïde d'amplitude $3,1$ divisions et de période $4$ divisions. Un voltmètre AC affiche $\SI{218}{V}$. La bouilloire porte la mention « \SI{220}{V} -- \SI{50}{Hz} ». Caractériser complètement la tension et conclure sur sa conformité.
\tcblower
\textbf{Étape 1 : tension maximale.}

On applique la formule $U_{max} = S_v \times y$ :
\begin{align*}
U_{max} &= S_v \times y \\
&= 100 \times 3,1 \\
&= \SI{310}{V}
\end{align*}
La tension maximale vaut $U_{max} = \SI{310}{V}$.

\textbf{Étape 2 : période et fréquence.}

Un motif complet occupe $4$ divisions avec un balayage de $\SI{5}{ms/div}$ :
\begin{align*}
T &= B \times x \\
&= 5 \times 4 \\
&= \SI{20}{ms}
\end{align*}
On convertit en secondes : $T = \SI{20}{ms} = \SI{0,020}{s} = 2 \times 10^{-2}\ \text{s}$.

La fréquence est l'inverse de la période :
\begin{align*}
f &= \dfrac{1}{T} \\
&= \dfrac{1}{0,020} \\
&= \SI{50}{Hz}
\end{align*}
La fréquence de la tension est $f = \SI{50}{Hz}$.

\textbf{Étape 3 : tension efficace.}

Pour une tension sinusoïdale :
\begin{align*}
U &= \dfrac{U_{max}}{\sqrt{2}} \\
&= \dfrac{310}{1,414} \\
&\approx \SI{219}{V}
\end{align*}
La tension efficace vaut environ $U \approx \SI{219}{V}$ (on arrondit au volt près).

\textbf{Étape 4 : comparaison avec le voltmètre.}

Le voltmètre en position AC affiche $U_{lu} = \SI{218}{V}$. Or un voltmètre en mode alternatif mesure la tension \emph{efficace}. Les deux valeurs ($\SI{219}{V}$ calculée et $\SI{218}{V}$ affichée) sont très proches : l'écart de $\SI{1}{V}$ correspond aux petites imprécisions de lecture sur l'écran (estimation des dixièmes de division) et sur l'appareil. Ce résultat était donc prévisible : il confirme que le voltmètre AC indique bien la tension efficace.

\textbf{Étape 5 : vérification de la relation $U_{max} = U\sqrt{2}$.}

Calculons $U\sqrt{2}$ avec la valeur mesurée par le voltmètre :
\begin{align*}
U\sqrt{2} &= 218 \times 1,414 \\
&\approx \SI{308}{V}
\end{align*}
On obtient environ $\SI{308}{V}$, valeur très proche de $U_{max} = \SI{310}{V}$ lu sur l'oscillogramme (écart inférieur à $1\ \%$). La relation $U_{max} = U\sqrt{2}$ est donc bien vérifiée expérimentalement.

\textbf{Étape 6 : conformité aux normes.}

Comparons les résultats aux normes du réseau camerounais et à l'étiquette de la bouilloire :
\begin{center}
\begin{tabularx}{\linewidth}{|l|X|X|X|}
\hline
\rowcolor{popblueL} \textbf{Grandeur} & \textbf{Valeur mesurée} & \textbf{Norme / étiquette} & \textbf{Conforme ?} \\
\hline
Tension efficace $U$ & $\SI{219}{V}$ & $\SI{220}{V}$ & Oui (écart $< 1\ \%$) \\
\hline
Fréquence $f$ & $\SI{50}{Hz}$ & $\SI{50}{Hz}$ & Oui \\
\hline
\end{tabularx}
\end{center}
La tension efficace et la fréquence correspondent aux valeurs nominales. La tension du quartier est donc \textbf{conforme} : la bouilloire peut fonctionner normalement.

\textbf{Étape 7 : rapport (modèle de rédaction).}

\emph{« Nous avons visualisé la tension à l'oscilloscope : l'amplitude de $3,1$ divisions à $\SI{100}{V/div}$ donne $U_{max} = \SI{310}{V}$, et le motif de $4$ divisions à $\SI{5}{ms/div}$ donne $T = \SI{20}{ms}$, soit $f = \SI{50}{Hz}$. La tension efficace calculée, $U = \dfrac{310}{\sqrt{2}} \approx \SI{219}{V}$, est confirmée par le voltmètre AC ($\SI{218}{V}$) et vérifie la relation $U_{max} = U\sqrt{2}$. Ces valeurs correspondent aux normes du réseau ENEO et à l'étiquette de la bouilloire : la tension du quartier est normale. Consignes de sécurité : (1) ne jamais toucher les bornes ou les fils dénudés d'un circuit sous tension ; (2) vérifier que le multimètre est bien réglé en position voltmètre AC sur un calibre adapté avant de le brancher, et le brancher toujours en dérivation, circuit hors tension si possible. »}
\end{exoresolu}

\courssub{Bilan de l'activité}

Cette activité t'a permis de mobiliser l'ensemble des savoir-faire de la leçon dans une situation authentique de la vie courante :
\begin{itemize}
\item tu as \textbf{exploité un oscillogramme} de bout en bout : lecture des réglages, mesure des divisions, calcul de $U_{max}$, de $T$ puis de $f$ ;
\item tu as compris le rôle du \textbf{voltmètre en position AC} : il affiche directement la tension efficace, ce qui évite le calcul $U = \dfrac{U_{max}}{\sqrt{2}}$ ;
\item tu as \textbf{vérifié expérimentalement} la relation fondamentale $U_{max} = U\sqrt{2}$, avec un écart inférieur à $1\ \%$ entre mesure et calcul ;
\item tu as \textbf{confronté des mesures à des normes} : les valeurs « \SI{220}{V} -- \SI{50}{Hz} » inscrites sur les appareils électriques vendus au Cameroun correspondent à la tension efficace et à la fréquence du réseau ENEO ;
\item tu as rédigé un \textbf{rapport scientifique} court, en justifiant chaque étape et en rappelant les consignes de sécurité : le courant du secteur est dangereux, toute mesure sur une installation réelle exige prudence, matériel adapté et respect des règles.
\end{itemize}

\begin{savaistu}[Pourquoi 220 V et 50 Hz ?]
Le réseau électrique camerounais, comme la plupart des réseaux africains et européens, distribue une tension alternative de valeur efficace \SI{220}{V} et de fréquence \SI{50}{Hz} : le courant change donc de sens $100$ fois par seconde ! Cette tension est produite par les alternateurs des centrales hydroélectriques (Song Loulou, Édéa, Lagdo) et des centrales thermiques. Les appareils sont construits pour fonctionner à ces valeurs : une tension trop faible fait mal chauffer la bouilloire, une tension trop forte peut griller les appareils. C'est pourquoi ENEO s'engage à maintenir la tension dans une marge étroite autour de \SI{220}{V}.
\end{savaistu}

\newpage

\coursec{Méthodes et exercices résolus}

\coursec{Les tensions alternatives sinusoïdales}

\begin{savaistu}[Le « courant de la prise » : une histoire de normes]
Au Cameroun, comme en Europe, le réseau ENEO distribue une tension alternative sinusoïdale de \textbf{fréquence 50 Hz} et de \textbf{tension efficace 220 V}. Ce choix historique (la « guerre des courants » entre Edison et Tesla/Westinghouse) s'est imposé car le courant alternatif se transporte sur de longues distances avec moins de pertes (grâce aux transformateurs) que le courant continu. La fréquence de 50 Hz est un compromis : assez élevée pour éviter le scintillement des lampes, assez basse pour limiter les pertes dans les lignes. Un groupe électrogène mal réglé (fréquence ou tension instable) peut endommager les appareils sensibles (téléviseurs, réfrigérateurs, ordinateurs).
\end{savaistu}

\begin{definition}[Tensions continue et alternative]
\begin{itemize}
\item Une \cle{tension continue} (symbole $=$ ou DC) garde une \textbf{valeur constante} et un \textbf{signe constant} au cours du temps. Elle est fournie par les piles, batteries, panneaux solaires.
\item Une \cle{tension alternative} (symbole $\sim$ ou AC) \textbf{change de signe} régulièrement (positive puis négative). Si sa courbe a la forme d'une sinusoïde, c'est une \cle{tension alternative sinusoïdale}. Elle est fournie par le réseau ENEO, les groupes électrogènes, les dynamos.
\end{itemize}
\end{definition}

\begin{experience}[Visualiser des tensions à l'oscilloscope]
\textbf{Matériel :} oscilloscope, pile 4,5 V, prise secteur (via transformateur de sécurité ou sonde atténuatrice 10:1 sous surveillance), dynamo de vélo, fils de connexion.
\textbf{Protocole :}
\begin{enumerate}
\item Brancher la pile aux entrées Y et GND de l'oscilloscope. Régler le balayage (TIME/DIV) sur 1 ms/div et la sensibilité (VOLTS/DIV) sur 1 V/div. Observer et dessiner.
\item Brancher la dynamo de vélo (en la faisant tourner à la main). Observer et dessiner.
\item \textbf{Sous surveillance du professeur uniquement :} visualiser la tension du secteur avec une sonde atténuatrice 10:1. Régler VOLTS/DIV sur 100 V/div (efficace 1000 V/div avec sonde 10:1) et TIME/DIV sur 5 ms/div. Observer et dessiner.
\end{enumerate}
\textbf{Observations :} Pile : trait horizontal au-dessus de l'axe. Dynamo et Secteur : courbes en « vagues » régulières traversant l'axe (sinusoïdes).
\textbf{Interprétation :} La pile délivre une tension continue. Le secteur et la dynamo délivrent des tensions alternatives sinusoïdales.
\textbf{Conclusion :} L'oscilloscope est l'instrument qui permet de \textbf{visualiser la forme} du signal (continue ou alternative) et de \textbf{mesurer} ses caractéristiques (période, tension maximale).
\end{experience}

\begin{popfigure}
\begin{tikzpicture}[scale=0.8, every node/.style={font=\small}]
% Grille oscilloscope
\draw[popmuted, step=1cm] (0,0) grid (10,6);
\draw[popline, thick] (0,3) -- (10,3) node[right] {$0\ \text{V}$}; % Axe zéro
\draw[popline, thick] (0,0) -- (0,6); % Axe temps
% Tension continue (Pile)
\draw[popblue, very thick] (0.5,4.5) -- (4.5,4.5);
\node[popblue, above] at (2.5,4.8) {\textbf{Tension continue (Pile)}};
\node[popblue, left] at (0.5,4.5) {$U_0$};
% Tension alternative (Secteur)
\draw[poporange, very thick, domain=5:9.5, samples=100, smooth] plot (\x, {3 + 2*sin(360*(\x-5))});
\draw[poporange, thick, <->] (5,3) -- (5,5) node[midway, right] {$U_{\text{max}}$};
\draw[poporange, thick, <->] (5,1) -- (5,5) node[midway, left] {$U_{\text{cc}}$};
\draw[poporange, thick, <->] (5,3) -- (6.57,3) node[midway, below] {$T$};
\node[poporange, above] at (7.25,5.5) {\textbf{Tension alternative sinusoïdale (Secteur)}};
% Graduations
\foreach \x in {0,1,...,10} \node[below, popmuted] at (\x,0) {\x};
\foreach \y in {0,1,...,6} \node[left, popmuted] at (0,\y) {\y};
\node[below] at (5,0) {Divisions horizontales (temps)};
\node[left, rotate=90] at (0,3) {Divisions verticales (tension)};
\end{tikzpicture}
\legende{Écran d'oscilloscope : à gauche, tension continue (pile) ; à droite, tension alternative sinusoïdale (secteur ENEO) avec repérage de $T$, $U_{\text{max}}$ et $U_{\text{cc}}$.}
\end{popfigure}

\courssub{Méthode 1 : Distinguer une tension continue d'une tension alternative}

\begin{methode}[Reconnaître le type d'une tension]
\begin{enumerate}
\item Observer le \cle{sens} et la \cle{valeur} de la tension au cours du temps (sur un oscillogramme ou un schéma).
\item Si la tension garde une valeur constante et un signe constant (toujours positive ou toujours négative) : c'est une tension \cle{continue}. Symbole : $=$ ou « DC ».
\item Si la tension change de signe régulièrement (positive puis négative) : c'est une tension \cle{alternative}. Symbole : $\sim$ ou « AC ».
\item Si, de plus, la courbe a la forme d'une sinusoïde (vagues régulières et symétriques par rapport à l'axe zéro) : c'est une tension \cle{alternative sinusoïdale}.
\item Conclure en citant la source : pile, batterie, panneau solaire (continue) ; prise de courant ENEO, groupe électrogène, dynamo (alternative sinusoïdale).
\end{enumerate}
\textbf{Réflexe :} le symbole $\sim$ sur un appareil signifie toujours « courant alternatif ».
\end{methode}

\begin{exoresolu}[Identifier des tensions]
On branche successivement un oscilloscope aux bornes : (a) d'une pile de \SI{4,5}{V} ; (b) d'une prise de courant ENEO (via sonde 10:1) ; (c) d'une dynamo de vélo actionnée à la main. On observe : (a) une droite horizontale au-dessus de l'axe ; (b) une sinusoïde ; (c) une sinusoïde. Pour chaque cas, préciser la nature de la tension et justifier.
\tcblower
\textbf{Cas (a) — la pile :} l'oscillogramme est une droite horizontale située au-dessus de l'axe des tensions : la tension garde une valeur constante ($\SI{4,5}{V}$) et un signe constant (positif). C'est donc une \textbf{tension continue}. C'est normal : une pile est un générateur de tension continue.

\textbf{Cas (b) — la prise ENEO :} l'oscillogramme est une sinusoïde : la tension varie au cours du temps en changeant de signe de façon régulière et symétrique. C'est une \textbf{tension alternative sinusoïdale}. Le réseau ENEO distribue une tension alternative de fréquence $\SI{50}{Hz}$.

\textbf{Cas (c) — la dynamo de vélo :} l'oscillogramme est également une sinusoïde : c'est une \textbf{tension alternative sinusoïdale}. La rotation de l'aimant devant la bobine produit une tension qui change de sens à chaque demi-tour.

\textbf{Conclusion :} (a) tension continue ; (b) et (c) tensions alternatives sinusoïdales. Seule la forme de la courbe permet de conclure, pas l'appareil utilisé pour la mesurer.
\end{exoresolu}

\courssub{Méthode 2 : Mesurer une période sur un oscillogramme}

\begin{definition}[Période et Fréquence]
Pour une tension alternative sinusoïdale :
\begin{itemize}
\item La \cle{période} $T$ (en secondes, s) est la durée du \textbf{motif élémentaire} le plus petit qui se répète (ex: d'une crête à la crête suivante).
\item La \cle{fréquence} $f$ (en hertz, Hz) est le nombre de périodes par seconde : $f = \dfrac{1}{T}$.
\end{itemize}
\end{definition}

\begin{methode}[Déterminer la période $T$]
\begin{enumerate}
\item Repérer sur l'écran un \cle{motif élémentaire} : la plus petite portion de courbe qui se répète identique à elle-même (par exemple d'une crête à la crête suivante, ou d'un passage par zéro montant au suivant).
\item Compter le nombre de divisions $x$ (horizontales) occupées par un motif complet.
\item Relever le \cle{balayage} (sensibilité horizontale) $b$, en \si{ms/div} ou \si{s/div} (bouton TIME/DIV).
\item Calculer : $T = x \times b$.
\item Convertir le résultat en secondes : $\SI{1}{ms} = \SI{e-3}{s}$.
\item \textbf{Précision :} mesurer sur plusieurs périodes ($n$ motifs sur $x_{\text{total}}$ divisions) puis diviser : $T = \dfrac{x_{\text{total}} \times b}{n}$.
\end{enumerate}
\end{methode}

\begin{exoresolu}[Lecture d'une période]
Sur l'écran d'un oscilloscope, on observe une tension sinusoïdale. Le balayage est réglé sur $b = \SI{5}{ms/div}$. Un motif complet occupe $4$ divisions horizontales.
\begin{enumerate}
\item Calculer la période $T$ de cette tension.
\item L'écran comporte $10$ divisions horizontales : combien de périodes observe-t-on ?
\end{enumerate}
\tcblower
\textbf{1. Calcul de la période.}

On applique la formule $T = x \times b$ avec $x = 4$ div et $b = \SI{5}{ms/div}$ :
\[ T = 4 \times 5 = \SI{20}{ms} \]
Conversion en secondes : $T = 20 \times 10^{-3} = \SI{0,02}{s}$.

\textbf{2. Nombre de périodes visibles.}

L'écran large de $10$ divisions correspond à une durée totale :
\[ \Delta t = 10 \times 5 = \SI{50}{ms} \]
Le nombre de périodes affichées est :
\[ n = \dfrac{\Delta t}{T} = \dfrac{50}{20} = 2,5 \]
On observe donc \textbf{2 périodes et demie} sur l'écran, ce qui est cohérent avec un motif de $4$ divisions ($10 / 4 = 2,5$).

\textbf{Conclusion :} la période de la tension est $T = \SI{20}{ms} = \SI{0,02}{s}$, et l'écran en affiche 2,5.
\end{exoresolu}

\courssub{Méthode 3 : Calculer une fréquence}

\begin{propriete}[Relation Période - Fréquence]
Pour tout signal périodique :
\[ f = \dfrac{1}{T} \qquad \text{avec } T \text{ en secondes (s) et } f \text{ en hertz (Hz)} \]
\[ T = \dfrac{1}{f} \]
\textbf{Attention :} $T$ doit \textbf{impérativement} être exprimé en secondes avant d'appliquer la formule.
\end{propriete}

\begin{methode}[Passer de la période à la fréquence]
\begin{enumerate}
\item Exprimer la période $T$ en \textbf{secondes} (conversion obligatoire si $T$ est en ms : diviser par 1000).
\item Appliquer la définition : $f = \dfrac{1}{T}$.
\item Exprimer le résultat en \cle{hertz} (Hz).
\item Vérifier l'ordre de grandeur : pour le réseau ENEO, on doit trouver $f = \SI{50}{Hz}$.
\end{enumerate}
\end{methode}

\begin{exoresolu}[Fréquence du réseau et d'un groupe électrogène]
\begin{enumerate}
\item La tension d'une prise ENEO a pour période $T = \SI{20}{ms}$. Calculer sa fréquence.
\item Un groupe électrogène défectueux délivre une tension de fréquence $f' = \SI{40}{Hz}$. Calculer sa période $T'$.
\item Comparer : lequel des deux générateurs produit la tension qui varie le plus rapidement ?
\end{enumerate}
\tcblower
\textbf{1. Fréquence de la tension ENEO.}

Conversion : $T = \SI{20}{ms} = 20 \times 10^{-3} = \SI{0,02}{s}$.

Application de la formule :
\[ f = \dfrac{1}{T} = \dfrac{1}{0,02} = \SI{50}{Hz} \]
La tension du réseau ENEO effectue $50$ périodes par seconde : c'est bien la valeur officielle au Cameroun.

\textbf{2. Période de la tension du groupe électrogène.}

On utilise la relation inverse :
\[ T' = \dfrac{1}{f'} = \dfrac{1}{40} = \SI{0,025}{s} = \SI{25}{ms} \]

\textbf{3. Comparaison.}

$f' = \SI{40}{Hz} < f = \SI{50}{Hz}$ : la tension du groupe effectue moins de périodes par seconde. C'est donc la tension \textbf{ENEO} qui varie le plus rapidement. Une fréquence anormale peut d'ailleurs endommager certains appareils sensibles (moteurs asynchrones, horloges électroniques, variateurs de vitesse).

\textbf{Conclusion :} $f = \SI{50}{Hz}$ pour ENEO ; $T' = \SI{25}{ms}$ pour le groupe ; la tension ENEO varie plus vite.
\end{exoresolu}

\courssub{Méthode 4 : Mesurer la tension maximale sur un oscillogramme}

\begin{definition}[Tension maximale et tension crête-à-crête]
\begin{itemize}
\item La \cle{tension maximale} $U_{\text{max}}$ (ou $U_{\text{crête}}$) est la valeur absolue maximale atteinte par la tension. Elle se mesure de l'\textbf{axe zéro} à la \textbf{crête} (sommet) de la sinusoïde. Unité : volt (V).
\item La \cle{tension crête-à-crête} $U_{\text{cc}}$ est la différence entre la crête positive et la crête négative : $U_{\text{cc}} = 2 \times U_{\text{max}}$.
\end{itemize}
\end{definition}

\begin{methode}[Déterminer la tension maximale $U_{\text{max}}$]
\begin{enumerate}
\item Repérer la \cle{sensibilité verticale} $s$ (en \si{V/div}), réglée sur le bouton « VOLTS/DIV ». \textbf{Attention :} si une sonde atténuatrice 10:1 est utilisée, la sensibilité effective est 10 fois celle affichée (ex: 10 V/div affiché $\rightarrow$ 100 V/div réel).
\item Mesurer le nombre de divisions $y$ entre l'axe horizontal (position zéro, trace sans signal) et le sommet (crête) de la sinusoïde.
\item Calculer : $U_{\text{max}} = y \times s$.
\item Vérifier que le zéro (trace sans signal) est bien sur l'axe central avant la mesure.
\end{enumerate}
\textbf{Réflexe :} $U_{\text{max}}$ se mesure uniquement à l'oscilloscope ; un voltmètre en position AC ne la donne pas (il donne $U_{\text{eff}}$).
\end{methode}

\begin{exoresolu}[Lecture de la tension maximale]
On visualise la tension d'une prise de courant à l'aide d'un oscilloscope muni d'une sonde atténuatrice 10:1. La sensibilité verticale affichée est de \SI{10}{V/div} (donc sensibilité réelle $s = \SI{100}{V/div}$). La crête de la sinusoïde se situe à $3,1$ divisions au-dessus de l'axe central.
\begin{enumerate}
\item Calculer la tension maximale $U_{\text{max}}$.
\item Quelle est la valeur crête à crête $U_{\text{cc}}$ ?
\end{enumerate}
\tcblower
\textbf{1. Tension maximale.}

La tension maximale est la hauteur entre l'axe des zéros et le sommet de la courbe :
\[ U_{\text{max}} = y \times s = 3,1 \times 100 = \SI{310}{V} \]

\textbf{2. Valeur crête à crête.}

La sinusoïde est symétrique par rapport à l'axe : la crête négative se situe à $3,1$ divisions en dessous de l'axe. La distance totale entre les deux crêtes est donc $2 \times 3,1 = 6,2$ div :
\[ U_{\text{cc}} = 2 \times U_{\text{max}} = 2 \times 310 = \SI{620}{V} \]

\textbf{Conclusion :} $U_{\text{max}} = \SI{310}{V}$ et $U_{\text{cc}} = \SI{620}{V}$. On retrouve la valeur connue : la tension maximale du réseau $\SI{220}{V}$ (efficace) vaut environ $\SI{311}{V}$.
\end{exoresolu}

\courssub{Méthode 5 : Exploiter la relation entre tension efficace et tension maximale}

\begin{propriete}[Valeur efficace d'une tension sinusoïdale]
Pour une \textbf{tension alternative sinusoïdale} uniquement, la tension efficace $U_{\text{eff}}$ (celle mesurée par un voltmètre en position AC) et la tension maximale $U_{\text{max}}$ (celle mesurée à l'oscilloscope) sont liées par :
\[ U_{\text{max}} = U_{\text{eff}} \times \sqrt{2} \qquad \text{soit} \qquad U_{\text{eff}} = \dfrac{U_{\text{max}}}{\sqrt{2}} \]
avec $\sqrt{2} \approx 1,41$.
\textbf{Propriété :} On a toujours $U_{\text{max}} > U_{\text{eff}}$.
\end{propriete}

\begin{methode}[Utiliser $U_{\text{max}} = U_{\text{eff}}\sqrt{2}$]
\begin{enumerate}
\item Identifier la grandeur connue : la tension \cle{efficace} $U_{\text{eff}}$ (lue au voltmètre en position AC, ou indiquée sur les appareils : « 220 V ») ou la tension \cle{maximale} $U_{\text{max}}$ (lue à l'oscilloscope).
\item Appliquer la relation, valable \textbf{uniquement pour une tension sinusoïdale} :
\[ U_{\text{max}} = U_{\text{eff}} \times \sqrt{2} \quad \text{ou} \quad U_{\text{eff}} = \dfrac{U_{\text{max}}}{\sqrt{2}} \]
\item Prendre $\sqrt{2} \approx 1,41$.
\item Vérifier la cohérence : on doit toujours avoir $U_{\text{max}} > U_{\text{eff}}$.
\item Préciser dans la rédaction quel appareil mesure quoi : voltmètre $\rightarrow U_{\text{eff}}$ ; oscilloscope $\rightarrow U_{\text{max}}$.
\end{enumerate}
\end{methode}

\begin{exoresolu}[Tension efficace et tension maximale du réseau]
La tension efficace délivrée par une prise de courant au Cameroun est $U_{\text{eff}} = \SI{220}{V}$.
\begin{enumerate}
\item Calculer la tension maximale $U_{\text{max}}$ de cette tension.
\item Un élève mesure cette tension avec un voltmètre en position « AC ». Quelle valeur affiche-t-il ? Justifier.
\item Un oscilloscope branché (avec sonde adaptée) sur la même prise montre une crête à $U_{\text{max}} = \SI{311}{V}$. Retrouver la valeur efficace.
\end{enumerate}
\tcblower
\textbf{1. Calcul de $U_{\text{max}}$.}

La tension du réseau est sinusoïdale, donc la relation $U_{\text{max}} = U_{\text{eff}}\sqrt{2}$ s'applique :
\[ U_{\text{max}} = 220 \times \sqrt{2} \approx 220 \times 1,414 \approx \SI{311}{V} \]

\textbf{2. Affichage du voltmètre.}

Un voltmètre utilisé en position « AC » (alternatif) mesure par construction la \textbf{tension efficace}. Il affiche donc :
\[ U_{\text{eff}} = \SI{220}{V} \]
C'est d'ailleurs pour cela qu'on dit couramment « la tension du secteur est de 220 V » : il s'agit toujours de la valeur efficace.

\textbf{3. Vérification à partir de l'oscilloscope.}

On applique la relation inverse :
\[ U_{\text{eff}} = \dfrac{U_{\text{max}}}{\sqrt{2}} = \dfrac{311}{1,414} \approx \SI{220}{V} \]
On retrouve exactement la valeur de départ : les deux mesures (voltmètre et oscilloscope) sont cohérentes.

\textbf{Conclusion :} $U_{\text{max}} \approx \SI{311}{V}$ ; le voltmètre affiche $\SI{220}{V}$ (valeur efficace) ; la relation $U_{\text{max}} = U_{\text{eff}}\sqrt{2}$ est vérifiée.
\end{exoresolu}

\courssub{Méthode 6 : Exploiter complètement un oscillogramme}

\begin{methode}[Bilan complet d'un oscillogramme]
Pour tirer toutes les informations d'un oscillogramme de tension sinusoïdale, suivre toujours le même plan :
\begin{enumerate}
\item \textbf{Nature} : la courbe change de signe et a une forme de sinusoïde $\rightarrow$ tension alternative sinusoïdale.
\item \textbf{Période} : $T = x \times b$ ($x$ divisions d'un motif, $b$ balayage), convertie en secondes.
\item \textbf{Fréquence} : $f = \dfrac{1}{T}$ en hertz.
\item \textbf{Tension maximale} : $U_{\text{max}} = y \times s$ ($y$ divisions de l'axe à la crête, $s$ sensibilité verticale réelle).
\item \textbf{Tension efficace} : $U_{\text{eff}} = \dfrac{U_{\text{max}}}{\sqrt{2}}$.
\item \textbf{Conclusion rédigée} : présenter les cinq résultats avec leurs unités.
\end{enumerate}
\textbf{Réflexe BEPC :} chaque étape non justifiée (formule absente, conversion oubliée, unité manquante) fait perdre des points : écris toujours la formule littérale avant l'application numérique.
\end{methode}

\begin{exoresolu}[Oscillogramme complet — type BEPC]
Un technicien visualise la tension aux bornes d'un groupe électrogène. Réglages : balayage $b = \SI{2}{ms/div}$ ; sensibilité verticale $s = \SI{50}{V/div}$. Sur l'écran, un motif complet occupe $10$ divisions horizontales et la crête atteint $2,8$ divisions verticales au-dessus de l'axe.
\begin{enumerate}
\item Quelle est la nature de cette tension ?
\item Déterminer sa période $T$ puis sa fréquence $f$.
\item Déterminer sa tension maximale $U_{\text{max}}$ puis sa tension efficace $U_{\text{eff}}$.
\item Ce groupe peut-il alimenter correctement des appareils prévus pour le réseau $220\ \text{V} - 50\ \text{Hz}$ ?
\end{enumerate}
\tcblower
\textbf{1. Nature de la tension.}

L'oscillogramme est une sinusoïde (motif régulier, symétrique, changeant de signe) : c'est une \textbf{tension alternative sinusoïdale}.

\textbf{2. Période et fréquence.}

Période :
\[ T = x \times b = 10 \times 2 = \SI{20}{ms} = \SI{0,02}{s} \]
Fréquence :
\[ f = \dfrac{1}{T} = \dfrac{1}{0,02} = \SI{50}{Hz} \]

\textbf{3. Tensions maximale et efficace.}

Tension maximale :
\[ U_{\text{max}} = y \times s = 2,8 \times 50 = \SI{140}{V} \]
Tension efficace :
\[ U_{\text{eff}} = \dfrac{U_{\text{max}}}{\sqrt{2}} = \dfrac{140}{1,414} \approx \SI{99}{V} \]

\textbf{4. Compatibilité avec les appareils.}

La fréquence est correcte ($\SI{50}{Hz}$), mais la tension efficace n'est que d'environ $\SI{99}{V}$, très inférieure aux $\SI{220}{V}$ attendus. Les appareils seront \textbf{sous-alimentés} : lampes faibles, moteurs qui peinent, démarrent mal ou surchauffent. Ce groupe est mal réglé (régime moteur trop faible, régulateur de tension défaillant) et ne convient pas en l'état.

\textbf{Conclusion :} tension alternative sinusoïdale ; $T = \SI{20}{ms}$ ; $f = \SI{50}{Hz}$ ; $U_{\text{max}} = \SI{140}{V}$ ; $U_{\text{eff}} \approx \SI{99}{V}$ ; le groupe ne peut pas alimenter correctement des appareils prévus pour $220\ \text{V} - 50\ \text{Hz}$.
\end{exoresolu}

\courssub{Méthode 7 : Appliquer les consignes de sécurité}

\begin{attention}[Danger : Tension secteur 220 V]
La tension efficace du secteur ($\SI{220}{V}$) correspond à une tension maximale de $\approx \SI{311}{V}$. Elle est \textbf{mortelle} (risque d'électrocution, fibrillation ventriculaire, brûlures internes). En manipulation scolaire, seules les tensions continues de \textbf{très basse tension de sécurité (TBTS)} $\leq \SI{24}{V}$ (piles, batteries, transformateurs de sécurité) sont autorisées sans surveillance rapprochée.
\end{attention}

\begin{methode}[Travailler en sécurité avec la tension du secteur (TP encadré)]
\begin{enumerate}
\item \textbf{Identifier le danger :} la tension du secteur ($\SI{220}{V}$ eff.) est mortelle.
\item \textbf{Oscilloscope :} ne jamais brancher directement la prise de courant sur un oscilloscope sans \cle{sonde atténuatrice 10:1} (ou 100:1) et sans l'encadrement explicite du professeur. Vérifier que la masse de l'oscilloscope (prise de terre) ne crée pas de court-circuit.
\item \textbf{Voltmètre / Multimètre :} vérifier la position du sélecteur : calibre \textbf{AC ($\sim$)} et calibre \textbf{immédiatement supérieur} à la valeur attendue (ex: $500$ V ou $750$ V AC) \textbf{avant} le branchement.
\item \textbf{Manipulation :} ne jamais toucher les parties métalliques des pointes de touche pendant la mesure ; avoir les \textbf{mains sèches} ; ne pas utiliser de rallonge ou de prise endommagée (fils dénudés, boîtier cassé) ; ne pas manipuler sous la pluie ou pieds nus dans l'eau.
\item \textbf{En cas d'accident :} \textbf{couper d'abord le courant} (disjoncteur général, débrancher la prise) avant de porter secours. Ne pas toucher la victime tant qu'elle est sous tension.
\end{enumerate}
\textbf{Réflexe BEPC :} citer au moins deux consignes concrètes (ex: « calibre AC 500 V », « mains sèches », « couper le courant avant secours ») vaut mieux qu'une phrase vague du type « faire attention ».
\end{methode}

\begin{exoresolu}[Situation de sécurité domestique]
Pendant une coupure ENEO, M. Etoundi branche un groupe électrogène. Sa fille Amina, élève de 3ème, veut mesurer la tension délivrée avec le multimètre de son père.
\begin{enumerate}
\item Sur quelle position (AC ou DC) et sur quel calibre doit-elle régler le multimètre ? Justifier.
\item Citer deux consignes de sécurité qu'elle doit respecter.
\item Le multimètre affiche $\SI{235}{V}$. Que représente cette valeur ? Calculer la tension maximale correspondante.
\end{enumerate}
\tcblower
\textbf{1. Réglage du multimètre.}

Le groupe électrogène délivre une tension \textbf{alternative} : Amina doit donc choisir la position \textbf{AC} (symbole $\sim$). La tension attendue est voisine de $\SI{220}{V}$ : elle doit choisir un calibre \textbf{immédiatement supérieur}, par exemple $500$ V ou $750$ V AC, pour ne pas endommager l'appareil et mesurer avec une bonne résolution.

\textbf{2. Consignes de sécurité.}

Par exemple :
\begin{itemize}
\item Avoir les \textbf{mains sèches} et ne pas toucher les parties métalliques des pointes de touche pendant la mesure.
\item Vérifier que les \textbf{fils et la prise ne sont pas dénudés} et que le groupe est posé au sec, à l'abri de la pluie.
\item Ne pas effectuer la mesure les pieds dans l'eau ou pieds nus.
\end{itemize}
(Deux consignes distinctes suffisent.)

\textbf{3. Signification de la valeur affichée.}

Un voltmètre en position AC mesure la \textbf{tension efficace} : $U_{\text{eff}} = \SI{235}{V}$.

La tension étant sinusoïdale :
\[ U_{\text{max}} = U_{\text{eff}}\sqrt{2} = 235 \times 1,414 \approx \SI{332}{V} \]

\textbf{Conclusion :} multimètre en AC, calibre supérieur à $220$ V ; mains sèches et fils en bon état ; la valeur affichée ($\SI{235}{V}$) est la tension efficace, et la tension maximale vaut environ $\SI{332}{V}$.
\end{exoresolu}

\begin{attention}[Les 5 erreurs qui coûtent des points au BEPC]
\begin{enumerate}
\item \textbf{Oublier la conversion des millisecondes en secondes} avant de calculer $f = \dfrac{1}{T}$. Avec $T = \SI{20}{ms}$, écrire $f = \dfrac{1}{20} = \SI{0,05}{Hz}$ est faux : il faut $f = \dfrac{1}{0,02} = \SI{50}{Hz}$.
\item \textbf{Mesurer $U_{\text{max}}$ de crête à crête.} La tension maximale se mesure de l'\textbf{axe des zéros à la crête} ; la valeur crête à crête vaut $2\,U_{\text{max}}$. Confondre les deux double le résultat.
\item \textbf{Appliquer $U_{\text{max}} = U_{\text{eff}}\sqrt{2}$ à n'importe quelle tension.} Cette relation n'est valable que pour une tension \textbf{alternative sinusoïdale} ; elle est fausse pour une tension continue, une tension en créneaux, ou une tension alternative non sinusoïdale.
\item \textbf{Affirmer que le voltmètre mesure $U_{\text{max}}$.} En position AC, le voltmètre mesure la tension \textbf{efficace} ; seul l'oscilloscope permet d'obtenir $U_{\text{max}}$. Le « 220 V » du réseau est une valeur efficace.
\item \textbf{Rédiger sans formule littérale ni unité.} Au BEPC, il faut écrire la formule ($T = x \times b$, $f = 1/T$, $U_{\text{max}} = y \times s$), puis l'application numérique, puis le résultat \textbf{avec son unité} (s, Hz, V) et une phrase de conclusion.
\end{enumerate}
\end{attention}

\begin{exercice}[Entraînement 1 : Lecture d'oscillogramme]
On visualise une tension alternative sinusoïdale sur un oscilloscope réglé sur : balayage $b = \SI{1}{ms/div}$, sensibilité $s = \SI{5}{V/div}$.
L'oscillogramme montre 3 périodes complètes sur la largeur de l'écran (10 divisions). La crête positive est à 2,5 divisions au-dessus de l'axe central.
\begin{enumerate}
\item Déterminer la période $T$ et la fréquence $f$ de cette tension.
\item Déterminer $U_{\text{max}}$ et $U_{\text{eff}}$.
\end{enumerate}
\end{exercice}

\begin{corrige}[Exercice 1]
\textbf{1.} Largeur écran = 10 div $\times$ 1 ms/div = 10 ms. 3 périodes dans 10 ms $\rightarrow$ $T = 10 / 3 \approx \SI{3,33}{ms} = \SI{3,33e-3}{s}$.
$f = 1 / T = 1 / 0,00333 \approx \SI{300}{Hz}$.
\textbf{2.} $U_{\text{max}} = 2,5 \times 5 = \SI{12,5}{V}$.
$U_{\text{eff}} = 12,5 / 1,414 \approx \SI{8,84}{V}$.
\end{corrige}

\begin{exercice}[Entraînement 2 : Choix du calibre]
Un élève veut mesurer la tension aux bornes d'une alimentation de PC portable (sortie \SI{19}{V} continu) avec un multimètre. Les calibres DC disponibles sont : 200 mV, 2 V, 20 V, 200 V, 600 V.
\begin{enumerate}
\item Quel calibre doit-il choisir ? Justifier.
\item S'il se trompe et choisit le calibre 2 V, que risque-t-il ?
\end{enumerate}
\end{exercice}

\begin{corrige}[Exercice 2]
\textbf{1.} La tension est continue ($\approx \SI{19}{V}$). Il faut choisir le calibre DC \textbf{immédiatement supérieur} à 19 V : le calibre \textbf{20 V}.
\textbf{2.} Sur le calibre 2 V, l'affichage saturera (souvent « 1 » ou « OL » pour OverLoad). L'appareil n'est généralement pas détruit pour une tension de 19 V sur un calibre 2 V (protection interne), mais la mesure est impossible. \textbf{Attention :} sur un calibre courant (A), une tension de 19 V ferait fondre le fusible interne.
\end{corrige}

\begin{popfigure}
\begin{tikzpicture}[scale=0.7, every node/.style={font=\small}]
% Schéma simplifié face avant oscilloscope
\draw[popdark, thick, rounded corners=5pt] (0,0) rectangle (12,8);
% Écran
\draw[popblueL, thick] (1,4) rectangle (7,7.5);
\draw[popmuted, step=0.5cm] (1,4) grid (7,7.5);
\draw[popline, thick] (1,5.75) -- (7,5.75); % Axe zéro
\draw[poporange, thick, domain=1:7, samples=100, smooth] plot (\x, {5.75 + 1.2*sin(360*(\x-1))});
% Boutons
\draw[popdark, fill=popcream] (8.5,6.5) circle (0.6) node {VOLTS/DIV};
\draw[popdark, fill=popcream] (10.5,6.5) circle (0.6) node {TIME/DIV};
\draw[popdark, fill=popcream] (8.5,5) circle (0.6) node {POS V};
\draw[popdark, fill=popcream] (10.5,5) circle (0.6) node {POS H};
\draw[popdark, fill=popcream] (8.5,3.5) circle (0.6) node {TRIG};
\draw[popdark, fill=popcream] (10.5,3.5) circle (0.6) node {AUTO};
% Prise entrée
\draw[popdark, fill=black] (1,1) circle (0.3) node[above=2pt, white] {Y};
\draw[popdark, fill=black] (3,1) circle (0.3) node[above=2pt, white] {GND};
\node[popdark, below] at (6,0) {Face avant simplifiée d'un oscilloscope analogique/numérique};
\end{tikzpicture}
\legende{Repérage des commandes principales d'un oscilloscope : VOLTS/DIV (sensibilité verticale $s$), TIME/DIV (balayage $b$), POS V/H (positionnement), TRIG (déclenchement).}
\end{popfigure}

\begin{aretenir}[Bilan de la leçon : Les tensions alternatives sinusoïdales]
\begin{itemize}
\item \textbf{Nature :} Tension qui change de signe régulièrement, forme sinusoïdale. Symbole $\sim$. Source : ENEO, groupe électrogène, dynamo.
\item \textbf{Visualisation :} Oscilloscope (écran, boutons VOLTS/DIV, TIME/DIV). Permet de voir la forme et de mesurer $T$ et $U_{\text{max}}$.
\item \textbf{Période $T$ (s) :} $T = x \times b$ ($x$ div horizontales d'un motif, $b$ balayage en s/div). \textbf{Toujours convertir en secondes.}
\item \textbf{Fréquence $f$ (Hz) :} $f = \dfrac{1}{T}$ (avec $T$ en s). Réseau Cameroun : $f = \SI{50}{Hz}$ ($T = \SI{20}{ms}$).
\item \textbf{Tension maximale $U_{\text{max}}$ (V) :} $U_{\text{max}} = y \times s$ ($y$ div verticales axe $\rightarrow$ crête, $s$ sensibilité V/div réelle). Mesurée \textbf{uniquement} à l'oscilloscope.
\item \textbf{Tension efficace $U_{\text{eff}}$ (V) :} $U_{\text{eff}} = \dfrac{U_{\text{max}}}{\sqrt{2}} \approx \dfrac{U_{\text{max}}}{1,41}$. Mesurée au voltmètre position AC. Relation valable \textbf{uniquement pour sinusoïdale}.
\item \textbf{Sécurité :} Secteur 220 V eff. (311 V max) = MORTEL. Calibre AC supérieur à la valeur attendue. Sonde atténuatrice obligatoire à l'oscilloscope. Mains sèches, fils intacts. Couper le courant avant secours.
\end{itemize}
\end{aretenir}

\newpage

\coursec{Exercices}

\courssub{Je vérifie mes connaissances}

\begin{exercice}[QCM : période et fréquence]
Une tension alternative sinusoïdale a une fréquence $f = \SI{100}{Hz}$. Sa période vaut :
\begin{enumerate}
\item[a)] $\SI{100}{s}$ ;
\item[b)] $\SI{0,01}{s}$ ;
\item[c)] $\SI{10}{s}$.
\end{enumerate}
Recopie la bonne réponse et justifie ton choix par un calcul.
\end{exercice}

\begin{exercice}[Vrai ou faux]
Réponds par \textbf{vrai} ou \textbf{faux}, puis justifie chaque réponse par une phrase.
\begin{enumerate}
\item La tension délivrée par une pile est une tension alternative.
\item Le voltmètre en position \textbf{AC} mesure la tension efficace d'une tension sinusoïdale.
\item La tension maximale d'une tension sinusoïdale est toujours plus petite que sa tension efficace.
\item L'oscilloscope permet de visualiser l'allure d'une tension en fonction du temps.
\end{enumerate}
\end{exercice}

\begin{exercice}[Définitions]
Définis chacun des termes suivants en une phrase, en précisant son symbole et son unité :
\begin{enumerate}
\item la période d'une tension alternative ;
\item la fréquence d'une tension alternative ;
\item la tension maximale ;
\item la tension efficace.
\end{enumerate}
\end{exercice}

\begin{exercice}[Calculs directs]
La tension du secteur au Cameroun a une fréquence $f = \SI{50}{Hz}$.
\begin{enumerate}
\item Calcule sa période $T$.
\item Sa valeur efficace est $U = \SI{220}{V}$. Calcule sa tension maximale $U_{\max}$. On donne : $U_{\max} = U\sqrt{2}$ et $\sqrt{2} \approx 1,41$.
\end{enumerate}
\end{exercice}

\courssub{Je m'entraîne}

\begin{exercice}[Lecture d'un oscillogramme]
Sur l'écran d'un oscilloscope, on observe une tension sinusoïdale. Les réglages sont : sensibilité verticale $S_v = \SI{2}{V/div}$ et balayage $B = \SI{5}{ms/div}$. L'amplitude de la courbe est de $3$ divisions et un motif complet occupe $4$ divisions horizontalement.
\begin{enumerate}
\item Calcule la tension maximale $U_{\max}$.
\item Calcule la période $T$, puis la fréquence $f$.
\item Calcule la tension efficace $U$.
\end{enumerate}
\end{exercice}

\begin{exercice}[Voltmètre et oscilloscope]
Un voltmètre branché en position \textbf{AC} aux bornes d'un GBF indique $U = \SI{6}{V}$.
\begin{enumerate}
\item Quelle tension maximale $U_{\max}$ observera-t-on à l'oscilloscope ?
\item Avec une sensibilité verticale de $\SI{2}{V/div}$, quelle déviation verticale (en divisions) cela représente-t-il ?
\end{enumerate}
\end{exercice}

\begin{exercice}[Cohérence de deux mesures]
Un élève mesure une tension sinusoïdale de deux façons : à l'oscilloscope, il lit $U_{\max} = \SI{14}{V}$ ; au voltmètre en position \textbf{AC}, il lit $U = \SI{9,9}{V}$.
\begin{enumerate}
\item Rappelle la relation entre $U_{\max}$ et $U$ pour une tension sinusoïdale.
\item Calcule $U\sqrt{2}$ avec $U = \SI{9,9}{V}$.
\item Les deux mesures sont-elles cohérentes ? Justifie.
\end{enumerate}
\end{exercice}

\begin{exercice}[Sécurité domestique]
\begin{enumerate}
\item Cite trois consignes à respecter avant de brancher un appareil électrique sur une prise du secteur.
\item Explique pourquoi il est dangereux de manipuler une prise de courant avec les mains mouillées.
\end{enumerate}
\end{exercice}

\courssub{Je me prépare au Bac}

\begin{exercice}[Le groupe électrogène du quartier (type BEPC)]
Pendant une coupure d'ENEO, M. Kamdem branche un groupe électrogène. Il veut vérifier la tension délivrée à l'aide d'un oscilloscope. Réglages : sensibilité verticale $S_v = \SI{100}{V/div}$, balayage $B = \SI{5}{ms/div}$. Sur l'écran, la courbe sinusoïdale a une amplitude de $3,1$ divisions et un motif occupe $4$ divisions.
\begin{enumerate}
\item La tension délivrée par le groupe est-elle continue ou alternative ? Justifie à partir de l'allure de la courbe.
\item Calcule la tension maximale $U_{\max}$ délivrée par le groupe.
\item Calcule la période $T$ puis la fréquence $f$ de cette tension. Le groupe respecte-t-il la fréquence du réseau camerounais ($\SI{50}{Hz}$) ?
\item Calcule la tension efficace $U$. Un appareil prévu pour $\SI{220}{V}$ peut-il être branché sans risque ?
\item Cite deux précautions à prendre lors de l'utilisation d'un groupe électrogène à la maison.
\end{enumerate}
\end{exercice}

\begin{exercice}[Réglage d'un GBF au laboratoire (type BEPC)]
Au laboratoire du collège, un professeur demande à ses élèves de régler un GBF pour obtenir une tension sinusoïdale de valeur efficace $U = \SI{12}{V}$ et de fréquence $f = \SI{200}{Hz}$. L'oscilloscope est réglé sur $S_v = \SI{5}{V/div}$ et $B = \SI{1}{ms/div}$.
\begin{enumerate}
\item Calcule la tension maximale $U_{\max}$ attendue.
\item Calcule la période $T$ de cette tension.
\item Quelle amplitude, en divisions, les élèves doivent-ils observer sur l'écran ?
\item Quelle largeur, en divisions, doit occuper un motif complet sur l'écran ?
\item Un élève mesure au voltmètre en position \textbf{DC} aux bornes du GBF. Que va-t-il observer ? Expliquer pourquoi cette mesure n'a pas de sens ici.
\end{enumerate}
\end{exercice}

\begin{exercice}[Installation électrique d'une maison à Yaoundé (type BEPC)]
La maison de la famille Ngo Bell est alimentée par le réseau ENEO : tension efficace $U = \SI{220}{V}$, fréquence $f = \SI{50}{Hz}$. Un technicien contrôle l'installation avec un oscilloscope réglé sur $S_v = \SI{100}{V/div}$ et $B = \SI{2,5}{ms/div}$.
\begin{enumerate}
\item Calcule la tension maximale $U_{\max}$ du secteur.
\item Calcule la période $T$ du secteur.
\item Détermine l'amplitude en divisions et la largeur d'un motif en divisions que le technicien doit observer si l'installation est normale.
\item Le technicien constate qu'une prise de la cuisine est endommagée et qu'un fil est dénudé près de l'évier. Explique pourquoi cette situation est particulièrement dangereuse et propose deux mesures de sécurité.
\item Pourquoi dit-on que la tension du secteur est \emph{alternative} alors que celle d'une batterie de moto est \emph{continue} ? Illustre ta réponse en décrivant l'allure de chaque courbe observée à l'oscilloscope.
\end{enumerate}
\end{exercice}

\newpage

\coursec{Corrigés des exercices}

\coursec{Exercices corrigés — Tensions alternatives sinusoïdales}

\courssub{Rappels et définitions}

\begin{corrige}[Exercice 1 — QCM : Période et fréquence ($f = \SI{100}{Hz}$)]
La bonne réponse est \textbf{b)} $T = \SI{0,01}{s}$.

\textbf{Justification :} la période et la fréquence sont liées par la relation $T = \dfrac{1}{f}$. Avec $f = \SI{100}{Hz}$ :
\[ T = \frac{1}{f} = \frac{1}{100} = \SI{0,01}{s}. \]
La réponse a) est fausse car elle confond $T$ et $f$ ; la réponse c) correspondrait à $f = \SI{0,1}{Hz}$.
\end{corrige}

\begin{corrige}[Exercice 2 — Vrai ou faux : notions de base]
\begin{enumerate}
\item \textbf{Faux.} La tension délivrée par une pile est une tension \cle{continue} : elle garde une valeur constante au cours du temps, elle ne change pas de signe.
\item \textbf{Vrai.} En position \textbf{AC} (courant alternatif), le voltmètre indique directement la \cle{tension efficace} $U$ de la tension sinusoïdale.
\item \textbf{Faux.} C'est l'inverse : $U_{\max} = U\sqrt{2}$, avec $\sqrt{2} \approx 1,41 > 1$, donc la tension maximale est toujours \emph{plus grande} que la tension efficace.
\item \textbf{Vrai.} L'oscilloscope affiche la courbe $u(t)$ : il permet de visualiser l'allure (forme, amplitude, période) d'une tension en fonction du temps.
\end{enumerate}
\end{corrige}

\begin{corrige}[Exercice 3 — Définitions à connaître pour le BEPC]
\begin{enumerate}
\item \textbf{Période} (symbole $T$, unité : la seconde, s) : c'est la durée d'un motif complet, c'est-à-dire le plus petit intervalle de temps au bout duquel la tension se reproduit identique à elle-même.
\item \textbf{Fréquence} (symbole $f$, unité : le hertz, Hz) : c'est le nombre de périodes (de motifs) par seconde ; elle est liée à la période par $f = \dfrac{1}{T}$.
\item \textbf{Tension maximale} (symbole $U_{\max}$, unité : le volt, V) : c'est la plus grande valeur atteinte par la tension au cours d'une période (l'amplitude de la courbe).
\item \textbf{Tension efficace} (symbole $U$, unité : le volt, V) : c'est la valeur indiquée par le voltmètre en position \textbf{AC} ; pour une tension sinusoïdale, $U = \dfrac{U_{\max}}{\sqrt{2}}$.
\end{enumerate}
\end{corrige}

\begin{corrige}[Exercice 4 — Calculs directs : le secteur camerounais]
\begin{enumerate}
\item La période est l'inverse de la fréquence :
\[ T = \frac{1}{f} = \frac{1}{50} = \SI{0,02}{s} = \SI{20}{ms}. \]
\item On applique la relation $U_{\max} = U\sqrt{2}$ :
\[ U_{\max} = 220 \times 1,41 = \SI{310,2}{V} \approx \SI{310}{V}. \]
\end{enumerate}
\textbf{Remarque :} quand on dit « le secteur délivre $\SI{220}{V}$ », il s'agit de la valeur \emph{efficace} ; la tension atteint en réalité environ $\SI{310}{V}$ au maximum.
\end{corrige}

\courssub{Exploitation d'oscillogrammes}

\begin{popfigure}
\begin{tikzpicture}[scale=0.75, >=Stealth]
% Grille 10 div x 6 div
\draw[popline, step=1cm] (0,0) grid (10,6);
% Axes centraux
\draw[popdark, thick] (0,3) -- (10,3) node[right] {$t$};
\draw[popdark, thick] (0,0) -- (0,6) node[above] {$u$};
% Trait zéro
\draw[popdark, dashed] (0,3) -- (10,3);
% Sinusoïde : amplitude 1.5 div (3 div crête-à-crête), période 4 div
\draw[popblue, thick, samples=200, domain=0:10] plot (\x, {3 + 1.5*sin(90*\x)});
% Repérage amplitude
\draw[poporange, <->] (1, 3) -- (1, 4.5) node[midway, left, poporange] {$3$ div};
\draw[poporange, <->] (1, 3) -- (1, 1.5) node[midway, left, poporange] {$3$ div};
% Repérage période
\draw[poporange, <->] (0, 1) -- (4, 1) node[midway, below, poporange] {$4$ div $= T$};
% Légendes paramètres
\node[popblue, anchor=north west] at (0.2, 5.8) {$S_v = 2\,\text{V/div}$};
\node[popblue, anchor=north west] at (0.2, 5.3) {$B = 5\,\text{ms/div}$};
\end{tikzpicture}
\legende{Oscillogramme type pour l'exercice 5 : lecture de l'amplitude (3 div) et de la période (4 div).}
\end{popfigure}

\begin{corrige}[Exercice 5 — Lecture d'un oscillogramme ($S_v = \SI{2}{V/div}$, $B = \SI{5}{ms/div}$)]
\begin{enumerate}
\item \textbf{Tension maximale :} on multiplie l'amplitude en divisions par la sensibilité verticale :
\[ U_{\max} = 3 \times S_v = 3 \times 2 = \SI{6}{V}. \]
\item \textbf{Période :} on multiplie la largeur d'un motif par le balayage :
\[ T = 4 \times B = 4 \times 5 = \SI{20}{ms} = \SI{0,02}{s}. \]
\textbf{Fréquence :}
\[ f = \frac{1}{T} = \frac{1}{0,02} = \SI{50}{Hz}. \]
\item \textbf{Tension efficace :}
\[ U = \frac{U_{\max}}{\sqrt{2}} = \frac{6}{1,41} \approx \SI{4,3}{V}. \]
\end{enumerate}
\textbf{Point de vigilance :} bien convertir les millisecondes en secondes ($\SI{20}{ms} = \SI{0,02}{s}$) avant de calculer la fréquence.
\end{corrige}

\begin{corrige}[Exercice 6 — Voltmètre et oscilloscope : aller-retour ($U = \SI{6}{V}$, $S_v = \SI{2}{V/div}$)]
\begin{enumerate}
\item Le voltmètre en position \textbf{AC} indique la tension efficace $U = \SI{6}{V}$. La tension maximale est :
\[ U_{\max} = U\sqrt{2} = 6 \times 1,41 \approx \SI{8,5}{V}. \]
\item La déviation verticale en divisions est :
\[ y = \frac{U_{\max}}{S_v} = \frac{8,5}{2} \approx 4,2 \text{ divisions}. \]
\end{enumerate}
\textbf{Point de vigilance :} ne pas confondre les deux relations : on \emph{multiplie} par $\sqrt{2}$ pour passer de $U$ à $U_{\max}$, on \emph{divise} par $\sqrt{2}$ pour passer de $U_{\max}$ à $U$.
\end{corrige}

\begin{corrige}[Exercice 7 — Cohérence de deux mesures ($U_{\max,\text{osc}} = \SI{14}{V}$, $U_{\text{volt}} = \SI{9,9}{V}$)]
\begin{enumerate}
\item Pour une tension sinusoïdale : $U_{\max} = U\sqrt{2}$, c'est-à-dire $U = \dfrac{U_{\max}}{\sqrt{2}}$.
\item Calculons la tension maximale attendue à partir de la valeur efficace :
\[ U\sqrt{2} = 9,9 \times 1,41 \approx \SI{13,96}{V} \approx \SI{14}{V}. \]
\item \textbf{Oui, les deux mesures sont cohérentes} : on retrouve $U\sqrt{2} \approx \SI{14}{V}$, qui est exactement la valeur de $U_{\max}$ lue à l'oscilloscope. La relation $U_{\max} = U\sqrt{2}$ est bien vérifiée (aux arrondis de lecture près).
\end{enumerate}
\end{corrige}

\courssub{Sécurité et contextes réels (type BEPC)}

\begin{corrige}[Exercice 8 — Sécurité domestique : prise de courant et eau]
\begin{enumerate}
\item Trois consignes avant de brancher un appareil sur une prise du secteur :
\begin{itemize}
\item vérifier que la fiche et le cordon ne sont pas endommagés (pas de fil dénudé) ;
\item avoir les mains bien sèches ;
\item vérifier que l'appareil est prévu pour la tension du secteur ($\SI{220}{V}$) et ne pas surcharger la prise avec trop d'appareils (éviter les multiprises surchargées).
\end{itemize}
\item L'eau du robinet (eau non pure) conduit le courant électrique. Avec les mains mouillées, la résistance du corps diminue fortement : le courant peut alors traverser le corps et provoquer une \cle{électrisation}, voire la mort par arrêt du cœur. C'est pourquoi il est interdit de toucher une prise ou un interrupteur avec les mains humides.
\end{enumerate}
\end{corrige}

\begin{popfigure}
\begin{tikzpicture}[scale=0.75, >=Stealth]
\draw[popline, step=1cm] (0,0) grid (10,6);
\draw[popdark, thick] (0,3) -- (10,3);
\draw[popdark, dashed] (0,3) -- (10,3);
\draw[popblue, thick, samples=200, domain=0:10] plot (\x, {3 + 1.55*sin(90*\x)}); % Amplitude ~3.1 div
\draw[poporange, <->] (1, 3) -- (1, 4.55) node[midway, left, poporange] {$3,1$ div};
\draw[poporange, <->] (0, 1) -- (4, 1) node[midway, below, poporange] {$4$ div};
\node[popblue, anchor=north west] at (0.2, 5.8) {$S_v = 100\,\text{V/div}$};
\node[popblue, anchor=north west] at (0.2, 5.3) {$B = 5\,\text{ms/div}$};
\end{tikzpicture}
\legende{Oscillogramme du groupe électrogène (Exercice 9) : amplitude 3,1 div, période 4 div.}
\end{popfigure}

\begin{corrige}[Exercice 9 — Le groupe électrogène du quartier (type BEPC)]
\begin{enumerate}
\item La tension est \textbf{alternative sinusoïdale} : la courbe observée à l'oscilloscope est une sinusoïde qui prend alternativement des valeurs positives et négatives, alors qu'une tension continue donnerait une droite horizontale.
\item \textbf{Tension maximale :}
\[ U_{\max} = 3,1 \times S_v = 3,1 \times 100 = \SI{310}{V}. \]
\item \textbf{Période :}
\[ T = 4 \times B = 4 \times 5 = \SI{20}{ms} = \SI{0,02}{s}. \]
\textbf{Fréquence :}
\[ f = \frac{1}{T} = \frac{1}{0,02} = \SI{50}{Hz}. \]
Le groupe respecte bien la fréquence du réseau camerounais.
\item \textbf{Tension efficace :}
\[ U = \frac{U_{\max}}{\sqrt{2}} = \frac{310}{1,41} \approx \SI{220}{V}. \]
Un appareil prévu pour $\SI{220}{V}$ peut donc être branché sans risque.
\item Deux précautions : placer le groupe \emph{à l'extérieur} de la maison (les gaz d'échappement contiennent du monoxyde de carbone, gaz mortel et invisible) ; ne pas le remplir d'essence lorsqu'il est chaud ou en marche (risque d'incendie) ; on peut aussi citer : ne pas le surcharger, le poser sur un sol sec et stable.
\end{enumerate}
\textbf{Barème indicatif (sur 10) :} question 1 : 1 pt ; question 2 : 2 pts ; question 3 : 3 pts ($T$ : 1 pt, $f$ : 1 pt, comparaison : 1 pt) ; question 4 : 2 pts ; question 5 : 2 pts.

\textbf{Points de vigilance :} citer la relation $U_{\max} = U\sqrt{2}$ avant de l'appliquer ; convertir $\SI{20}{ms}$ en $\SI{0,02}{s}$ ; conclure explicitement par une phrase de comparaison avec $\SI{50}{Hz}$ et $\SI{220}{V}$.
\end{corrige}

\begin{popfigure}
\begin{tikzpicture}[scale=0.75, >=Stealth]
\draw[popline, step=1cm] (0,0) grid (10,6);
\draw[popdark, thick] (0,3) -- (10,3);
\draw[popdark, dashed] (0,3) -- (10,3);
\draw[popblue, thick, samples=200, domain=0:10] plot (\x, {3 + 1.7*sin(180*\x)}); % f=200Hz -> T=5ms -> 5 div at 1ms/div. Amplitude 3.4 div
\draw[poporange, <->] (1, 3) -- (1, 4.7) node[midway, left, poporange] {$3,4$ div};
\draw[poporange, <->] (0, 1) -- (5, 1) node[midway, below, poporange] {$5$ div};
\node[popblue, anchor=north west] at (0.2, 5.8) {$S_v = 5\,\text{V/div}$};
\node[popblue, anchor=north west] at (0.2, 5.3) {$B = 1\,\text{ms/div}$};
\end{tikzpicture}
\legende{Réglage GBF visé (Exercice 10) : amplitude 3,4 div, période 5 div.}
\end{popfigure}

\begin{corrige}[Exercice 10 — Réglage d'un GBF au laboratoire (type BEPC : $U = \SI{12}{V}$, $f = \SI{200}{Hz}$, $S_v = \SI{5}{V/div}$, $B = \SI{1}{ms/div}$)]
\begin{enumerate}
\item \textbf{Tension maximale attendue :}
\[ U_{\max} = U\sqrt{2} = 12 \times 1,41 \approx \SI{16,9}{V} \approx \SI{17}{V}. \]
\item \textbf{Période :}
\[ T = \frac{1}{f} = \frac{1}{200} = \SI{0,005}{s} = \SI{5}{ms}. \]
\item \textbf{Amplitude en divisions :}
\[ y = \frac{U_{\max}}{S_v} = \frac{16,9}{5} \approx 3,4 \text{ divisions}. \]
\item \textbf{Largeur d'un motif en divisions :}
\[ x = \frac{T}{B} = \frac{5}{1} = 5 \text{ divisions}. \]
\item En position \textbf{DC}, le voltmètre mesure la valeur \emph{moyenne} de la tension. Or une tension sinusoïdale prend autant de valeurs positives que négatives : sa valeur moyenne est nulle. Le voltmètre indiquera donc environ $\SI{0}{V}$, ce qui ne correspond à rien d'utile ici : pour mesurer une tension alternative, il faut utiliser la position \textbf{AC} (qui donne la tension efficace).
\end{enumerate}
\textbf{Barème indicatif (sur 10) :} question 1 : 2 pts ; question 2 : 2 pts ; question 3 : 2 pts ; question 4 : 2 pts ; question 5 : 2 pts.

\textbf{Points de vigilance :} dans les questions 3 et 4, on fait le raisonnement « inverse » de la lecture d'oscillogramme : on \emph{divise} par la sensibilité ou le balayage ; pour la question 5, expliciter que la moyenne d'une sinusoïde est nulle.
\end{corrige}

\begin{popfigure}
\begin{tikzpicture}[scale=0.75, >=Stealth]
\draw[popline, step=1cm] (0,0) grid (10,6);
\draw[popdark, thick] (0,3) -- (10,3);
\draw[popdark, dashed] (0,3) -- (10,3);
\draw[popblue, thick, samples=200, domain=0:10] plot (\x, {3 + 1.55*sin(45*\x)}); % T=20ms -> 8 div at 2.5ms/div. Amp 3.1 div
\draw[poporange, <->] (1, 3) -- (1, 4.55) node[midway, left, poporange] {$3,1$ div};
\draw[poporange, <->] (0, 1) -- (8, 1) node[midway, below, poporange] {$8$ div};
\node[popblue, anchor=north west] at (0.2, 5.8) {$S_v = 100\,\text{V/div}$};
\node[popblue, anchor=north west] at (0.2, 5.3) {$B = 2,5\,\text{ms/div}$};
\end{tikzpicture}
\legende{Oscillogramme secteur attendu (Exercice 11) : amplitude 3,1 div, période 8 div.}
\end{popfigure}

\begin{corrige}[Exercice 11 — Installation électrique d'une maison à Yaoundé (type BEPC : $U = \SI{220}{V}$, $f = \SI{50}{Hz}$, $S_v = \SI{100}{V/div}$, $B = \SI{2,5}{ms/div}$)]
\begin{enumerate}
\item \textbf{Tension maximale du secteur :}
\[ U_{\max} = U\sqrt{2} = 220 \times 1,41 = \SI{310,2}{V} \approx \SI{310}{V}. \]
\item \textbf{Période du secteur :}
\[ T = \frac{1}{f} = \frac{1}{50} = \SI{0,02}{s} = \SI{20}{ms}. \]
\item \textbf{Amplitude attendue en divisions :}
\[ y = \frac{U_{\max}}{S_v} = \frac{310}{100} \approx 3,1 \text{ divisions}. \]
\textbf{Largeur d'un motif :}
\[ x = \frac{T}{B} = \frac{20}{2,5} = 8 \text{ divisions}. \]
\item Un fil dénudé près de l'évier est particulièrement dangereux car l'eau conduit le courant : en touchant le fil ou l'eau souillée, le courant peut traverser le corps et provoquer une électrisation mortelle. Deux mesures de sécurité : couper immédiatement le courant au disjoncteur général avant toute intervention ; faire réparer la prise et isoler le fil par un technicien qualifié (ne jamais toucher le fil à mains nues, encore moins mouillées).
\item La tension du secteur est \emph{alternative} car elle change de signe et de valeur en permanence : à l'oscilloscope, on observe une sinusoïde qui monte et descend régulièrement autour de zéro. La tension d'une batterie de moto est \emph{continue} car elle garde une valeur constante et un signe fixe : à l'oscilloscope, on observe une droite horizontale (au-dessus de zéro).
\end{enumerate}
\textbf{Barème indicatif (sur 10) :} question 1 : 1,5 pt ; question 2 : 1,5 pt ; question 3 : 2 pts ; question 4 : 2,5 pts ; question 5 : 2,5 pts.

\textbf{Points de vigilance :} bien écrire les relations littérales ($U_{\max} = U\sqrt{2}$, $T = \dfrac{1}{f}$) avant l'application numérique ; dans la question 5, la justification doit s'appuyer sur l'\emph{allure des courbes} (sinusoïde contre droite horizontale) et sur le changement de signe, pas seulement sur des définitions apprises par cœur.
\end{corrige}

\newpage

\coursec{Fiche bilan}

\begin{popfigure}
\begin{tikzpicture}[mindmap, grow cyclic, every node/.style={concept, font=\small},
root concept/.append style={concept color=popblue, font=\bfseries, minimum size=2.6cm},
level 1/.append style={level distance=3.4cm, sibling angle=51.4, minimum size=2.2cm}]
\node[root concept]{Tensions alternatives sinusoïdales}
child[concept color=poporange]{node[align=center]{Deux types de tension\\ continue : constante\\ alternative : change de signe}}
child[concept color=popgreen]{node[align=center]{L'oscilloscope\\ visualise $u(t)$\\ balayage + sensibilité}}
child[concept color=poppurple]{node[align=center]{Période $T$ (s)\\ $T = N \times B$\\ plus petit motif répété}}
child[concept color=poppink]{node[align=center]{Fréquence $f$ (Hz)\\ $f = \dfrac{1}{T}$\\ réseau : \SI{50}{Hz}}}
child[concept color=popgold]{node[align=center]{Tension maximale $U_{\max}$ (V)\\ $U_{\max} = n \times s$}}
child[concept color=popblue!60]{node[align=center]{Tension efficace $U$ (V)\\ $U = \dfrac{U_{\max}}{\sqrt{2}}$\\ voltmètre en AC}}
child[concept color=popdark]{node[align=center]{Sécurité\\ $U = \SI{220}{V}$ : dangereux\\ ne jamais toucher les fils}};
\end{tikzpicture}
\legende{Carte mentale de la leçon : les sept idées clés sur les tensions alternatives sinusoïdales.}
\end{popfigure}

\begin{bilan}
\textbf{1. Définitions clés}
\begin{itemize}
\item Une \cle{tension continue} garde une valeur constante au cours du temps (pile, batterie).
\item Une \cle{tension alternative} change de valeur et de signe au cours du temps ; elle est \cle{sinusoïdale} si sa courbe est une sinusoïde (secteur ENEO, groupe électrogène).
\item L'\cle{oscilloscope} est l'appareil qui permet de visualiser une tension en fonction du temps : l'écran affiche l'\cle{oscillogramme}.
\item La \cle{période} $T$ est la durée du plus petit motif qui se répète ; elle s'exprime en secondes (s).
\item La \cle{fréquence} $f$ est le nombre de périodes par seconde ; elle s'exprime en hertz (Hz).
\item La \cle{tension maximale} $U_{\max}$ est la plus grande valeur atteinte par la tension.
\item La \cle{tension efficace} $U$ est la valeur indiquée par un voltmètre en mode alternatif (AC).
\end{itemize}

\textbf{2. Formules à connaître par cœur}
\begin{center}
\begin{tabularx}{\linewidth}{|l|X|l|}\hline
\rowcolor{popblueL} \textbf{Grandeur} & \textbf{Formule} & \textbf{Unité} \\ \hline
Période (oscillogramme) & $T = N \times B$ ($N$ : nombre de divisions d'un motif ; $B$ : balayage) & seconde (s) \\ \hline
Fréquence & $f = \dfrac{1}{T}$ & hertz (Hz) \\ \hline
Tension maximale & $U_{\max} = n \times s$ ($n$ : divisions du pic ; $s$ : sensibilité verticale) & volt (V) \\ \hline
Tension efficace & $U = \dfrac{U_{\max}}{\sqrt{2}}$ & volt (V) \\ \hline
\end{tabularx}
\end{center}

\textbf{3. Méthodes express}
\begin{itemize}
\item \textbf{Lire un oscillogramme :} repérer le balayage $B$ (en ms/div) et la sensibilité $s$ (en V/div), compter les divisions, puis appliquer $T = N \times B$ et $U_{\max} = n \times s$.
\item \textbf{Convertir :} $\SI{1}{ms} = \SI{e-3}{s}$ ; penser à convertir le balayage en secondes avant de calculer $f$.
\item \textbf{Passer de $U_{\max}$ à $U$ :} diviser par $\sqrt{2} \approx 1{,}414$. Exemple : $U_{\max} \approx \SI{311}{V}$ donne $U \approx \SI{220}{V}$.
\item \textbf{Sécurité :} la tension du secteur (\SI{220}{V}, \SI{50}{Hz}) est dangereuse ; ne jamais toucher un fil dénudé, ne pas surcharger les rallonges, couper le courant avant toute intervention.
\end{itemize}
\end{bilan}

\begin{aretenir}[Checklist avant le BEPC]
\begin{itemize}
\item[$\square$] Je sais distinguer une tension continue d'une tension alternative sur un oscillogramme.
\item[$\square$] Je sais nommer l'appareil qui visualise une tension : l'oscilloscope.
\item[$\square$] Je sais définir la période $T$ et la mesurer avec $T = N \times B$.
\item[$\square$] Je sais convertir les millisecondes en secondes sans me tromper.
\item[$\square$] Je sais calculer la fréquence avec $f = \dfrac{1}{T}$ et donner l'unité (Hz).
\item[$\square$] Je connais la fréquence du réseau : $f = \SI{50}{Hz}$ (soit $T = \SI{20}{ms}$).
\item[$\square$] Je sais mesurer la tension maximale avec $U_{\max} = n \times s$.
\item[$\square$] Je sais calculer la tension efficace avec $U = \dfrac{U_{\max}}{\sqrt{2}}$.
\item[$\square$] Je sais que le voltmètre en mode AC affiche la tension efficace.
\item[$\square$] Je connais la valeur du secteur : $U \approx \SI{220}{V}$.
\item[$\square$] Je cite au moins deux règles de sécurité électrique à la maison.
\item[$\square$] Je rédige mes calculs avec les unités à chaque étape.
\end{itemize}
\end{aretenir}

\findecours

\end{document}
